\documentclass[11pt,a4paper]{article}

% Packages
\usepackage[utf8]{inputenc}
\usepackage[T1]{fontenc}
\usepackage{amsmath,amssymb,amsthm}
\usepackage{mathtools}
\usepackage{graphicx}
\usepackage{hyperref}
\usepackage{cleveref}
\usepackage{booktabs}
\usepackage{array}
\usepackage{xcolor}
\usepackage{microtype}
\usepackage[numbers]{natbib}

% Page layout
\usepackage[margin=1in]{geometry}
\emergencystretch=1em
\tolerance=1000

% Theorem environments
\theoremstyle{plain}
\newtheorem{theorem}{Theorem}[section]
\newtheorem{lemma}[theorem]{Lemma}
\newtheorem{corollary}[theorem]{Corollary}
\newtheorem{proposition}[theorem]{Proposition}
\newtheorem{conjecture}[theorem]{Conjecture}

\theoremstyle{definition}
\newtheorem{definition}[theorem]{Definition}
\newtheorem{remark}[theorem]{Remark}

% Notation shortcuts
\newcommand{\F}{\mathbb{F}}
\newcommand{\LL}{\mathcal{L}}
\newcommand{\E}{\mathbb{E}}
\newcommand{\Prb}{\Pr}
\DeclareMathOperator{\Lagr}{Lagr}
\DeclareMathOperator{\SDA}{SDA}
\DeclareMathOperator{\VSTAT}{VSTAT}
\DeclareMathOperator{\Adv}{Adv}
\DeclareMathOperator{\Var}{Var}
\newcommand{\etal}{\textit{et al.}}
\newcommand{\R}{\mathbb{R}}
\newcommand{\LSN}{\mathsf{LSN}}

\title{Lagrangian Subspace Noise: \\ Exact Statistical-Query Lower Bounds and the Linear-Reduction Landscape}
\author{Kwanghoo Choo\thanks{ORCID: \href{https://orcid.org/0009-0005-5682-8098}{0009-0005-5682-8098}. Email: \texttt{gattochoo@gmail.com}. \textcopyright~2026 Kwanghoo Choo. Licensed under CC BY 4.0.}\\ \normalsize Independent Researcher}
\date{June 2026}

\begin{document}

\maketitle

\begin{abstract}
Lagrangian Subspace Noise (LSN) is a variant of Learning Parity with Noise whose secret is a Lagrangian subspace of a symplectic space over $\F_2$ rather than a linear function. It is a classical, membership-oracle problem in the family of the stabilizer-decoding / symplectic-LPN formulations recently proposed as quantum-native hardness sources by Khesin \etal\ and Lu \etal; its precise relation to those computational formulations is separated below. We give the first \emph{exact} (asymptotic-error-free) statistical analysis of its membership formulation, and a near-complete map of its linear-reduction landscape.

A character-sum method makes the symplectic constraint exactly computable. The pairwise correlation of two LSN distributions is \emph{exactly} $\tfrac{(1-2p)^2}{p(1-p)}\,2^{\,j-2n}$ (coefficient $4/3$ at $p=1/4$), with no asymptotic error term, where $j$ is the dimension of the secrets' intersection; and every all-ones subset moment of the isotropic row ensemble has an exact closed form, so a fixed-size bundle of rows deviates from the unconstrained LPN ensemble by $\Theta(4^{-n})$ --- the symplectic relation $S_A=0$ is invisible to any constant-size statistical test. From the exact correlations we obtain an \emph{unconditional} Statistical Query (SQ) lower bound of $\Omega(2^n)$ queries on a worst-case promise problem via an explicit symplectic spread, and an \emph{unconditional} $2^{2n-O(1)}$ average-case bound (via a size-constrained Delsarte certificate on the big cell and an augmentation/averaging transfer to all Lagrangians; only a sharp universal constant remains conjectural); the same method yields exact correlations and a $2^{c_p n}$-query bound for distinguishing the \emph{averaged} symplectic-LPN ensemble --- in both the Bernoulli and the externally-relevant depolarizing-symplectic noise models, the latter being the object whose hardness Khesin \etal\ and Lu \etal\ reduce from constant-rate LPN --- whose off-diagonal likelihood-ratio correlation is an exact, exponentially small negative quantity in place of the $0$ of the unconstrained ensemble. Moreover, an inspection of that external reduction shows its depolarizing step is removable, so the independent-Bernoulli sympLPN we analyze inherits a constant-rate-LPN hardness \emph{floor} --- a lower bound, not a claim of novelty; independence from the established families remains open.

For reductions, we introduce a \emph{transport} invariant $C^{\top}MC=0$ inherited from the secret's isotropy and combine it with an entropy-support counting bound to close three of the four cells of the linear-reduction landscape (public-matrix, fixed, conditionally-uniform adaptive) unconditionally, and the deterministic half of the fourth, leaving a single precisely-stated open cell (randomized marginal-adaptive). This is independent of, and complementary to, the concurrent decoding-formulation barriers of Lu \etal: our transport invariant and exact-moment analysis do not use their entropy argument. Our statistical-query and information-theoretic analysis works in the membership formulation --- its natural form --- while the linear-reduction landscape and the hardness floor above concern the computational sympLPN formulation; we separate the membership, batch, and sympLPN formulations by information-theoretic floors. While the LSN problem belongs to the recently proposed stabilizer-decoding / symplectic-LPN family rather than constituting a new hardness family of its own, the contributions developed here---the first exact, asymptotic-error-free statistical analysis of this problem class and a near-complete linear-reduction landscape---are new; every claim is a proved theorem, an explicitly labeled piece of empirical evidence, or an explicitly labeled conjecture.
\end{abstract}

\paragraph{\textbf{Keywords.}} Post-quantum cryptography, Learning Parity with Noise, statistical query model, symplectic geometry, lower bounds, black-box reductions.

\section{Introduction}
\label{sec:intro}

\subsection{Background}
\label{subsec:background}

The Learning Parity with Noise (LPN) problem, introduced by Blum \etal\ \cite{BFKL94}, asks to recover a secret linear function $s \in \F_2^k$ given noisy inner products $(a_i, \langle a_i, s \rangle \oplus e_i)$. Despite decades of cryptanalytic effort, no polynomial-time algorithm is known for LPN with constant noise rate. The best classical attacks run in sub-exponential time $2^{O(k / \log k)}$ using the Blum--Kalai--Wasserman (BKW) algorithm \cite{BKW03} and its modern refinements \cite{LF06,Lyu05,EKM17,GJL14,Kir11}. Grover's algorithm reduces the naive key-search cost from $2^k$ to $2^{k/2}$ \cite{Gro96}, and quantum variants of BKW remain the dominant quantum threat; no exponential quantum speed-up over the best classical methods is known.

LPN has been the foundation for a wide range of cryptographic primitives, including OWFs, PKE, OT, and MPC \cite{ABW10,Ale11,DGM12,BCGI18}. Its main theoretical limitation is the absence of a worst-case hardness reduction comparable to Regev's LWE reduction from lattice problems \cite{Reg05}. LPN hardness is therefore treated as a \emph{primitive assumption}, accepted on the basis of sustained cryptanalytic resistance rather than reducibility.

Post-quantum cryptography currently rests on a handful of established hardness families: lattices (LWE/SIS), error-correcting codes (LPN/syndrome decoding), multivariate quadratic equations (MQ), hash functions, and isogenies, with tensor and group-action problems emerging as further candidates. Shor's algorithm \cite{Sho94} breaks the integer-factoring and discrete-logarithm assumptions underlying classical public-key cryptography, motivating the search for quantum-resistant alternatives. A genuinely new family --- structurally distinct and not reducible to the existing ones --- has long been sought. LSN is a recent candidate in this spirit; this paper studies its hardness landscape without asserting its standing among the established families; what is new here is the \emph{method}---the first exact, asymptotic-error-free statistical analysis of this problem class, together with a near-complete linear-reduction landscape.

\subsection{The LSN Problem}
\label{subsec:lsn}

We propose \emph{Lagrangian Subspace Noise} (LSN); its name and motivation come from the \emph{Learning Stabilizers with Noise} problem of \cite{PQS26} and the \emph{symplectic LPN} problem of \cite{KLP+25,LPQR26}, and the precise relation between our membership formulation and theirs is discussed in \Cref{subsec:two-forms}. The secret is not a linear function but a Lagrangian subspace $L \subset \F_2^{2n}$ of a symplectic vector space. The learner receives samples
\[
(a_i, b_i) \quad \text{where} \quad b_i = \mathbf{1}_L(a_i) \oplus e_i,
\]
with $a_i \sim \F_2^{2n}$ uniform, $e_i \sim \operatorname{Bernoulli}(p)$, and $\mathbf{1}_L$ the indicator of $L$. The secret space is the Lagrangian Grassmannian:
\begin{equation}
\label{eq:lagr-count}
|\Lagr(2n, \F_2)| = \prod_{i=1}^n (2^i + 1) = 2^{n(n+1)/2 + O(1)}.
\end{equation}
This is exponentially larger than LPN's $2^n$ secrets, yet it is not an unstructured space: every secret satisfies the strong isotropy constraint $L = L^{\perp_\Omega}$. It is exactly this tension --- exponential size versus rigid structure --- that makes LSN both attractive as a hardness candidate and technically difficult to analyze.

Recent independent work has put the stabilizer-decoding form of the problem on a firm footing: Poremba--Quek--Shor \cite{PQS26} introduced the Learning-Stabilizers-with-Noise problem (which contains LPN as a special case) and proved a worst-case-to-average-case reduction for a restricted (non-standard) variant of the problem; Khesin \etal\ \cite{KLP+25} proved that constant-rate LPN reduces to \emph{their} stabilizer-decoding LSN even for a single logical qubit; and Lu \etal\ \cite{LPQR26} established strong barriers against reducing symplectic LPN to standard LPN (for linear reductions, in the constant-rate sample regime $m=\Theta(n)$; see \Cref{sec:barriers}). The precise relation between their LSN and our membership formulation, and the broader question of LSN's independence from the established families, are stated as explicit open positioning items in \Cref{subsec:two-forms}. This concurrent progress motivates the present study of LSN's own statistical and reduction structure.

\subsection{Contributions}
\label{subsec:contributions}

Our central technical tool is a character-sum method that renders the symplectic isotropy constraint $L = L^{\perp_\Omega}$ exactly computable, producing asymptotic-error-free correlations and moments where prior analyses give only bounds. From the exact correlations we derive statistical-query lower bounds; from the exact moments, the constant-size indistinguishability of the isotropic ensemble from unconstrained LPN; and from a transport invariant inherited from isotropy, a near-complete landscape of linear reductions to LPN. This situates our work alongside, but methodologically apart from, the concurrent results of Khesin \etal\ \cite{KLP+25}, Poremba--Quek--Shor \cite{PQS26}, and Lu \etal\ \cite{LPQR26}, which establish the average-case cryptographic hardness of the public-matrix stabilizer-decoding / symplectic-LPN problem (and, in \cite{LPQR26}, bar linear reductions to LPN through an entropy argument). We instead analyze the membership oracle directly --- computing its exact moment and correlation structure --- and close the reduction landscape through a transport invariant inherited from isotropy, independent of and complementary to that entropy method.

\begin{enumerate}
    \item \textbf{Exact SQ lower bound (\Cref{sec:sq}).} We prove an unconditional $\Omega(2^n)$ SQ query lower bound on a worst-case promise problem (an explicit symplectic spread), and an \emph{unconditional} average-case $q \geq 2^{2n - O(1)}$ bound (\Cref{thm:main-sq-cond}), via the big-cell theorem and an augmentation/averaging transfer to all Lagrangians; only the sharp pencil-extremality constant (\Cref{conj:pencil}) remains conjectural, and would merely improve the VSTAT parameter. For the sympLPN formulation itself we prove exact likelihood-ratio correlations (\Cref{thm:symplpn-corr}; the isotropic conditioning contributes only an exponentially small negative off-diagonal term) and a $2^{c_p n}$-query SQ bound at constant VSTAT strength (\Cref{cor:symplpn-sq}), which extends with the same exact structure to the depolarizing-symplectic noise model that carries the external constant-rate-LPN hardness reduction of \cite{KLP+25,LPQR26} (\Cref{thm:symplpn-depol}): the very object their reductions prove at least as hard as constant-rate LPN also obeys our exact SQ lower bound, for every constant $p\in(0.317,\tfrac34)$. Inspecting that reduction further shows its depolarizing step is removable, so our independent-Bernoulli sympLPN (\Cref{def:symplpn}) \emph{itself} inherits the constant-rate-LPN hardness (\Cref{prop:bernoulli-inherits}); it is thus simultaneously LPN-hard, subject to our exact SQ bound, and covered by the linear-reduction barrier landscape. Both membership bounds rest on an exact pairwise-correlation formula
    \[
    \langle D_L, D_{L'} \rangle = \frac{(1-2p)^2}{p(1-p)} \cdot 2^{j-2n},
    \]
    where $j = \dim(L \cap L')$ and $\langle D_L, D_{L'} \rangle$ denotes the pairwise correlation. For $p=1/4$ the coefficient is exactly $4/3$. The averaging transfer (\Cref{thm:avg-transfer}) makes the average-correlation bound hold for \emph{every} subset of size $\ge|\Lagr(2n,\F_2)|/2^{2n-3}$ unconditionally (so the statistical dimension is $\ge 2^{2n-3}$); only the sharp constant of \Cref{conj:pencil} remains conjectural. The results hold for adaptive SQ algorithms.

    \item \textbf{Reduction barriers (\Cref{sec:barriers}).} We prove that linear feature maps give zero advantage against LSN; polynomial feature maps of degree $< n$ have error at least $2^{-n}$; and adaptive linear SQ algorithms have advantage at most $(1-2p)2^{-n}$ at any noise rate $p$.
    
    \textbf{Linear-reduction landscape.} Transport theorems (\Cref{thm:transport-fullrank} and \Cref{thm:transport-nearfull}) close all public-$B$ reductions; Theorem~D.1 of \cite{LPQR26} with the Fannes--Csisz\'ar inequality closes all fixed $B$; the entropy-support bound (\Cref{lem:m1}) and the reachability theorem (\Cref{thm:any-b-uniformity}) close conditionally-uniform adaptive $B$; and a support-counting bound (\Cref{thm:deterministic-marginal-adaptive}) closes the deterministic marginal-adaptive case. Within the last (randomized marginal-adaptive) cell, a confined-noise support bound (\Cref{lem:confined-g}) closes every reduction whose conditional noise occupies sub-exponentially many low-dimensional cages, and a solvable-sample-size argument (\Cref{prop:diffuse-solvable}) closes near-uniform outputs --- in particular the extreme uniform-$B$ strategy --- leaving only a concentrated-diffuse residual (\Cref{open:marginal-adaptive}). The operative barrier throughout is statistical distance: the output must stay far from a \emph{usable} LPN instance. It is \emph{not} an information bound on the secret --- we correct an earlier ``$I(x;y\mid C)=o(n)$'' target, showing instead $I(x;y\mid C)=\Theta(n)$ (the leakage is symplectic-structured, unusable by an LPN solver) --- so the open core is the distance bound \Cref{lem:m2}, not a mutual-information bound.
    
    \textbf{Two formulations.} We separate the membership and batch/sympLPN formulations and prove information-theoretic floors for the former (\Cref{sec:info-floors}).

    \item \textbf{Exact moments of the isotropic row ensemble (\Cref{sec:moments}).} For the symplectic-LPN public matrix ensemble (uniform full-rank matrices with pairwise symplectically orthogonal columns), we prove an exact closed form for \emph{every} all-ones subset moment $m_j$ ($1\le j\le 2n$). Writing $u = 2^{2n-2}$, the low-order cases give
    \[
    m_2 = \frac{(2n-1)u^2-(4n-3)u}{4(2n-1)(4u^2-5u+1)}, \qquad
    m_3 = \frac{u(u-4)}{16(4u^2-5u+1)},
    \]
    and in general $|m_j - (1/4)^j| = \Theta(4^{-n})$ for each fixed $j$. Consequently every fixed-size bundle of rows of the isotropic ensemble converges to the unconstrained LPN ensemble at rate $O(4^{-n})$: conditioning on the symplectic relation $S_A=0$ does not help a constant-size statistical test.

    \item \textbf{The standard-model gap, explicitly delineated (\Cref{sec:open}).} We isolate precisely what a hardness proof would still require --- a non-vacuous LPN(low-noise) $\to$ LSN or LWE $\to$ LSN reduction --- and classify every claim in the paper as theorem, evidence, or conjecture.
\end{enumerate}

\paragraph{Applications.} Cryptographic constructions built on LSN are developed in separate companion work --- a succinct argument system (LSN-SNARK), together with a discussion of why a key-encapsulation mechanism is deferred --- and are not treated here; the present paper concerns only the hardness landscape of the problem itself.

\section{Preliminaries}
\label{sec:prelim}

\subsection{Symplectic Vector Spaces}
\label{subsec:symplectic}

Let $V = \F_2^{2n}$ with standard symplectic form
\[
\Omega(x,y) = \sum_{i=1}^n (x_i y_{n+i} + x_{n+i} y_i),
\]
\[
\text{i.e. } \Omega(x,y) = x^{\!\top} J y \text{ with } J = \begin{psmallmatrix} 0 & I_n \\ I_n & 0 \end{psmallmatrix} \text{ (char 2)}.
\]
A subspace $L \subset V$ is \emph{isotropic} if $\Omega|_{L \times L} = 0$, and \emph{Lagrangian} if it is isotropic of dimension $n$; equivalently $L = L^{\perp_\Omega}$. We write $\Lagr(2n, \F_2)$ for the Lagrangian Grassmannian.

\begin{proposition}[Lagrangian count \cite{AG04}]
$|\Lagr(2n, \F_2)| = \prod_{i=1}^n (2^i + 1) = 2^{n(n+1)/2 + O(1)}$.
\end{proposition}

For $L, L' \in \Lagr(2n)$ let $j = \dim(L \cap L')$ denote the intersection dimension; every value $0 \le j \le n$ occurs. The distribution of $j$ is governed by the $q$-binomial coefficients at $q=2$:
\begin{equation}
\label{eq:j-distribution}
\Pr[j=k] = \frac{{n \brack k}_2 \cdot 2^{(n-k)(n-k+1)/2}}{|\Lagr(2n,\F_2)|}.
\end{equation}

\subsection{Fourier Transform on the Symplectic Space}
\label{subsec:fourier}

The unnormalized $\Omega$-Fourier transform of $f \colon \F_2^{2n} \to \R$ is
\[
\mathcal{F}_\Omega[f](\xi) = \sum_{x \in \F_2^{2n}} f(x) (-1)^{\Omega(x, \xi)}.
\]

\begin{lemma}[Self-duality]
For every $L \in \Lagr(2n,\F_2)$,
\[
\mathcal{F}_\Omega[\mathbf{1}_L] = 2^n \cdot \mathbf{1}_L.
\]
\end{lemma}

\begin{proof}
If $\xi \in L$, then $\Omega(x,\xi)=0$ for all $x \in L$, so the sum is $|L|=2^n$. If $\xi \notin L$, choose $x_0 \in L$ with $\Omega(x_0,\xi)=1$; the involution $x \mapsto x+x_0$ pairs terms of opposite sign, so the sum is $0$.
\end{proof}

\subsection{Statistical Query Model}
\label{subsec:sq}

Let $\mathcal{D}$ be a class of distributions and $D_0$ a reference distribution. The \emph{statistical dimension with average correlation} \cite{FGR+17} is
\[
\SDA(B(\mathcal{D},D_0),\gamma) = \max\Bigl\{d : \forall S \subseteq \mathcal{D}, |S| \geq \frac{|\mathcal{D}|}{d} \Rightarrow \frac{1}{|S|^2}\sum_{D,D' \in S} |\langle D,D' \rangle| \leq \gamma\Bigr\}.
\]

\begin{theorem}[Feldman \etal\ \cite{FGR+17}]
\label{thm:feldman}
If $\SDA(B(\mathcal{D},D_0),\gamma) = d$, any SQ algorithm distinguishing $\mathcal{D}$ from $D_0$ with success probability $\alpha > 1/2$ requires $q \geq (2\alpha - 1)d$ queries to $\VSTAT(1/(3\gamma))$.
\end{theorem}

The theorem applies to randomized \emph{adaptive} SQ algorithms, so our lower bound will inherit adaptivity automatically.

\subsection{Notation}
\label{subsec:notation}

We write $D_L$ for the LSN distribution with secret $L$ and noise rate $p$, and $D_0$ for the noise-only distribution in which $b \sim \operatorname{Bernoulli}(p)$ independently of $a$. All logarithms are base $2$ unless stated otherwise.


\section{Related Work}
\label{sec:related}

\paragraph{LPN foundations and cryptanalysis.}
LPN was introduced by Blum, Furst, Kearns, and Lipton \cite{BFKL94} as a foundation for cryptographic primitives. The dominant classical attack is BKW \cite{BKW03}, later optimized with fast Walsh--Hadamard transforms \cite{LF06}, covering codes \cite{GJL14}, and hybrid information-set-decoding techniques \cite{EKM17}. Lyubashevsky \cite{Lyu05} gave a polynomial-sample variant, and a comprehensive unified attack framework was developed by Esser, K{\"u}bler, and May \cite{EKM17}. Quantum speed-ups are limited to Grover \cite{Gro96} and quantum variants of BKW; no exponential improvement over the best classical methods is known.

\paragraph{LPN variants.}
Many structured variants have been proposed to improve efficiency: Ring-LPN \cite{HKL+12} and the Learning with Rounding (LWR) problem \cite{BPR12}. All of these retain a secret space of size $2^{O(n)}$. LSN differs fundamentally in that its secret space has size $2^{n^2/2 + O(n)}$ while remaining polynomial-time sampleable via symplectic linear algebra.

\paragraph{Symplectic and quantum coding.}
Stabilizer codes, the quantum analogue of classical linear codes, are naturally described by Lagrangian subspaces of $\F_2^{2n}$ \cite{Got97,CRSS98,NC00}. Quantum state tomography of stabilizer states under depolarizing noise is a closely related learning problem \cite{GL22}. The recent independent work of Khesin \etal\ \cite{KLP+25} proved that constant-rate LPN reduces to \emph{their} stabilizer-decoding LSN even for $k=1$ (the relation to our membership formulation is an open positioning item; \Cref{subsec:two-forms}), while Lu \etal\ \cite{LPQR26} proved that linear reductions cannot reduce sympLPN to LPN in the constant-rate sample regime $m=\Theta(n)$; for $m=\omega(n)$ their bound shows error weight $>1/2-\delta$, which they note is not sufficient to fully rule out decodability (though they expect the barrier to extend to all $m=\operatorname{poly}(n)$).

\paragraph{Statistical query hardness.}
The SQ model was introduced by Kearns \cite{Kea98} and developed by Blum \etal\ \cite{BFJ+94}. Feldman, Grigorescu, Reyzin, Vempala, and Xiao \cite{FGR+17} gave the general SDA framework we use, and Diakonikolas \etal\ \cite{DKS17} applied it to high-dimensional robust estimation. The known quantum approaches over classical samples (Grover search and quantum BKW) give no exponential speed-up over the best classical methods, so they do not affect the asymptotic security parameter. We do not, however, prove a quantum lower bound: the SQ bound is evidence, not proof, of quantum hardness (cf.\ \Cref{conj:quantum}).

\paragraph{Code-based cryptography.}
Code-based cryptography dates to McEliece \cite{McE78} and Prange \cite{Pra62}; modern code-based KEM candidates include Classic McEliece \cite{ABD+21}, HQC \cite{AAC+22}, and BIKE. NIST has standardized the lattice-based ML-KEM (FIPS~203, derived from CRYSTALS-Kyber) and selected HQC for additional standardization, while Classic McEliece remains a candidate rather than a FIPS KEM standard. LSN sits in this tradition --- its hardness is a decoding-flavoured assumption --- but the secret is a subspace rather than a codeword, and the public structure is symplectic rather than linear-random.

\section{The LSN Problem}
\label{sec:lsn-problem}

\begin{definition}[Membership-LSN$_{n,m,p}$]
\label{def:lsn}
Let $L \subset \F_2^{2n}$ be uniformly random Lagrangian. The LSN distribution $D_L$ over $\F_2^{2n} \times \F_2$ is generated by sampling $a \sim \F_2^{2n}$, $e \sim \operatorname{Bernoulli}(p)$, and outputting $(a, b)$ where $b = \mathbf{1}_L(a) \oplus e$. Given $m$ independent samples from $D_L$, recover $L$.
\end{definition}

\paragraph{Search vs.~decision.} All hardness results in this paper are stated for the \emph{decisional} problem: distinguishing $D_L$ from the noise-only distribution $D_0$ (where $b \sim \operatorname{Bernoulli}(p)$ independently of $a$). Search-to-decision for membership-LSN does \emph{not} follow automatically from the LPN template: the secret is a Lagrangian subspace, not a linear function, and verifying a candidate $L'$ against a decision oracle only checks consistency --- it does not by itself search the $2^{\Theta(n^2)}$-size secret space. We do not prove a search-to-decision reduction here, and treat decisional LSN as the primitive assumption.

\subsection{Two Formulations}
\label{subsec:two-forms}

Definition~\ref{def:lsn} is the \emph{membership} (oracle) formulation: the query point $a$ is freshly random for every sample, and the adversary sees it only together with its label. This is the natural form for statistical-query and information-theoretic analysis.

Cryptographic constructions may use the \emph{batch} membership formulation, where query points are fixed in advance; the external KLP/LPQR hardness results instead concern the computational sympLPN / stabilizer-decoding formulation with public isotropic matrices (and, in the multi-sample version, fresh public bases per sample).

\begin{definition}[Batch-LSN$_{n,m,p}$]
\label{def:lsn-batch}
Let $L \subset \F_2^{2n}$ be uniformly random Lagrangian, and let $A = (a_1,\dots,a_m) \in (\F_2^{2n})^m$ be a fixed sequence of query points (public parameters). The batch distribution $D_L^{(A)}$ outputs labels $b = (b_1,\dots,b_m)$ where $b_j = \mathbf{1}_L(a_j) \oplus e_j$ with independent $e_j \sim \operatorname{Bernoulli}(p)$. Given $A$ and $b$, recover $L$.
\end{definition}

The columns of $A$ may be uniform random, pseudorandom (as in applications that derive $A$ from a short seed), or structured.

\begin{definition}[sympLPN$_{n,p}$ \cite{KLP+25,LPQR26}]
\label{def:symplpn}
Let $A \in \F_2^{2n \times n}$ be a public matrix whose columns are \emph{isotropic}: $S_A := A^{\top}\Omega A = 0$ (every pair of columns is orthogonal under the symplectic form). Let $x \in \F_2^n$ be a secret vector, $e \sim \operatorname{Bernoulli}(p)^{2n}$ a noise vector, and $y = Ax + e \in \F_2^{2n}$. Given $(A,y)$, recover $x$.
This is the special case $k=n$ of LPQR26's general sympLPN$(k,n,p)$; following LPQR26, the error is drawn from the symplectic representation of the depolarizing distribution, and for simplicity of analysis we state the problem with independent Bernoulli noise. This preserves all barrier results and, as we show in \Cref{prop:bernoulli-inherits}, the Bernoulli formulation also inherits the external constant-rate-LPN hardness directly (the depolarizing step in the reduction of \cite{KLP+25} is removable).
\end{definition}

\paragraph{Relationship.}
\begin{enumerate}
\item \textbf{Batch-LSN with uniform $A$ $\equiv$ membership-LSN.} When the columns of $A$ are uniform and independent, the adversary's view $(A,b)$ in batch-LSN is identical in distribution to $m$ independent samples from the membership oracle. Hence SQ lower bounds and information-theoretic floors proved for membership-LSN transfer \emph{verbatim} to batch-LSN with uniform~$A$.
\item \textbf{Pseudorandom $A$.} Applications may generate $A$ pseudorandomly from a short seed; security then reduces to batch-LSN with uniform $A$ plus one PRG step.
\item \textbf{sympLPN is a different problem.} True sympLPN (\Cref{def:symplpn}) has linear labels $y = Ax + e$ and secret $x$; the isotropic constraint $S_A = 0$ is a public quadratic relation on the public parameters. It is not equivalent to batch-LSN (membership labels) without a pinned bridge theorem. LPQR26 study sympLPN; our transport theorems (\Cref{sec:barriers}) use only the $S_A = 0$ structure on the public matrix and therefore apply to any formulation with isotropic columns, regardless of the label structure.
\item \textbf{Open positioning item: our membership-LSN vs.\ \linebreak[0] the LSN of \cite{KLP+25,LPQR26}.} Their (classical) LSN has a \emph{public} matrix $[A \mid B]$, a junk register $x$, and a secret logical string $y$. Here the $n$ columns of $A$ are pairwise symplectically orthogonal and span a Lagrangian, and the $k$ columns of $B$ are likewise mutually isotropic and jointly independent of $A$, so the concatenated matrix is \emph{not} itself isotropic (since $n$ is the maximal isotropic dimension). Each sample is a noisy codeword $[A \mid B] \cdot [x; y] + e$ under symplectic depolarizing noise, with a fresh public matrix per sample in the multi-sample variant $\mathrm{LSN}^m$. Our membership-LSN has no public matrix and the secret is the Lagrangian itself. The two formulations are not known to be equivalent, and we do not claim the external hardness results (e.g.\ constant-rate LPN $\le$ their LSN) for our membership formulation. Establishing or refuting this bridge is an open precision item (cf.\ Open Problems, \Cref{sec:open}).
\end{enumerate}

Throughout the paper we state which formulation each theorem addresses; unless noted otherwise, ``LSN'' refers to the membership form. \emph{Computational} hardness and reduction claims --- worst-case-to-average-case obstructions, the quantum-hardness conjecture, reductions to or from LPN --- concern the \emph{computational} batch/sympLPN formulation; the membership form is statistically trivial at polynomial sample size (\Cref{sec:info-floors}) and is the object of the statistical-query and information-theoretic analysis, not of a computational hardness assumption.

\subsection{Information-Theoretic Floors}
\label{sec:info-floors}

The per-sample information content of membership-LSN is exponentially small, which makes the membership form statistically trivial at polynomial sample complexity and \emph{re-points} the hardness question to the batch formulation.

\begin{proposition}[Per-sample mutual information]
\label{prop:per-sample-mi}
For membership-LSN with a fixed noise rate $p\in(0,\tfrac12)$ bounded away from $0$ and $\tfrac12$ (the constant-noise regime, e.g.\ $p=\tfrac14$),
\[
I\bigl(L; (a,b)\bigr) \;=\; \Theta(2^{-n}).
\]
(At $p=\tfrac12$ the label is pure noise and $I=0$; the leading $\Theta(2^{-n})$ term carries the factor $\log_2\tfrac{1-p}{p}$, which vanishes there.)
\end{proposition}

\begin{proof}
Given $L$, the label is $b = \mathbf{1}_L(a) \oplus e$ with $e \sim \operatorname{Bernoulli}(p)$, so $H(b \mid a, L) = H_2(p)$ exactly.  For $a \neq 0$, a uniformly random $a$ lies in $L$ with probability $(2^n-1)/(2^{2n}-1) = 2^{-n} + O(2^{-2n})$, hence $\Pr[b=1 \mid a] = p + 2^{-n}(1-2p) + O(2^{-2n})$; for $a=0$ we have $\Pr[b=1 \mid a=0] = 1-p$, and its entropy is $H_2(1-p) = H_2(p)$, but this occurs with probability $2^{-2n}$ and contributes negligibly.  Thus $H(b \mid a) = H_2(p + 2^{-n}(1-2p)) + O(2^{-2n})$ and the mutual information is $\Theta(2^{-n})$ by the Taylor expansion $H_2(p+\varepsilon) = H_2(p) + \varepsilon\log_2\frac{1-p}{p} + O(\varepsilon^2)$.
\end{proof}

\begin{proposition}[$\chi^2$ divergence and sample complexity]
\label{prop:chi2-sample}
For any fixed Lagrangian $L$, the $\chi^2$ divergence from $D_0$ is
\[
\chi^2(D_L \| D_0) \;=\; \frac{(1-2p)^2}{p(1-p)} \cdot 2^{-n}.
\]
Consequently any distinguisher (even computationally unbounded) needs $m = \Omega(2^n)$ samples to tell $D_L$ from $D_0$ with constant advantage, and any search algorithm needs $m = \Omega(n^2 \cdot 2^n)$ samples to recover $L$ with constant probability.
\end{proposition}

\begin{proof}
From the likelihood-ratio calculation in \Cref{app:corr}, the squared likelihood ratio has expectation $\E_{D_0}[\ell_L^2] = \frac{(1-2p)^2}{p(1-p)} \cdot 2^{-n}$.  Since $\E_{D_0}[\ell_L]=0$, this is exactly $\chi^2(D_L \| D_0)$.  The standard $\chi^2$-based sample-complexity bound (Le Cam's method) gives $m = \Omega(1/\chi^2) = \Omega(2^n)$ for decision.  For search, Fano's inequality with $|\Lagr| = 2^{n(n+1)/2+O(1)}$ and per-sample information $\Theta(2^{-n})$ yields $m = \Omega(n^2 \cdot 2^n)$.
\end{proof}

\paragraph{Consequence for reductions.} Propositions~\ref{prop:per-sample-mi} and~\ref{prop:chi2-sample} imply that \emph{membership-LSN with $\operatorname{poly}(n)$ samples is information-theoretically secure}: no reduction, linear or non-linear, adaptive or static, can recover $L$ or distinguish $D_L$ from $D_0$ with only polynomially many samples.  This makes ``membership-LSN $\not\to$ LPN with $\operatorname{poly}(n)$ samples'' trivially true and cryptographically empty.  The hardness assumption relevant to applications is \emph{batch-LSN with pseudorandom $A$} (\Cref{def:lsn-batch}).  The barrier studied by LPQR26 is \emph{sympLPN} (\Cref{def:symplpn}), a decoding problem with linear labels and public isotropic matrix.

\subsection{Distance Distribution and Expected Intersection}
\label{subsec:distance}

The correlation between two LSN distributions depends only on the dimension $j = \dim(L \cap L')$ of the intersection of their secrets. The distribution of $j$ is:

\begin{theorem}[Distance distribution]
\label{thm:distance}
For uniform $L, L' \in \Lagr(2n, \F_2)$,
\[
\Pr[j=k] = \frac{{n \brack k}_2 \cdot 2^{(n-k)(n-k+1)/2}}{|\Lagr(2n,\F_2)|}.
\]
\end{theorem}

\noindent\emph{Numerical evaluation (evidence).} The moments converge rapidly: $\E[j]\to 0.7645$ and $\operatorname{Var}(j)\to 0.5963$ as $n\to\infty$ (matched to $4$ dp by $n=20$).

The rapid decay of \Cref{eq:j-distribution} means that two random Lagrangians typically intersect in dimension $0$ or $1$; intersections of dimension $\geq 14$ occur with probability $< 2^{-100}$ for every $n \geq 14$ (the tail $\Pr[j=k] \approx 2^{-k(k+1)/2}$ is essentially independent of $n$).

\subsection{The Diluted-LPN View}
\label{subsec:diluted-lpn}

In the Siegel chart $L = \{(u, Mu) : u \in \F_2^n\}$ with $M \in \operatorname{Sym}(n,\F_2)$, a sample $(a,b)$ with $a=(u,v)$ and $b=1$ reveals $n$ linear equations $v = Mu$ in the $n(n+1)/2$ unknown entries of $M$.  The catch is that true positives are drowned in noise.

\begin{proposition}[Dilution of true positives]
\label{prop:dilution}
Conditioned on the label $b=1$, the posterior probability that the query point lies in $L$ is
\[
\Pr[a \in L \mid b=1] \;=\; \frac{(1-p)\,2^{-n}}{p + (1-2p)\,2^{-n}} \;=\; \Theta(2^{-n}).
\]
At $p=1/4$ the dilution constant tends to $3$ as $n\to\infty$ (the conditional probability is $\approx 3\cdot 2^{-n}$ for large $n$); for $n=3$ it equals exactly $3/10$.
\end{proposition}

\begin{proof}
Bayes' rule: $\Pr[a \in L \mid b=1] = \Pr[b=1 \mid a \in L]\Pr[a \in L] / \Pr[b=1] = (1-p)\cdot 2^{-n} / (p + (1-2p)2^{-n})$.
\end{proof}

\Cref{prop:dilution} reframes batch-LSN search as an \emph{$\varepsilon$-diluted LPN} problem over $\operatorname{Sym}(n,\F_2)$: with probability $\varepsilon \approx 3\cdot 2^{-n}$ a sample yields $n$ clean linear equations in $M$; otherwise it yields noise.  This is LPN with secret dimension $n(n+1)/2$ at an extreme dilution rate.  The computational question becomes:

\paragraph{Enrichment question.}
\label{q:enrichment}
Can any efficient selection rule (using past labeled samples and the symplectic structure) produce query points with $\Pr[a \in L] \geq (1+\delta)\,2^{-n}$ for non-negligible $\delta$?  In other words, can the $2^{-n}$ dilution be beaten?

A negative answer would cleanly quantify how the hardness of LSN exceeds that of LPN; a positive answer would yield a concrete attack improvement, and evidence that LSN reduces to an existing assumption.  The SQ lower bounds of \Cref{sec:sq} suggest that enrichment is impossible for SQ-implementable selection rules: any such rule is a linear combination of SQ queries, and the correlation bounds limit its advantage to $2^{-n}$. The following lemma makes the reduction from enrichment to distinguishing explicit.

\begin{lemma}[Enrichment boosts SQ advantage]
\label{lem:enrichment-sq}
Let $\mathcal{R}$ be a selection rule that, using past samples from the membership-LSN oracle, outputs a query $a \in \F_2^{2n}$ such that $\Pr[a \in L] = (1+\delta)\,2^{-n}$ for some $\delta \ge 0$. Then the statistical query $\phi(a,b) = (-1)^b$ has advantage
\[
  \operatorname{Adv}(\phi) = 2|1-2p|(1+\delta)2^{-n}.
\]
In particular, the gain over the uniform-query baseline ($\delta=0$) is $\Delta\operatorname{Adv} = 2|1-2p|\cdot\delta\cdot 2^{-n}$.
\end{lemma}

\begin{proof}
Under the null model ($b=e$) we have $\E[\phi] = \E[(-1)^e] = 1-2p$. Under the structured model, $\Pr[b=1 \mid a] = p + \mathbf{1}_L(a)(1-2p)$, so $\E[\phi \mid a] = 1-2\Pr[b=1 \mid a] = (1-2p)(1-2\mathbf{1}_L(a))$. Averaging over $a \sim \mathcal{R}$ gives $\E[\phi] = (1-2p)(1-2\Pr[a \in L]) = (1-2p)(1-2(1+\delta)2^{-n})$. The advantage is $|\E[\phi]-\E_{\text{null}}[\phi]| = |1-2p| \cdot 2(1+\delta)2^{-n}$.
\end{proof}

A $\delta$-enrichment would scale the per-query advantage by $(1+\delta)$; for selection rules implementable by linear SQ this exceeds the per-query ceiling of \Cref{thm:linear-sq}, while general (non-linear, multi-sample) selection rules remain open.  Empirical probes confirm the predicted scaling $\delta \approx c\,m\,2^{-n}$ with $c=\Theta(1)$.  The exact constant depends on the prior, the sample-size regime, and the estimator; at $n=65$ and $m=22{,}528$ the enrichment is $\delta \le 2^{-50}$, far below any cryptographically relevant threshold.

\subsection{Parameters}
\label{subsec:parameters}

\Cref{tab:parameters} gives illustrative statistical-query levels from \Cref{thm:main-sq-cond} under the sharp normalization ($c=0$ and the sharp constant $5$ of \Cref{conj:pencil}); these are not concrete-security estimates. (\Cref{thm:main-sq-cond} is itself unconditional with absolute constant $24$; the sharp constant only tightens these illustrative levels.)

\begin{table}[ht]
\centering
\caption{Illustrative LSN statistical-query-exponent levels for $p=1/4$: the dimension $n$ at which the SQ-query lower bound reaches a target exponent $\lambda$ (i.e.\ $\log_2 q_{\min}\approx\lambda$). These are SQ-model query exponents, \emph{not} concrete bit-security levels.}
\label{tab:parameters}
\begin{tabular}{@{}cccc@{}}
\hline
Target exponent $\lambda$ & $n$ & $\log_2(q_{\min})^{\text{(a)}}$ & $|\Lagr|$ \\
\hline
$80$  & 41  & 80.4  & $\sim 2^{862}$  \\
$128$ & 65  & 128.4 & $\sim 2^{2146}$ \\
$192$ & 97  & 192.4 & $\sim 2^{4754}$ \\
$256$ & 129 & 256.4 & $\sim 2^{8386}$ \\
\hline
\end{tabular}
\vspace{4pt}
\noindent {\footnotesize $^{\text{(a)}}$ The $q_{\min}$ figures use \Cref{thm:main-sq-cond} under its sharp-constant normalization (\Cref{conj:pencil}); the spread bound (\Cref{thm:main-sq-uncond}) gives $\Omega(2^n)$. The \emph{unconditional} constant $24$ ($c=3$) gives $\log_2 q_{\min}\approx 2n-\log_2 24\approx 2n-4.59$, about three bits below the sharp-normalized figures shown. These figures bound \emph{statistical-query} algorithms and are evidence, not a concrete-security guarantee.}
\end{table}

The $q_{\min}$ figures follow from the \emph{sharp-normalized} bound $q \ge \tfrac13 \cdot 2^{2n-c}$ of \Cref{thm:main-sq-cond} (which is unconditional with absolute constant $24$; for these illustrative levels we instead use the sharp target constant $5$ of \Cref{conj:pencil}, normalizing its $c$ to $0$), built on the exact pairwise-correlation formula of Lemma~\ref{lem:exact-corr}: at $c=0$, $\log_2 q_{\min} = 2n - \log_2 3 \approx 2n - 1.585$, so the $\lambda=80$ line is met at $n=41$ with $0.4$ bits of slack, while an unknown $c>0$ would shift every level down uniformly. Because the multiplicative constant $\tfrac13$ costs about $1.6$ bits, parameters with $2n < \lambda + 1.6$ are not advised without a matching upper bound.

\section{Decoder Landscape}
\label{sec:decoders}

We review the natural decoder families and explain why each fails against LSN at constant noise.

\paragraph{Linear decoders.}
A real-linear query has the form $q(a,b) = \langle w, a \rangle + c \cdot b$. Its expectation under $D_L$ is
\[
\E_{D_L}[q] = c \cdot \Pr[b=1] = c \cdot \bigl(p + 2^{-n}(1-2p)\bigr),
\]
which is independent of $L$. Hence real-linear decoders have exactly zero $L$-identification power. F$_2$-linear queries can achieve advantage at most $(1-2p)2^{-n}$ (\Cref{thm:linear-sq}).

\paragraph{Polynomial decoders.}
The indicator $\mathbf{1}_L$ has an exact degree-$n$ polynomial representation
\[
\mathbf{1}_L(x) = \prod_{j=1}^n (1 + \langle w_j, x \rangle)
\]
where $\{w_j\}$ is a basis for the annihilator of $L$. After an invertible linear change of variables $\mathbf 1_L$ is the monomial $\prod_{j=1}^n z_j$ (degree exactly $n$), whose distance from $\mathrm{RM}(n-1,2n)$ is $2^{-n}$; hence every polynomial of degree $<n$ differs from $\mathbf 1_L$ on at least a $2^{-n}$ fraction of inputs.

\paragraph{List decoding.}
The list size needed to enumerate all candidate Lagrangians is $|\Lagr| = 2^{n^2/2 + O(n)}$. Even with pruning, list decoding is infeasible for $n \geq 41$.

\paragraph{BKW and ISD.}
A naive chart-linearization treating the Lagrangian as a $\Theta(n^2)$-parameter secret suggests BKW-style costs on the order of $2^{\Theta(n^2 / \log n)}$ --- attack evidence, not a proved lower bound against all BKW variants. Information-set decoding does not apply directly because the code ensemble (Lagrangian indicators) is highly structured rather than random linear.  We screened both attack families empirically: positive-basis ISD at $n=5$ succeeds in only $3/10$ trials even with $50\,000$ attempts, far above the $2^{2n}=1024$ brute-force scale; a BKW-style bucket-pair screen shows no scalable signal at $n=6$--$8$ because the label-xor noise growth $p\mapsto2p(1-p)$ kills any remaining structure after one or two rounds.

\paragraph{Structural attacks.}
For very low noise ($p = o(1/n)$), the span of the positive samples reveals $L$. At constant noise $p=1/4$, \cite{LPQR26} and our experiments (\emph{evidence}) indicate that structural attacks are ineffective: the span of positives has full dimension after $O(n)$ samples.  A direct span-of-positives decoder recovers $L$ perfectly at $p=0$ (sanity check) and fails at $p=1/4$ for every measured trial ($n=3,4,5$), because false positives span the ambient space.

\paragraph{Empirical sample-complexity lower bound.}
We ran brute-force experiments for $n=3,4,5$ using transvection-orbit enumeration of the full Lagrangian family.  In each trial we sampled $L\leftarrow\Lagr$ uniformly, drew $m$ noisy samples from $\LSN_{L,1/4}$, and let a brute-force decoder enumerate all Lagrangians, picking the one with the highest agreement.  The recovery success rates (200 trials for $n=3$, 50 for $n=4$, 20 for $n=5$) were:
\begin{center}
\begin{tabular}{@{}cccccc@{}}
\toprule
$n$ & $|\Lagr|$ & $m=2^{2n-2}$ & $m=2^{2n-1}$ & $m=2^{2n}$ & $m=2\cdot2^{2n}$ \\
\midrule
3 & 135 & $11.5\%$ & $18.5\%$ & $46.5\%$ & $80.0\%$ \\
4 & 2{,}295 & $6.0\%$ & $24.0\%$ & $62.0\%$ & $98.0\%$ \\
5 & 75{,}735 & $5.0\%$ & $20.0\%$ & $75.0\%$ & $100\%$~(20/20) \\
\bottomrule
\end{tabular}
\end{center}
Two trends are visible.  First, the success rate at the floor $m=2^{2n}$ \emph{increases} with $n$ ($46.5\%$, $62\%$, $75\%$), consistent with $2^{2n}$ being the asymptotic threshold.  Second, halving the budget to $m=2^{2n-1}$ already drops recovery below $25\%$ even at $n=5$.  These experiments empirically support an exponential-in-$n$ sample-complexity scale for the brute-force ML decoder.  A Rust cross-check at $n=5$ (compact count-aggregated scorer, 20 trials) gave consistent results: $0.25$ at $m=512$, $0.90$ at $m=1024$, $1.00$ at $m=2048$.  We note that at these small $n$ the scales $2^{2n}$ and the information-theoretic recovery floor $\Theta(n^2 2^n)$ of \Cref{prop:chi2-sample} are numerically indistinguishable, so the data do not discriminate between them; asymptotically the ML recovery threshold is governed by the $\Theta(n^2 2^n)$ floor, while $2^{2n}$ is the SQ-query scale of \Cref{sec:sq}. The experiments do not rule out a more sample-efficient decoder---the rigorous statement is the SQ query bound.

\section{Statistical Query Lower Bound}
\label{sec:sq}

This section contains the main technical result: an exponential SQ lower bound with an exact correlation formula.

\subsection{Exact Pairwise Correlation}
\label{subsec:exact-corr}

Let $D_0$ be the noise-only distribution. The likelihood ratio of $D_L$ with respect to $D_0$ is
\[
\ell_L(a,b) = \frac{dD_L}{dD_0}(a,b) - 1
            = \mathbf{1}_L(a) \cdot \beta \cdot \Bigl(\mathbf{1}_{b=1} - \mathbf{1}_{b=0} \frac{p}{1-p}\Bigr),
\]
where $\beta = (1-2p)/p$.

\begin{lemma}[Exact pairwise correlation]
\label{lem:exact-corr}
For $L,L' \in \Lagr(2n,\F_2)$ with $j = \dim(L \cap L')$,
\begin{equation}
\label{eq:exact-corr}
\langle D_L, D_{L'} \rangle
= \frac{(1-2p)^2}{p(1-p)} \cdot 2^{j-2n}.
\end{equation}
\end{lemma}

\begin{proof}
Because $a$ and $b$ are independent under $D_0$, the inner product factors:
\[
\langle D_L, D_{L'} \rangle
= \E_a[\mathbf{1}_L(a) \mathbf{1}_{L'}(a)] \cdot
  \E_b\Bigl[\beta^2 \Bigl(\mathbf{1}_{b=1} - \mathbf{1}_{b=0}\frac{p}{1-p}\Bigr)^2\Bigr].
\]
The first factor is $|L \cap L'| / 2^{2n} = 2^{j-2n}$. The second factor is
\[
\beta^2 \Bigl(p + (1-p)\frac{p^2}{(1-p)^2}\Bigr)
= \frac{(1-2p)^2}{p^2} \cdot \frac{p}{1-p}
= \frac{(1-2p)^2}{p(1-p)}.
\]
Multiplying yields \Cref{eq:exact-corr}. For $p=1/4$ the coefficient is exactly $4/3$.
\end{proof}

\subsection{Average Correlation}
\label{subsec:avg-corr}

Taking expectation over $j$ gives the average pairwise correlation.

\begin{lemma}[Average correlation]
\label{lem:avg-corr}
Let $C_n = \E[2^j]$ computed from \Cref{thm:distance} (diagonal-inclusive). Then exactly
\[
C_n = \frac{2^{n+1}}{2^n+1} \;\nearrow\; 2,
\qquad
\rho_{\mathrm{avg}}
= \frac{(1-2p)^2}{p(1-p)} \cdot C_n \cdot 2^{-2n}.
\]
(The closed form follows from \Cref{thm:distance} by the $q$-binomial theorem: $\sum_{\ell} {n \brack \ell}_2 2^{\ell(\ell-1)/2} = \prod_{i=0}^{n-1}(1+2^i)$, so $C_n = 2^n \cdot 2\prod_{i=1}^{n-1}(2^i+1) / \prod_{i=1}^{n}(2^i+1) = 2^{n+1}/(2^n+1)$.)
\end{lemma}

\paragraph{Convention (two averages).}\label{conv:two-averages}
We reserve $C_n=2^{n+1}/(2^n+1)$ for the \emph{kernel} (combinatorial) average $\E[2^{\dim(L\cap L')}]$, and write $\rho_{\mathrm{avg}}^{\mathrm{SQ}}:=\tfrac{(1-2p)^2}{p(1-p)}\,2^{-2n}C_n$ for the \emph{distributional} SQ correlation average $\rho_{\mathrm{avg}}$ of \Cref{lem:avg-corr} (the two names coincide; we write $\rho_{\mathrm{avg}}^{\mathrm{SQ}}$ only where it must be distinguished from the kernel average $C_n$). Because the correlation kernel $2^{\dim(L\cap L')}$ is scaled by the fixed factor $\tfrac{(1-2p)^2}{p(1-p)}2^{-2n}$, every subset-to-average \emph{ratio} is identical in the two normalizations; in the combinatorial big-cell and averaging arguments below $\rho_{\mathrm{avg}}$ therefore denotes $C_n$, while the SQ query theorems (\Cref{thm:main-sq-cond}) are stated in $\rho_{\mathrm{avg}}^{\mathrm{SQ}}$ (so a combinatorial $\rho(S)\le 24\,\rho_{\mathrm{avg}}=24\,C_n$ bound yields $\VSTAT(1/(72\,\rho_{\mathrm{avg}}^{\mathrm{SQ}}))$).

\subsection{Unconditional SQ Bound: The Symplectic Spread}
\label{subsec:spread-bound}

We first give an unconditional exponential lower bound on a concrete worst-case promise problem. The bound is weaker in both query count and VSTAT strength than the average-case bound of \Cref{thm:main-sq-cond}, but it is entirely elementary, resting only on an explicit symplectic spread (no association-scheme machinery).

Let $V = \F_{2^n} \times \F_{2^n}$ with symplectic form $\omega((a,b),(c,d)) = \operatorname{Tr}(ad) + \operatorname{Tr}(bc)$, where $\operatorname{Tr}$ is the trace from $\F_{2^n}$ to $\F_2$. The \emph{Desarguesian symplectic spread} $\mathcal{S}^* = \{L_\lambda\}_{\lambda \in \F_{2^n}} \cup \{L_\infty\}$ consists of $d_0 = 2^n+1$ Lagrangians
\[
L_\lambda = \{(x, \lambda x) : x \in \F_{2^n}\}, \qquad L_\infty = \{0\} \times \F_{2^n},
\]
which are pairwise transversal: $\dim(L \cap L') = 0$ for all distinct $L, L' \in \mathcal{S}^*$.

\begin{theorem}[Unconditional spread bound]
\label{thm:main-sq-uncond}
Let $\kappa = (1-2p)^2/(p(1-p))$ and fix $1 \leq t \leq n-1$. Set $\bar\gamma_t = \kappa \cdot 2^{-(2n-t)}$. Every subset $T \subseteq \mathcal{S}^*$ with $|T| \geq 2^{n-t}$ has diagonal-inclusive average correlation at most $2\bar\gamma_t$. Consequently $\SDA(\mathcal{S}^*, 2\bar\gamma_t) \geq 2^t$, and any SQ algorithm that is correct on every $L \in \mathcal{S}^*$ requires $\Omega(2^t)$ queries to $\VSTAT(1/(6\bar\gamma_t))$.
\end{theorem}

At $t = n-1$ this yields an unconditional $\Omega(2^n)$ query lower bound at VSTAT strength $O(2^n)$.

\begin{proof}
For distinct $L, L' \in \mathcal{S}^*$ we have $j = \dim(L \cap L') = 0$, so by Lemma~\ref{lem:exact-corr} the pairwise correlation is exactly $\langle D_L, D_{L'}\rangle = \kappa \cdot 2^{-2n} = \bar\gamma_t \cdot 2^{-t} \leq \bar\gamma_t$. The diagonal terms contribute $\beta = \kappa \cdot 2^{-n}$ each, so for $|T| \geq 2^{n-t}$
\[
\frac{1}{|T|^2}\sum_{D,D' \in T} |\langle D,D'\rangle|
\;\leq\; \frac{|T|^2 \cdot \bar\gamma_t + |T| \cdot \beta}{|T|^2}
\;\leq\; \bar\gamma_t + \frac{\beta}{2^{n-t}}
\;\leq\; 2\bar\gamma_t.
\]
Hence $\SDA(\mathcal{S}^*, 2\bar\gamma_t) \geq |\mathcal{S}^*| / 2^{n-t} \geq 2^t$. Applying \Cref{thm:feldman} with $\alpha = 2/3$ gives the query bound.
\end{proof}

\paragraph{Remarks.}
(i) \textbf{Worst-case promise only.} The spread has measure $|\mathcal{S}^*|/|\Lagr| \sim 2^n / 2^{n(n+1)/2} \to 0$. \Cref{thm:main-sq-uncond} does not cover the average-case problem; it protects only the promise problem ``$L \in \mathcal{S}^*$?''.
(ii) \textbf{VSTAT trade-off.} The spread bound is $\Omega(2^n)$ queries at $\VSTAT(O(2^n))$; the average-case bound below is $2^{2n-O(1)}$ queries at $\VSTAT(\Theta(2^{2n}))$. The two parameters are not interchangeable.
(iii) The subfamily restriction is valid for the promise problem: any algorithm solving LSN over the full family is correct on $\mathcal{S}^*$ a fortiori.

\subsection{The Average-Case SQ Bound}
\label{subsec:main-sq}

The stronger $2^{2n}$-scale bound is unconditional (via the big-cell theorem and the averaging transfer of \Cref{thm:avg-transfer}, yielding \Cref{thm:main-sq-cond}); only its sharp constant rests on a structural conjecture about extremal subsets.

\begin{conjecture}[Pencil extremality]
\label{conj:pencil}
There is an absolute constant $c$ such that every subset $S \subseteq \Lagr(2n,\F_2)$ of size $|S| \geq |\Lagr(2n,\F_2)| / 2^{2n-c}$ has kernel average $\rho(S) \le 5\,C_n$ (equivalently, average SQ correlation at most $5\,\rho_{\mathrm{avg}}^{\mathrm{SQ}}$).
\end{conjecture}

The pencil side of the picture is now exact. For a $k$-dimensional isotropic core $W$, the pencil $S_W = \{L : W \subseteq L\}$ is in bijection with $\Lagr(W^{\perp_\Omega}/W) \cong \Lagr(2(n-k),\F_2)$, and $\dim(L \cap L') = k + \dim(\bar L \cap \bar L')$ inside it; combining with the exact $C_m = 2^{m+1}/(2^m+1)$ of \Cref{lem:avg-corr} gives the exact pencil ratio
\[
\frac{\bar\rho(S_W)}{C_n} = \frac{2^k\,C_{n-k}}{C_n} = \frac{2^n+1}{2^{n-k}+1}
\qquad (1 \le k \le n),
\]
diagonal-inclusive, for every $n$. In particular the $k=1$ ratio tends to $2$ from below and the $k=2$ ratio equals $4 - 3/(2^{n-2}+1)$, tending to $4$ from below --- so any threshold below $4\,C_n$ is ruled out by $k=2$ pencils, exactly as the conjecture's constant anticipates. For $k \ge 3$ the ratio tends to $2^k \ge 8$, but $|S_W| = |\Lagr(2(n-k),\F_2)|$ falls below the conjectured scale $|\Lagr(2n,\F_2)|/2^{2n}$ for every $n \ge 3$, so only $k = 1, 2$ pencils are scale-relevant. A proof of the conjecture itself would require classifying near-extremal subsets beyond pencils. Exact small-case verification (evidence, not proof) supports it: at $n=2$, exhaustive enumeration of all $2^{15}$ subsets gives maximum ratio $5/2 < 5$ (attained only by singletons; the size-$3$ maximisers are exactly the $k=1$ pencils); at $n=3$, the exact maxima over all subsets of sizes $3$ and $4$ are $3$ and $81/32$ (the size-$3$ maximum attained by $k=2$ pencils; the size-$4$ value $81/32$ is attained by near-$2$-pencils (a $k=2$ pencil augmented by one further Lagrangian) \emph{and} other configurations, not by $k=2$ pencils themselves, which have only $3$ members), and pre-registered random/greedy search over larger sizes found nothing above ratio $3$; and at $n=5$ ($|\Lagr(10,\F_2)|=75735$, enumerated exactly), random/local/annealing search at the conjectured scale found nothing above ratio $\approx 2.17$, with the scale-relevant $k=2$ pencil at $11/3$ --- no subset exceeding the conjectured $5$.

\paragraph{A degree reformulation, and what the conjecture reduces to.} The average correlation has an exact \emph{degree} form. Writing $d_S(v) = |\{L \in S : v \in L\}|$ for the number of subset members through a point $v \in \F_2^{2n}$ (so $d_S(0) = M$), and using $2^{\dim(L \cap L')} = |L \cap L'| = \sum_v \mathbf 1[v \in L]\,\mathbf 1[v \in L']$,
\[
\rho(S) \;=\; \frac{1}{M^2}\sum_{v \in \F_2^{2n}} d_S(v)^2 .
\]
Two exact facts follow. First, if $\Omega(u,v) \neq 0$ then no Lagrangian contains both $u$ and $v$, so $d_S(u) + d_S(v) \le M$; hence any two points of degree $> M/2$ are $\Omega$-orthogonal and the high-degree points span an isotropic subspace. Second, the common core $W_* = \bigcap_{L \in S} L$ is isotropic and $S \subseteq S_{W_*}$, so $M \le |\Lagr(2(n-\dim W_*),\F_2)|$; since $|\Lagr(2n,\F_2)|/|\Lagr(2(n-k),\F_2)| = \prod_{i=n-k+1}^{n}(2^i+1) = 2^{kn - k(k-1)/2 + O(1)}$, the scale $M \ge |\Lagr|/2^{2n}$ forces $\dim W_* \le 2$ for $n > 3$, and quotienting by $W_*$ gives $\rho(S) = 2^{\dim W_*}\,\bar\rho(\bar S)$ with $\bar S \subseteq \Lagr(2(n-\dim W_*))$. The conjecture is therefore equivalent to a \emph{stability} statement for the symplectic dual polar space --- high average correlation forces near-containment in a $k \le 2$ pencil --- of the kind established for maximal intersecting families in dual polar spaces. Concretely, $2^r \le 4$ for $r \le 2$ and $2^r \le 8\binom{r}{3}_2$ for $r \ge 3$ (equality at $r = 3$), so with $e_S(D) = |\{L \in S : D \subseteq L\}|$ and the incidence identity $\sum_{L,L'}\binom{\dim(L \cap L')}{3}_2 = \sum_{D}e_S(D)^2$ (over the $3$-dimensional isotropic $D$; a $3$-dimensional $D \subseteq L$ is automatically isotropic), a constant-factor form of the conjecture (sufficient for the $2^{2n-O(1)}$ exponent) follows from a single \emph{supersaturation} bound,
\[
\Sigma_3(S) \;:=\; \sum_{\dim D = 3} e_S(D)^2 \;\le\; A\,M^2 \quad\text{(absolute constant $A$)} \;\Longrightarrow\; \rho(S) \le 4 + 8A :
\]
the total weight of $3$-dimensional coincidences among the members is, up to an absolute constant, no larger than for a pencil. (The implication uses $2^r \le 4 + 8\binom{r}{3}_2$ for every $r \ge 0$, with the two sides agreeing at $r = 2$; it is a \emph{sufficient} reduction --- proving the supersaturation bound with an absolute constant $A$ settles the conjecture \emph{up to that constant} (it yields the $2^{2n-O(1)}$ SQ exponent of \Cref{thm:main-sq-cond}; the literal constant $5$ would require the sharper $4+8A\le 5$ after normalization) --- that isolates the cubic incidence structure. We do not claim the converse: since $\binom{r}{3}_2 \ge 2^r/8$ for $r \ge 3$ and $\binom{r}{3}_2 = \Theta(2^{3r})$ grows strictly faster than $2^r$, a bound on the average correlation $\rho$ alone does not bound $\Sigma_3$, so $(\mathrm{SUPSAT})$ is a possibly-stronger handle rather than a literal equivalent of the conjecture.) The constant is lossy: at the evidence-level $A \approx 2$--$3$ (the worst clustered family we identify --- two non-orthogonal quotient line-pencils inside a fixed line-link, with exact $\Sigma_3/M^2 = (R_2(n{-}2)+R_1(n{-}2))/2$ where $R_t(m) = |\Lagr(2m,\F_2)|^{-1}\sum_{d} \binom{m}{d}_2 2^{\binom{m-d+1}{2}} \binom{t+d}{3}_2$ --- has the exact asymptotic $\Sigma_3/M^2 \to 92/21 \approx 4.38$, e.g.\ $403/153,\ 1891/561,\ 46609/12155$ at $n=6,7,8$) the bound certifies the ratio $\bar\rho(S)/\rho_{\mathrm{avg}} = \rho(S)/C_n$ only up to an absolute constant near $10$, not the conjectured constant $5$ of \Cref{conj:pencil} (a slack universal target: the $k=2$ pencils have ratio $4-3/(2^{n-2}+1)$, tending to $4$ from below, and $k=1$ tends to $2$, so $5$ is not a sharp finite-$n$ value); any absolute constant already suffices for the $2^{2n-O(1)}$ exponent of \Cref{thm:main-sq-cond}. Consistently, the spectral ratio bound gives only $1 + (\lambda_1/\lambda_0)(2^{2n}-1)$ with $\lambda_1/\lambda_0 = 1/(2^n+2)$ --- too weak --- as is the \emph{unconstrained} Delsarte linear program, which returns only the trivial single-vertex value $\rho = 2^n$ (ratio $(2^n+1)/2$); both omit the size constraint at the pencil scale (and the $3$-dimensional reformulation does not rescue the ratio bound --- the diagonal $\binom{n}{3}_2$ shift only worsens the resulting ratio bound). The \emph{size-constrained} Delsarte LP is sharper: with the membership normalization $a_0 = 1/M$ on the inner distance distribution $(a_i)$ --- so that every genuine $M$-subset is feasible in the program --- at the scale-relevant $k=2$-pencil size $M$ its optimum equals \emph{exactly} the $k=2$ pencil ratio $(2^n+1)/(2^{n-2}+1) = 4 - 3/(2^{n-2}+1)$ for every $n = 2,\dots,8$ (exact LP, rational: $\tfrac52,\,3,\,\tfrac{17}5,\,\tfrac{11}3,\,\tfrac{65}{17},\,\tfrac{43}{11},\,\tfrac{257}{65}$) --- a finite-$n$ exact certificate, for this fixed-size Delsarte relaxation, that the ratio stays below $4$ (hence under the slack target $5$); this is rigorous \emph{evidence} for \Cref{conj:pencil}, not a finite-$n$ proof of the full large-subset statement. Nor does it extend to a uniform-in-$n$ proof: the optimal Delsarte dual needs \emph{all} $n+1$ MacWilliams levels --- keeping only the lowest few (e.g.\ levels $j \le 3$) blows the certified bound up to $\approx 2^{n-2}$ ($\approx 16.3,\,32.3,\,64.3$ at $n = 6,7,8$, against true ratios still below $4$) --- so within this Delsarte-dual route the certificate is all-or-nothing, and the all-$n$ bound remains open. In fact the correlation operator $\sum_i 2^{n-i} A_i$ ($A_i$ the relation $\dim(L\cap L')=n-i$) equals the Gram matrix $\Phi\Phi^{\top}$ of the point--Lagrangian incidence $\Phi_{L,x}=\mathbf 1[x\in L]$ ($x\in\F_2^{2n}$; the $2^{2n}-1$ nonzero $x$ are the isotropic points of the polar space, the form being alternating over $\F_2$), since $2^{\dim(L\cap L')}=|L\cap L'|=\sum_x \mathbf 1[x\in L]\,\mathbf 1[x\in L']$; its range is spanned by the point-pencil indicators $\mathbf 1_{\{L\,:\,x\in L\}}$ (over the nonzero $x$; the $x=0$ incidence is the all-ones vector), an $\mathrm{Sp}(2n,2)$-submodule meeting only the two bottom cometric levels $V_0\oplus V_1$, so $\lambda_d=0$ for all $d\ge 2$ and all $n$. (Equivalently, the dual polar eigenvalues~\cite{BCN89} give $\lambda_d=\sum_i 2^{n-i}P_i(d)$ with a $q$-Pochhammer factor $\prod_{u=0}^{d-1}(1-2^{u-1})$ that vanishes for $d\ge 2$, since its $u=1$ term is $1-2^0=0$; this leaves $\lambda_0=2^{n+1}|\Lagr(2n,\F_2)|/(2^n+1)$ and $\lambda_1=\lambda_0/(2^n+2)$, verified exactly for $n\le 4$, where the spectra of $\sum_i 2^{n-i}A_i$ are $\{24,4,0\}$, $\{240,24,0,0\}$, and $\{4320,240,0,0,0\}$.) Consequently $\rho(S)=\rho_{\mathrm{avg}}+\lambda_1\,\lVert E_1\mathbf 1_S\rVert^2/M^2$ exactly ($E_1$ the projection onto $V_1$), and \Cref{conj:pencil} is \emph{equivalent} to the single first-level bound $\lVert E_1\mathbf 1_S\rVert^2 \le 4(\lambda_0/\lambda_1)\,M^2/|\Lagr(2n,\F_2)|$ at every scale $M\ge|\Lagr(2n,\F_2)|/2^{2n-c}$: only the degree-one case of a global-functions inequality is at stake, not the full level-$d$ tower, the task being to upgrade the trivial $\lVert E_1\mathbf 1_S\rVert^2\le M$ to a bound quadratic in $M$ (first-eigenspace, or EKR, stability). The point ($k{=}1$) pencils lie entirely in $V_0\oplus V_1$; the $k{=}2$ pencils are the scale-relevant extremal examples, and although a $2$-pencil indicator (the intersection of the two point pencils of a basis of the isotropic $2$-space) carries higher-level components, these are annihilated by the correlation operator, so only its $V_0\oplus V_1$ part enters $\rho$. The natural route to that first-level bound is the global-hypercontractivity framework for stability of association schemes~\cite{KLLM2024}: decompose $\mathbf 1_S$ into a global part, where the density-$2^{-2n}$ regime keeps the hypercontractive penalty $O(n)$ and the scheme's rapid eigenvalue decay absorbs it, and a junta part supported on a bounded union of pencils, each contributing $O(1)$ to $\rho$. Establishing the requisite degree-one hypercontractive inequality (the only level entering $\rho$) for the symplectic dual polar scheme $\mathrm{DSp}(2n,2)$ (known for the Johnson, symmetric-group, and Grassmann schemes, but not yet for dual polar spaces) and the matching junta-stability lemma would settle the conjecture; we leave this open. We note one obstruction specific to our regime: the high-dimensional-expander route to such an inequality is weakest at $q=2$. By the local-to-global (trickling-down) principle the local-spectral expansion of the type-$C_n$ dual polar building is controlled by its rank-$2$ residues, the generalized quadrangles $\mathrm{GQ}(q,q)$; for $\mathrm{GQ}(2,2)$ this expansion is a large constant ($\lambda \approx 1/2$ on the collinearity graph, $\approx 2/3$ on the incidence graph), decaying to $0$ only as $q \to \infty$, whereas the hypercontractivity theorems require $\lambda$ small. Good local expansion of symplectic dual polar spaces is thus a large-$q$ phenomenon, and $q=2$ is precisely the worst case. This rules out only the \emph{local-spectral} (random-walk) hypercontractivity of high-dimensional expanders; it does \emph{not} rule out the \emph{global-functions} hypercontractivity of Keevash--Lifshitz--Long--Minzer~\cite{KLLM2024}, whose global-functions framework (unlike the local-spectral route) does not rest on a local-expansion hypothesis, exploiting the intersection semilattice rather than local random walks; adapting it here would require a corresponding degree-one global inequality for $\mathrm{DSp}(2n,2)$, which we do not prove. The live route at $q=2$ is therefore to extend that global-functions framework from the Grassmann to the symplectic dual polar scheme, or to argue through the symplectic-Fourier self-duality $\mathcal F_\Omega[\mathbf 1_L] = 2^n \mathbf 1_L$ directly. An exact \emph{link recursion} isolates the combinatorial content. For an isotropic $2$-space $W$ write $m_W = |\{L \in S : W \le L\}|$ and let $S_W$ be the family of links $L/W$ in the rank-$(n-2)$ symplectic quotient $W^{\perp}/W$; since every $3$-dimensional space contains exactly $\binom{3}{2}_2 = 7$ two-dimensional subspaces,
\[
7\,\Sigma_3(S) \;=\; \sum_{\dim W = 2,\ W \text{ iso}} \bigl(\rho(S_W) - 1\bigr)\, m_W^2 .
\]
Thus (SUPSAT) is equivalent to a \emph{$2$-pencil packing} bound $\sum_W (\rho(S_W)-1)\,m_W^2 = O(M^2)$: the $2$-pencil shadows must pack with square-summable total mass. This is \emph{not} a clean global-versus-junta dichotomy. A near-$2$-pencil $\mathcal P_W \cup \{L_0\}$ --- the full pencil $\mathcal P_W = \{L : W \le L\}$ plus one Lagrangian $L_0$ transverse to $W$ --- has trivial common core ($W \cap L_0 = 0$) yet keeps $1-o(1)$ of its members in a single pencil, with $\Sigma_3 \asymp M^2$ contributed by the $2^{2n-4}$ shadows $D \supseteq W$; so the pencil contribution is genuine and any proof must allow it. The heavy-codegree threshold is moreover vacuous at scale: $e_S(D) \le |\Lagr(2(n-3),\F_2)| < M\,2^{3-n}$ for every $D$, so the entire mass lives at the $M\,2^{-n}$ level and a robust core controlling only $e_S(D) > M/2$ sees none of it. Because the links $S_W$ need not inherit the density or trivial-core hypotheses, the recursion is not a naive induction --- at the pencil scale $M \ge |\Lagr(2n,\F_2)|/2^{2n-c}$ the average point-link $S_p$ has relative size exactly the rank-$(n-1)$ threshold with $c$ decreased by $2$ (since $\mathbb E_p|S_p| = M/(2^n+1) = |\Lagr(2(n-1),\F_2)|/2^{2(n-1)-(c-2)}$), so any single relative-density parameter degrades at a fixed rate down the tower and the link hypothesis is exhausted after $O(c)$ descent steps; the square-summable packing bound is the precise open content, and the first-moment-times-maximum estimate misses it by the factor $2^{2n}$. Equivalently, in the point-shadow picture --- to which the symplectic-Fourier self-duality $\mathcal F_\Omega[\mathbf 1_L]=2^n\mathbf 1_L$ reduces, since it returns the identity $\rho(S)=1+M^{-2}\sum_{x\neq 0}d_S(x)^2$ rather than a separate handle --- the conjecture is a Kruskal--Katona-type \emph{support-expansion} bound: a family with no large $1$- or $2$-pencil must spread its point support beyond the $2^{2n-2}$ points of a single $2$-pencil. The point-support bound is tight --- uniform mass over exactly those $2^{2n-2}$ points gives $\sum_{x\neq 0}d_S(x)^2\to 4M^2$, hence $\rho\to 5$ (the slack target of \Cref{conj:pencil}; the $k=2$ pencils themselves reach only $4$) --- but a shadow lower bound on the point support does \emph{not} by itself suffice, because these scalar and self-dual constraints are blind to multiplicity. For $n\ge3$, fixed $\varepsilon>0$, and $M\ge(1+\varepsilon)M_2$, a $2$-pencil $S_W$ of size $M_2=|S_W|$, and a Lagrangian $L_0$ transverse to $W$, the nonnegative self-dual function $f=(M-M_2)\,\mathbf 1_{L_0}+d_{S_W}$ satisfies all of them --- $f(0)=M$, $\sum_x f(x)=M2^n$, $\mathcal F_\Omega f=2^n f$, $\operatorname{supp}(f)\supseteq W^{\perp}$, and $\max_{x\neq 0}f(x)<M$ (transversality of $L_0$ bounds every nonzero degree by $\max\{M_2,\,M-M_2+|\Lagr(2(n-3),\F_2)|\}<M$) --- yet $\sum_x f(x)^2\ge(M-M_2)^2\,2^n$, so $\rho(f)\ge(1-M_2/M)^2\,2^n=\Theta(2^n)$ and both $\rho(f)$ and the ratio $\rho(f)/\rho_{\mathrm{avg}}$ diverge. It is the degree function of a \emph{multiset} ($M-M_2$ copies of $L_0$ together with the pencil), not of any genuine subset. The missing input is \emph{distinctness}: a genuine $M$-subset obeys the point-projection bound $\sum_{x\in L'}d_S(x)=\sum_{L\in S}2^{\dim(L\cap L')}\le 2^n+(M-1)2^{n-1}$ for every Lagrangian $L'$ ($f$ certainly exceeds it once $M>2M_2$). More uniformly, $f$ is excluded already at the Delsarte diagonal level: the size-constrained program enforces $a_0=1/M$, and a multiset of multiplicities $k_i$ has $a_0=\sum_i k_i^2/M^2\ge 1/M$ with equality \emph{iff} the elements are distinct, whereas $f$ has $a_0=((M-M_2)^2+M_2)/M^2>1/M$ once its multiplicity $M-M_2\ge 2$ (automatic for large $n$). This is why those scalar/support/self-dual constraints \emph{alone} cannot prove \Cref{conj:pencil}, even though the size-constrained program is tight for $n\le 8$.

\paragraph{A symmetric-bilinear-forms reduction, and what it isolates.} A different reduction recovers a genuine product structure that the dual polar scheme lacks. Polarising $\F_2^{2n}=X\oplus Y$ with $Y$ a fixed Lagrangian, every Lagrangian transverse to $Y$ is the graph $L_A=\{(x,Ax):x\in X\}$ of a unique symmetric matrix $A\in\mathrm{Sym}_n(\F_2)$, and $\dim(L_A\cap L_B)=n-\mathrm{rank}(A-B)$. On this ``big cell'' --- a fraction $\prod_{i\ge 1}(1+2^{-i})^{-1}\approx 0.42$ of $\Lagr(2n,\F_2)$ --- the average correlation becomes the translation-invariant convolution $\rho(S)=|S|^{-2}\sum_{A,B\in S}2^{n-\mathrm{rank}(A-B)}=\sum_{x\in\F_2^n}\mathrm{CP}(Ax)$ ($\mathrm{CP}$ the collision probability of the evaluation $A\mapsto Ax$) on the abelian group $\mathrm{Sym}_n(\F_2)=\F_2^{n(n+1)/2}$, i.e.\ on the \emph{symmetric bilinear forms} translation association scheme~\cite{Schmidt2018}, which carries a full Fourier structure. The Fourier transform of the kernel $2^{n-\mathrm{rank}\,C}$ is supported only on the rank-$\le 2$ quadratic-form characters (squares $\ell^2$ and hyperbolic products $\ell m$, $\ell\neq m$), so \Cref{conj:pencil} on the big cell is \emph{exactly} the level-$2$ square-function bound $\sum_{\ell\neq 0}|\widehat{\mathbf 1_S}(\ell^2)|^2+2\sum_{\ell<m}|\widehat{\mathbf 1_S}(\ell m)|^2\le B_c\,2^n\alpha^2$ at density $\alpha=|S|/2^{n(n+1)/2}\ge 2^{-2n+c}$; over $\F_2$ the square characters $\ell^2=\ell$ are linear, so the degree-one part is automatically $\le 2^n\alpha^2$ and only the hyperbolic (genuinely degree-two) part is at stake. Indeed, writing $r=\mathrm{rank}(C)$, the hyperbolic character sum $H_n(C)=\sum_{\ell<m}(-1)^{\ell^{\top}Cm}$ equals $2^{2n-r-1}-3\cdot 2^{n-1}+1$ when $\mathrm{diag}(C)=0$ and $2^{2n-r-1}-2^n+1$ otherwise, so the degree-two mass admits the exact closed form $W_2(\mathbf 1_S)/(2^n\alpha^2)=\rho(S)-p_{\mathrm{alt}}(S)+(2^{1-n}-2)$, where $p_{\mathrm{alt}}(S)=\Pr_{A,B\in S}[\mathrm{diag}(A-B)=0]$ is the diagonal-collision probability; the diagonal (alternating) sector thus enters only as a \emph{subtracted} term and cannot inflate the degree-two mass. This is the \emph{symmetric} analogue of the global-hypercontractive (Bonami) inequality for functions on spaces of linear maps of Ellis--Kindler--Lifshitz~\cite{EKL2022}; it does \emph{not} follow from their theorem by restriction --- $\mathrm{Sym}_n$ has measure $2^{-n(n-1)/2}$ inside all $n\times n$ matrices, vacuating the ambient bound, and exposing a row also exposes the matching column --- and we leave it open. A random choice of $Y$ transfers the big-cell estimate to all of $\Lagr(2n,\F_2)$ with only an absolute-constant loss --- \emph{unconditionally}, with no concentration or global/junta hypothesis --- by the following augmentation argument. Write $S_Y=S\cap\{L:L\cap Y=0\}$ for the transverse slice and $s=2^{\,n(n+1)/2-2n+2}$ for the big-cell threshold size. The probability that both members of a pair are transverse to $Y$ depends only on $j=\dim(L\cap L')$: in a $W$-adapted symplectic basis with $L=\langle e_1,\dots,e_n\rangle$ and $L'=\langle e_1,\dots,e_j,f_{j+1},\dots,f_n\rangle$ (so $L\cap L'=\langle e_1,\dots,e_j\rangle$), a Lagrangian transverse to $L$ is the graph of a symmetric $A$, and $Y\cap L'=0$ iff the bottom-right $(n-j)\times(n-j)$ block of $A$ is nonsingular; hence $\pi(j):=\Pr_Y[L\cap Y=0=L'\cap Y]=p_{\mathrm{tr}}\prod_{i=1}^{\lceil(n-j)/2\rceil}(1-2^{-(2i-1)})$, where $p_{\mathrm{tr}}:=2^{n(n+1)/2}/|\Lagr(2n,\F_2)|$ is the single-transversality probability (not to be confused with the noise rate $p$), with $\pi_{\max}=\pi(n)=p_{\mathrm{tr}}$ and $\pi_{\max}/\pi_{\min}\le\prod_{i\ge1}(1-2^{-(2i-1)})^{-1}<\tfrac{12}{5}$. The exact moment identities $\mathbb E_Y\big[\sum_{L,L'\in S_Y}2^{\dim(L\cap L')}\big]=\sum_{L,L'\in S}\pi(j)\,2^{\dim(L\cap L')}$ and $\mathbb E_Y|S_Y|^2=\sum_{L,L'\in S}\pi(j)$ hold, with $\mathbb E_Y|S_Y|=|S|p_{\mathrm{tr}}$. For \emph{every} $Y$ the slice sum is capped at the fixed scale $s^2$: when $|S_Y|<s$, embed $S_Y$ in any distinct $s$-element subset $T\subseteq\{L:L\cap Y=0\}$ (possible since $2^{n(n+1)/2}\ge s$), apply the big-cell bound (\Cref{thm:bigcell}) to $T$, and discard the added nonnegative terms, giving $\sum_{L,L'\in S_Y}2^{\dim(L\cap L')}\le\sum_{L,L'\in T}2^{\dim(L\cap L')}\le5\,C_ns^2$; when $|S_Y|\ge s$ the big-cell theorem applies to $S_Y$ directly, $\sum_{L,L'\in S_Y}2^{\dim(L\cap L')}\le5\,C_n|S_Y|^2$. (The cap $\rho_{BC}(T)\le5\,C_n$ holds for every $n$, and we do \emph{not} rely on the limiting ratio at finite $n$: the displayed Delsarte certificate has value $<37/4$ (its ratio against the Schmidt-cell average tends to $37/8<5$ and is verified through $n=32$), while $5\,C_n=5\cdot2^{n+1}/(2^n+1)\ge160/17>37/4$ for all $n\ge4$ as the Lagrangian average increases in $n$; the small cases $n\le3$ are checked exhaustively, with worst average kernel $3$ and $9/2$ against $5\,C_n=8$ and $80/9$.) One-sided Chebyshev bounds the frequency of thin slices, $\Pr_Y[|S_Y|<s]\le\mathbb E_Y|S_Y|^2/(|S|p_{\mathrm{tr}}-s)^2$, so with $k:=|S|p_{\mathrm{tr}}/s=|S|\,2^{2n-2}/|\Lagr(2n,\F_2)|$ one obtains $\mathbb E_Y[\sum_{L,L'\in S_Y}2^{\dim(L\cap L')}]\le5\,C_n\,\mathbb E_Y|S_Y|^2\,(1+(k-1)^{-2})$. Dividing through by the $\pi$-weights ($\sum_S\pi(j)2^{\dim}\ge\pi_{\min}\sum_S2^{\dim}$ and $\sum_S\pi(j)\le\pi_{\max}|S|^2$) gives $\rho(S)\le(\pi_{\max}/\pi_{\min})\,5\,C_n\,(1+(k-1)^{-2})$; for $|S|\ge|\Lagr(2n,\F_2)|/2^{\,2n-c}$ with $c=3$ one has $k\ge2$, hence (using $\pi_{\max}/\pi_{\min}<\tfrac{12}{5}$ and $1+(k-1)^{-2}\le2$, so the constant is $<\tfrac{12}{5}\cdot5\cdot2=24$)
\[
  \rho(S)\ \le\ 24\,C_n\qquad\text{for every }S\text{ with }|S|\ge|\Lagr(2n,\F_2)|/2^{\,2n-3},
\]
an absolute constant for all $n$. (Throughout this argument $\rho_{\mathrm{avg}}$ denotes the $K$-average $C_n=2^{n+1}/(2^n+1)$, consistent with the identity $\rho(S)/\rho_{\mathrm{avg}}=\rho(S)/C_n$ used above; the conclusion is a normalization-independent \emph{ratio}, so the same bound $\rho(S)\le24\,\rho_{\mathrm{avg}}$ applies verbatim to the distributional SQ average $\rho_{\mathrm{avg}}^{\mathrm{SQ}}=\tfrac{(1-2p)^2}{p(1-p)}\,2^{-2n}C_n$ entering \Cref{thm:main-sq-cond}.) The cap crucially uses the \emph{distinctness} of $S_Y$ --- a multiset could not be embedded in a density-feasible distinct subset --- which is exactly why the multiplicity-blind scalar/self-dual obstructions noted above are evaded. The transversality count $\pi(j)$ is verified exhaustively for $n\le5$ (all $j$), the cap and the full chain Track-A for $n\le4$ on worst-case $S$. The evidence is favourable: exhaustive ($n=2$) and structured ($n\le 5$) search finds no subset above ratio $\approx 7$, and the slice-rank/polynomial-method route is provably vacuous here (the $2$-pencil $\{A:AU=0\}$, $\dim U=2$, is already a clique for $\mathrm{rank}\le n-2$, forcing any density-saving rank-difference count to be exponentially weak). The conjecture thus reduces, up to this verified averaging step, to the level-two ($d=2$) case of the \emph{symmetric} analogue of the Ellis--Kindler--Lifshitz restriction-global hypercontractive inequality, and that adaptation is largely set up. Symmetric $r$-restrictions are indexed by a \emph{single} subspace $R\le\F_2^n$ (the fibre $\{A:A|_R=T\}$ is an affine copy of $\mathrm{Sym}(\F_2^n/R)$, so $\ell^2$ and $\ell m$ are read off from $R=\langle\ell\rangle$ and $R=\langle\ell,m\rangle$); a one-line recursion --- verified exactly for $n\le 4$ --- bounds the level-two mass $W_2(\mathbf 1_S)$ on $\mathrm{Sym}_n$ by a local line term $L_e\le(2^{n+1}-3)\alpha^2$ plus $2\,\mathbb E_y\,W_2(\mathbf 1_{S_y})$ over the $\mathrm{Sym}_{n-1}$ line-fibres $S_y$; the diagonal/local part sits at the allowed $2^n\alpha^2$ scale; and the hypercontractive constant $q^{Cd^2}$ is an absolute $2^{O(1)}$ at $d=2$ (so it does not grow with $n$). The single missing input for \emph{that} route is the symmetric generalised-influence (rank/type-lowering) estimate for Schmidt's eigenspaces --- the analogue, for the symmetric scheme, of the globalness step that powers the proved linear-maps case~\cite{EKL2022}. We do not need it: the big-cell bound \emph{itself} is unconditional for every $n$, by an explicit size-constrained Delsarte certificate on Schmidt's scheme that bypasses the hypercontractive estimate entirely. The kernel's Fourier support being exactly the squares $\ell^2$ and the hyperbolic products $\ell m$, one has $\rho(S)=\rho_{\mathrm{avg}}+(W_1+2W_2)/(2^n\,|S|^2)$ with $W_1=\sum_{\ell\neq0}|\widehat{\mathbf 1_S}(\ell^2)|^2$ and $W_2=\sum_{\ell<m}|\widehat{\mathbf 1_S}(\ell m)|^2$, so the size-constrained Delsarte program on the rank--type refinement of $\mathrm{Sym}_n(\F_2)$ at the membership density $\alpha\ge 2^{-2n+2}$ upper-bounds $\rho(S)$ by its optimum. That optimum carries an explicit \emph{constant} dual certificate: with $\eta_i\equiv -5\cdot 2^{-2n}$ on the eigenspace support $\{(4,-)\}\cup\{(2a,\pm):a\ge 3\}\cup\{(2b{+}1,0):b\ge 2\}$ and $\lambda=8$, the dual slack $8-5\cdot 2^{-2n}\,g(k)-2^{\,n-\operatorname{rank}(k)}$ at each class $k$ --- where $g(k)=\sum_{i\in S}P^{(n)}[i,k]=-1-\sum_{i\in\mathrm{low}}P^{(n)}[i,k]$ is the sum of the $O(1)$ low-rank eigenvalues in closed form --- reduces to five scalar quadratics in $y=4^{\lfloor n/2\rfloor-r}$ (such as $5y^2-117y+636$ and $10y^2-117y+318$), each positive on every power of four (their root intervals are skipped by the powers of four), so the certificate is dual-feasible for \emph{all} $n$, with value $\to 37/4$ and hence ratio $\rho(S)/\rho_{\mathrm{avg}}\to 37/8=4.625<5$. We record the two results established by the argument above as labeled theorems.

\begin{theorem}[Big-cell square-function bound]\label{thm:bigcell}
At membership density $\alpha=|S|/2^{n(n+1)/2}\ge 2^{-2n+2}$ on the big cell, every subset satisfies the level-two square-function bound displayed above, equivalently $\rho_{BC}(S)\le 5\,C_n$. The explicit constant dual certificate ($\eta_i\equiv-5\cdot2^{-2n}$, $\lambda=8$) is dual-feasible for every $n$ --- its slack reduces to five scalar quadratics positive on every power of four, and we verified it exactly against the full dual-feasibility constraints, including the membership row, through $n=32$ (\Cref{app:bigcell-cert} records the five quadratics, the support, and the feasibility reduction explicitly).
\end{theorem}

\begin{theorem}[Averaging transfer]\label{thm:avg-transfer}
For every $S\subseteq\Lagr(2n,\F_2)$ with $|S|\ge|\Lagr(2n,\F_2)|/2^{\,2n-3}$, the augmentation argument above transfers \Cref{thm:bigcell} to all Lagrangians with an absolute-constant loss: $\rho(S)\le 24\,C_n$ (equivalently average SQ correlation at most $24\,\rho_{\mathrm{avg}}^{\mathrm{SQ}}$), unconditionally --- with no concentration or global/junta hypothesis. Consequently \Cref{thm:main-sq-cond} is unconditional.
\end{theorem}

\begin{theorem}[Main SQ lower bound]
\label{thm:main-sq-cond}
Any SQ algorithm distinguishing LSN (uniform $L$) from $D_0$ with probability $> 2/3$ requires $q \geq 2^{2n - O(1)}$ queries to $\VSTAT(1/(72\,\rho_{\mathrm{avg}}^{\mathrm{SQ}}))$.
\end{theorem}

\begin{proof}
By the averaging transfer (\Cref{thm:avg-transfer}, the augmentation argument with $c=3$), $\rho(S)\le 24\,C_n$ --- equivalently average SQ correlation at most $24\,\rho_{\mathrm{avg}}^{\mathrm{SQ}}$ by \Cref{conv:two-averages} --- so every subset of size at least $|\Lagr(2n,\F_2)| / 2^{2n-3}$ has SQ correlation at most $24\,\rho_{\mathrm{avg}}^{\mathrm{SQ}}$, so $\SDA(B(\{D_L\},D_0), 24\,\rho_{\mathrm{avg}}^{\mathrm{SQ}}) \geq 2^{2n-3}$. Applying \Cref{thm:feldman} with $\alpha = 2/3$ and $\gamma = 24\,\rho_{\mathrm{avg}}^{\mathrm{SQ}}$ yields $q \geq (4/3 - 1) \cdot 2^{2n-3} = 2^{2n - O(1)}$ queries to $\VSTAT(1/(72\,\rho_{\mathrm{avg}}^{\mathrm{SQ}}))$. (Conditionally on the sharp constant $5$ of \Cref{conj:pencil}, the bound improves to $5\,\rho_{\mathrm{avg}}^{\mathrm{SQ}}$ and $\VSTAT(1/(15\,\rho_{\mathrm{avg}}^{\mathrm{SQ}}))$.)
\end{proof}

\begin{theorem}[Adaptive SQ]
\Cref{thm:main-sq-cond} holds for randomized adaptive SQ algorithms.
\end{theorem}

\begin{proof}
\Cref{thm:feldman} applies to randomized adaptive SQ algorithms \cite[Theorem 3.7]{FGR+17}, so the bound is inherited directly.
\end{proof}

\subsection{Adaptive Linear SQ is Blocked}
\label{subsec:adaptive-linear}

\begin{theorem}[Adaptive linear SQ barrier]
\label{thm:linear-sq}
For any real-linear query $q(a,b) = \langle w, a \rangle + c \cdot b$, the expectation $\E_{D_L}[q]$ is independent of $L$. For any F$_2$-linear query $q(a,b)=\langle w,a\rangle \oplus c\cdot b$, the maximum distinguishing advantage against $D_0$ is $(1-2p)\,2^{-n}$; at $p=1/4$ this is $2^{-n-1}$.
\end{theorem}

\begin{proof}
For real-linear queries, direct computation gives that $\E_a[\langle w,a\rangle]$ is independent of $L$ (it equals $1/2$ for $w \neq 0$ and $0$ for $w=0$) and $\E_{D_L}[b]=p+2^{-n}(1-2p)$, which is independent of $L$. For an F$_2$-linear query $q(a,b)=\langle w,a\rangle\oplus c\cdot b$: if $c=0$ and $w\neq 0$ then $\E_{D_L}[q]=1/2$ for every $L$ (the constant query $w=0,c=0$ is trivially $L$-independent); if $c=1$ and $w=0$, then $q=b$ and the advantage over $D_0$ is exactly $(1-2p)2^{-n}$. If $c=1$ and $w\neq 0$, write $q=b\oplus\langle w,a\rangle$. Under $D_L$ we have
\[
\E_{D_L}[q] = \frac12 + (1-2p)\bigl(2^{-n}-2\E_a[\langle w,a\rangle\cdot\mathbf 1_L(a)]\bigr).
\]
If $w\in L^\perp$ then $\E_a[\langle w,a\rangle\cdot\mathbf 1_L(a)]=0$ and the advantage is $(1-2p)2^{-n}$. If $w\notin L^\perp$ then $\E_a[\langle w,a\rangle\cdot\mathbf 1_L(a)]=2^{-n-1}$ and the expectation collapses to $1/2$, giving advantage $0$. Thus no F$_2$-linear query has advantage larger than $(1-2p)2^{-n}$.
\end{proof}

\subsection{The sympLPN Formulation: Exact Correlations and an SQ Bound}
\label{subsec:symplpn-sq}

The bounds above concern the membership formulation. The same character-sum technique gives exact pairwise correlations --- and an SQ lower bound --- for the \emph{sympLPN} formulation itself, whose secret space is only $\F_2^n$, so no bound for it can exceed $2^n$. Here one sample is the full block $(A, y)$ with $A \sim \mathcal{A}_n$, $y = Ax + e$, $e \sim \mathrm{Bernoulli}(p)^{2n}$; the null $D_0$ draws $y$ uniform, independent of $A$.

\begin{theorem}[Exact sympLPN correlations]
\label{thm:symplpn-corr}
Let $\tau = (1-2p)^2$ and $\langle D_x, D_{x'} \rangle = \E_{D_0}[(\tfrac{D_x}{D_0}-1)(\tfrac{D_{x'}}{D_0}-1)]$. Under the isotropic ensemble,
\[
\langle D_x, D_x \rangle = (1+\tau)^{2n} - 1, \qquad
\langle D_x, D_{x'} \rangle = -\frac{(1+\tau)^{2n}-1}{2^{2n}-1} \quad (x \neq x').
\]
Moreover, for every non-empty $S \subseteq [2n]$, the bundle-parity query $h^{(S)}(A,y) = (-1)^{\sum_{i \in S} y_i}$ has advantage function $g_x(A) = \E_{D_x \mid A}[h^{(S)}] = (-1)^{\langle \mathbf 1_S, Ax\rangle}(1-2p)^{|S|}$ with $\E_A[g_x g_{x'}] = (1-2p)^{2|S|}$ if $x = x'$ and $-(1-2p)^{2|S|}/(2^{2n}-1)$ otherwise.
\end{theorem}

\begin{proof}
Conditioned on $A$ the likelihood ratio factorizes, $D_x(y \mid A)/D_0(y \mid A) = \prod_{i=1}^{2n} \bigl(1 + (1-2p)(-1)^{y_i + \langle a_i, x \rangle}\bigr)$, and averaging the product for the pair $(x,x')$ over $y \sim D_0$ leaves $\prod_{i=1}^{2n}(1 + \tau(-1)^{v_i})$ with $v = A(x - x')$. For $x = x'$ this is $(1+\tau)^{2n}$. For $x \neq x'$, expand $\prod_i (1 + \tau (-1)^{v_i}) = \sum_{S} \tau^{|S|} (-1)^{\langle \mathbf 1_S, v \rangle}$; the vector $v = Aw$ ($w = x - x' \neq 0$) is uniform over $L \setminus \{0\}$ for $L = \operatorname{colspan}(A)$. For a fixed Lagrangian the character sum is
\[
\sum_{v \in L} (-1)^{\langle \mathbf 1_S, v \rangle} = 2^n \cdot \mathbf 1_{\{J \mathbf 1_S \in L\}},
\]
where $J$ is the Gram matrix of $\Omega$ (the coordinate involution $i \mapsto i \pm n$): the Euclidean dual of $L$ is $J L^{\perp_\Omega} = JL$, and the twist by $J$ matters --- e.g.\ $L = \operatorname{span}\{e_1, e_2\}$ with $S = \{1\}$ has sum $0$, not $2^n$. Since $J \mathbf 1_S \neq 0$, $\Pr_L[J \mathbf 1_S \in L] = 1/(2^n+1)$ by $\mathrm{Sp}$-transitivity, whence for every non-empty $S$
\[
\E_A\bigl[(-1)^{\langle \mathbf 1_S, Aw \rangle}\bigr] = \frac{1}{2^n - 1}\Bigl(\frac{2^n}{2^n+1} - 1\Bigr) = -\frac{1}{2^{2n}-1}.
\]
Summing the expansion gives the off-diagonal value; the bundle-parity statement is the single-$S$ case. All values were verified by exhaustive enumeration of the full ensemble at $n = 2, 3$.
\end{proof}

Two readings. First, this answers the correlation-level form of Open Problem~\ref{open:statdim} for the sympLPN formulation: for the \emph{unconstrained} LPN ensemble the off-diagonal correlation is exactly $0$ (there $Aw$ is uniform over $\F_2^{2n}$), so the isotropic conditioning $S_A = 0$ contributes exactly $-((1+\tau)^{2n}-1)/(2^{2n}-1)$ --- negative, and exponentially smaller than the diagonal. The conditioning does not help an SQ adversary. Second, the diagonal $(1+\tau)^{2n}-1$ is exponentially \emph{large} (one sympLPN sample is a full $2n$-row block), which intrinsically limits the VSTAT strength of the resulting bound:

\begin{corollary}[sympLPN SQ bound]
\label{cor:symplpn-sq}
Let $\beta = (1+\tau)^{2n}-1$ and fix $0 \leq t \leq n-1$. Every subset of $\{D_x\}_{x \in \F_2^n}$ of size at least $2^{n-t}$ has diagonal-inclusive average correlation at most $\gamma_t = 2\beta\, 2^{t-n}$, so $\SDA(B(\{D_x\}, D_0), \gamma_t) \geq 2^t$, and any SQ algorithm distinguishing $\{D_x\}$ from $D_0$ with probability $> 2/3$ requires $\geq 2^t/3$ queries to $\VSTAT(2^{n-t}/(6\beta))$. The oracle is non-trivial ($\VSTAT$ strength $\geq 1$) for $t \leq c_p\, n - O(1)$ with $c_p = 1 - 2\log_2(1+\tau)$; at $p = 1/4$, $c_p = 1 - 2\log_2(5/4) \approx 0.356$, giving a $2^{c_p n - O(1)}$-query bound at constant VSTAT strength.
\end{corollary}

\begin{proof}
For $|T| \geq 2^{n-t}$ the off-diagonal terms contribute at most $\beta/(2^{2n}-1)$ each and the diagonal terms $\beta/|T| \leq \beta\, 2^{t-n}$, so the average is at most $\beta(2^{t-n} + (2^{2n}-1)^{-1}) \leq 2\beta\, 2^{t-n}$. Hence $\SDA \geq 2^n/2^{n-t} = 2^t$, and \Cref{thm:feldman} with $\alpha = 2/3$ gives the query bound. The strength condition $2^{n-t} \geq 6\beta$ rearranges to $t \leq n(1 - 2\log_2(1+\tau)) - O(1)$.
\end{proof}

We emphasize that \Cref{cor:symplpn-sq} bounds \emph{distinguishing} the averaged isotropic ensemble from the null, not \emph{recovery}: with a single \emph{revealed} public matrix $A$, the coordinate-mean queries $\E_y[(-1)^{y_i}] = (1-2p)(-1)^{(Ax)_i}$ recover $Ax$ in $2n$ constant-tolerance queries, and hence $x$, for every $A$ regardless of the code's minimum distance. The corollary therefore isolates the SQ-invisibility of the isotropic conditioning $S_A = 0$ in the averaged ensemble --- the point of \Cref{thm:symplpn-corr}, that $S_A=0$ does not help an SQ distinguisher relative to unconstrained LPN --- and makes no claim that fixed-matrix sympLPN is SQ-hard to solve; that problem is average-case \emph{decoding} of a random Lagrangian code, for which minimum distance gives identifiability, not SQ hardness.

The $2^n$ cap (small secret space) and the $c_p$ haircut (per-block $\chi^2$) are both intrinsic to the sympLPN formulation; the membership formulation, whose per-sample divergence is exponentially small, is where the $\Omega(2^n)$ and $2^{2n-O(1)}$ (both unconditional) SQ bounds live. The two formulations are complementary, not interchangeable.

\paragraph{The depolarizing-noise model carries the external hardness.} \Cref{def:symplpn} uses independent Bernoulli noise for analytic simplicity. The object actually carrying the external reductions of \cite{KLP+25,LPQR26} --- constant-rate LPN reduces to it (the formal \cite[Thm.\ 6.4]{KLP+25}, stated informally as their Thm.~1.6, then \cite[Thm.\ 4.1]{LPQR26}) --- is sympLPN with the depolarizing distribution's \emph{symplectic representation} as noise: the $n$ qubit-pairs $(e_{2i},e_{2i+1})$ are i.i.d.\ with $\Pr[(0,0)]=1-p$ and $\Pr[(1,0)]=\Pr[(0,1)]=\Pr[(1,1)]=p/3$. The same character-sum method gives an exact SQ bound for \emph{any} per-qubit-pair noise law --- governed by a single scalar $M_2$ --- so in particular the externally-hard depolarizing problem itself obeys an unconditional SQ lower bound.

\begin{theorem}[sympLPN SQ structure under any Pauli noise; the depolarizing instance]
\label{thm:symplpn-depol}
Let the noise be any per-qubit-pair-i.i.d.\ law $\mu$ on $\F_2^2$ (a Pauli channel $\mu^{\otimes n}$ in the symplectic representation), and set
\[
M_2 \;=\; \sum_{\sigma\in\F_2^2}\hat\mu(\sigma)^2 \qquad (\hat\mu(0)=1).
\]
Then the sympLPN correlations depend on $\mu$ \emph{only} through the scalar $M_2$:
\[
\langle D_x, D_x \rangle = M_2^{\,n} - 1, \qquad
\langle D_x, D_{x'} \rangle = -\frac{M_2^{\,n} - 1}{2^{2n}-1}\quad(x\neq x').
\]
Consequently every subset of $\{D_x\}$ of size $\ge 2^{n-t}$ has diagonal-inclusive average correlation $\le 2(M_2^n-1)\,2^{t-n}$, and (exactly as in \Cref{cor:symplpn-sq}) any SQ distinguisher requires $2^{c_p n-O(1)}$ queries at constant VSTAT strength with $c_p = 1-\log_2 M_2$. The Bernoulli model of \Cref{cor:symplpn-sq} is the case $M_2=(1+\tau)^2$, $\tau=(1-2p)^2$. The \emph{depolarizing} model carrying the external reductions --- $\Pr[(0,0)]=1-p$, $\Pr[(1,0)]=\Pr[(0,1)]=\Pr[(1,1)]=p/3$, whose three nonzero pair-Fourier coefficients are all equal to $\tau_d=1-\tfrac{4p}{3}$ --- is the case $M_2=1+3\tau_d^2$, giving $c_p^{\mathrm{dep}}=1-\log_2(1+3\tau_d^2)$, non-trivial precisely for $p>p_*:=\tfrac34\bigl(1-\tfrac1{\sqrt3}\bigr)\approx0.317$ ($\tau_d^2<\tfrac13$).
\end{theorem}

\begin{proof}
Write $\mathrm{mass}(e)=\prod_i\mu(\sigma_i(e))$ for the noise law, $\sigma_i(e)$ the $i$-th qubit-pair of $e$. Conditioning on $A$, the diagonal is $\langle D_x,D_x\rangle=2^{2n}\sum_e\mathrm{mass}(e)^2-1$; by Parseval on each $\F_2^2$ factor, $\sum_e\mathrm{mass}(e)^2=\prod_i\bigl(\tfrac14\sum_\sigma\hat\mu(\sigma)^2\bigr)=(M_2/4)^n$, so the diagonal is $4^n(M_2/4)^n-1=M_2^n-1$, independent of $A$. For the off-diagonal, $\langle D_x,D_{x'}\rangle=\sum_s\widehat{\mathrm{mass}}(s)^2(-1)^{\langle s,c\rangle}-1$ with $c=A(x{-}x')$ and $\widehat{\mathrm{mass}}(s)=\prod_i\hat\mu(\sigma_i(s))$; averaging over $c$ uniform on $\F_2^{2n}\setminus\{0\}$ factorizes per pair into $\sum_\sigma\hat\mu(\sigma)^2(-1)^{\langle\sigma,c_{\mathrm{pair}}\rangle}$, whose uniform-$c$ average is $\hat\mu(0)^2=1$, so the nonzero-$c$ average collapses to $-(M_2^n-1)/(2^{2n}-1)$ (here $c=A(x{-}x')$ is uniform on $\F_2^{2n}\setminus\{0\}$ by $\mathrm{Sp}$-transitivity of the isotropic ensemble on first columns). The SDA and \Cref{thm:feldman} steps are then identical to \Cref{cor:symplpn-sq}. Verified exactly at $n=2,3$ for the Bernoulli, depolarizing, asymmetric-independent, and a correlated pair law (and at $n=4$ for the off-diagonal).
\end{proof}

Thus the \emph{same} object that \cite{KLP+25,LPQR26} prove is at least as hard as constant-rate LPN also satisfies an exact $2^{\Theta(n)}$-query SQ lower bound, for every constant $p\in(p_*,\tfrac34)$: the external (worst-case-style) hardness and our averaged SQ-invisibility of the isotropic conditioning meet on one and the same problem. By \Cref{thm:symplpn-depol} the two noise models share the identical correlation \emph{structure} --- a diagonal $\beta=M_2^n-1$ with a flat off-diagonal $-\beta/(2^{2n}-1)$ --- and differ only in the scalar $M_2$ ($(1+\tau)^2$ for Bernoulli, $\tau=(1-2p)^2$; $1+3\tau_d^2$ for depolarizing). Structurally the $X$-coordinates of the depolarizing noise are themselves exactly i.i.d.\ $\mathrm{Bernoulli}(2p/3)$, and the depolarizing pair is the convolution of an independent $\mathrm{Bernoulli}(q)$ pair with a within-pair joint-flip $(h,h)$; a reduction \emph{equating} the two noise models on a single instance is nonetheless obstructed --- depolarizing is invariant under local Cliffords, which Bernoulli is not, and the joint-flip convolution is not invertible. The hardness inheritance below does not need such an equivalence.

\begin{proposition}[Bernoulli sympLPN inherits the constant-rate-LPN hardness]
\label{prop:bernoulli-inherits}
Constant-rate LPN reduces to sympLPN \emph{with independent Bernoulli noise} and uniform isotropic public matrix; hence \Cref{def:symplpn} as stated is at least as hard as constant-rate LPN.
\end{proposition}

\begin{proof}
We trace the reduction of \cite[Thm.~6.4]{KLP+25} from Decision LPN to Decision sympLPN and verify that its depolarizing-conversion step is not load-bearing, so the same reduction certifies Bernoulli sympLPN.

\emph{Step 1 (public matrix).} The reduction's public matrix $\mathbf K$ is within negligible total-variation distance of the uniform distribution over full-rank symplectically-orthogonal $2n\times n$ matrices: it is assembled from \cite[Lem.~6.2]{KLP+25} (a near-uniform full-rank symplectically-orthogonal extension of the LPN matrix), a \emph{noise-preserving} paired coordinate permutation $\pi\oplus\pi$ (a qubit relabelling, not a symplectic scramble), and an iterative completion of the remaining $n-\ell$ columns. This distribution is \emph{identical} for the structured (LPN-derived) and unstructured (uniform) cases --- the construction is noise-agnostic --- and matches the public-matrix law of \Cref{def:symplpn}.

\emph{Step 2 (Bernoulli noise, pre-conversion).} At this stage the noise $\mathbf e'$ is i.i.d.\ $\mathrm{Bernoulli}(q)$ on all $2n$ coordinates, the planted rate $p$ raised to $q=p+r-2pr$ by XOR-ing a fresh $\mathrm{Bernoulli}(r)$ (their noise-raising step). \cite{KLP+25} state verbatim that this sample ``is nearly distributed as a valid sympLPN structured sample, except [\dots] the error is Bernoulli instead of depolarizing'' --- i.e.\ it is exactly an instance of the \emph{Bernoulli} sympLPN of \Cref{def:symplpn}.

\emph{Step 3 (the depolarizing conversion is inert for decision).} Only afterwards does the reduction add a fresh $\mathbf h\sim\mathrm{Bernoulli}(q)^n$ to \emph{both} symplectic halves (their Cor.~2.15), converting the noise to depolarizing. The decision-correctness argument matches only the structured/unstructured output laws against the oracle's promise; it never invokes this conversion. Omitting the $\mathbf h$-step therefore leaves a valid reduction whose output is precisely the Step-2 Bernoulli sympLPN sample, and no symplectic randomization (which would correlate the Bernoulli noise) is required anywhere.

\emph{Step 4 (composition).} Each step incurs only negligible total-variation error and Step~1's public-matrix law is conversion-independent, so the composed map is a sound reduction from Decision LPN to Decision Bernoulli sympLPN. Hence \Cref{def:symplpn}, as stated, is at least as hard as constant-rate LPN.
\end{proof}

Consequently the Bernoulli sympLPN of \Cref{def:symplpn} is, simultaneously: at least as hard as constant-rate LPN (\Cref{prop:bernoulli-inherits}) --- both classically and, since the reduction is classical and polynomial-time, also \emph{quantumly} (a quantum sympLPN solver yields a quantum LPN solver) --- subject to our exact SQ lower bound (\Cref{thm:symplpn-depol} at $M_2=(1+\tau)^2$, giving \Cref{cor:symplpn-sq}), and covered by the linear-reduction barrier landscape of \Cref{sec:barriers}.

\paragraph{The self-dual code view: structural inertness and the reverse direction.} As a search problem, sympLPN is syndrome decoding of the code $L=\mathrm{Col}(A)$: the isotropy makes $A^{\top}\Omega$ a parity check ($(A^{\top}\Omega)A=A^{\top}\Omega A=0$), so $L=L^{\perp_\Omega}$ is \emph{symplectically self-dual} and the syndrome $A^{\top}\Omega y=A^{\top}\Omega e$ is the ordinary one --- the isotropy reveals nothing beyond standard syndrome decoding. Standard LPN decodes a random \emph{linear} code; sympLPN decodes a random \emph{self-dual} one. The structure appears cryptanalytically inert: the Lagrangian code's minimum distance is never smaller than a generic random rate-$\tfrac12$ code's --- in fact marginally \emph{larger} in mean (ratio $\approx 1.07$; exact enumeration at $n\le 4$, sampled at $n=5$) --- so information-set decoding sees no structural advantage; the only public signature, $A_X^{\top}A_Z$ symmetric in the $X/Z$ block form, is secret-free (hence reduction-inert, cf.\ \Cref{conj:source}); and $S_A=0$ is SQ-invisible (\Cref{thm:symplpn-depol}). Whether sympLPN \emph{reduces back} to LPN --- which, with \Cref{prop:bernoulli-inherits}, would make the two equivalent --- is exactly the question of whether syndrome decoding of random self-dual codes is equivalent to that of random linear codes, structurally the same obstruction as Ring-LWE versus plain LWE: the isotropic ensemble is a recognizable, measure-$2^{-\Theta(n^2)}$ set (the signature ``$A_X^{\top}A_Z$ symmetric'' is automatic for isotropic $A$ but holds with probability $\Theta(2^{-\binom n2})$ for a uniform matrix), so a matrix-converting reduction cannot feed it to a uniform-matrix oracle without tripping a fragile average-case oracle, and un-recognizing it requires a symplectic map that correlates the Bernoulli noise. We leave this equivalence open. Its linear case is precisely the marginal-adaptive barrier cell of \Cref{sec:barriers}: a linear reverse reduction is a map $B=g(A,R)$ outputting $(BA,\,By)$, and whether that output stays far from a \emph{usable} LPN instance is exactly the open \Cref{lem:m2} (\Cref{open:marginal-adaptive}). (Numerically the output can sit close to the LPN law at small $n$, but only below the solvability threshold $m\ge n/(1-H_2(p'))$, where the LPN target is itself near-uniform and unsolvable by \Cref{prop:diffuse-solvable} --- not a usable reduction.) The self-duality also makes every \emph{linear} syndrome check provably weak against such a reverse reduction: for any $B$ with $BA=C$ and any $w\in\ker(C^{\top})$, the effective syndrome $B^{\top}w$ lies in $\ker(A^{\top})=\Omega L$ (the standard complement equals $\Omega L$ by self-duality, with $\Omega$ a weight-preserving coordinate swap), so $\mathrm{wt}(B^{\top}w)\ge d_{\min}(L)=\Omega(n)$ with high probability over a uniform Lagrangian and the linear statistic $\langle w,y\rangle$ has bias $\le(1-2p)^{d_{\min}(L)}=2^{-\Omega(n)}$ (the only escape, $B^{\top}w=0$, holds for a $2^{-n}$ fraction of $w$ and requires knowing $B$). Any distinguisher witnessing \Cref{lem:m2} must therefore be \emph{non-linear} (e.g.\ minimum-syndrome-weight decoding), so the open obstruction is the joint low-rank structure of $Be$ rather than any linear parity test. Cryptanalytically the self-dual structure is far poorer than a ring --- no multiplication, ideals, or weight-preserving automorphisms --- and exposes no attack handle, so a random Lagrangian code is hard to tell from a random linear one and $\mathrm{sympLPN}\approx\mathrm{LPN}$ is the natural cryptanalytic expectation; but the same poverty removes the handles a reduction would use, so this expectation is not (yet) a proof, and a strict separation is not excluded.


\section{Exact Moments of the Isotropic Row Ensemble}
\label{sec:moments}

The barriers of \Cref{sec:barriers} concern the \emph{symplectic-LPN} formulation, whose public matrix $A \in \F_2^{2n \times n}$ is a uniform full-rank matrix with pairwise symplectically orthogonal columns (equivalently, a uniform ordered basis of a uniform Lagrangian; we call this the \emph{isotropic ensemble} $\mathcal{A}_n$). A natural question --- Open Problem~\ref{open:statdim} below --- is whether conditioning on the symplectic relation $S_A = 0$ changes the statistical behaviour of the rows of $A$ enough to affect SQ hardness. In this section we answer the question exactly, at the level of \emph{all} subset moments: we prove a closed form for every $m_j$ showing that the row ensemble of $\mathcal{A}_n$ deviates from the unconstrained LPN row ensemble by $\Theta(4^{-n})$, so that no statistical test of constant bundle size can exploit the symplectic conditioning.

Fix two distinct non-zero secrets $x, x' \in \F_2^n$. Over $\F_2$, any two distinct non-zero vectors are linearly independent, and $\mathrm{GL}(n,2)$ acts transitively on such ordered pairs; since the ensemble $\mathcal{A}_n$ is invariant under right multiplication by $\mathrm{GL}(n,2)$ (change of basis of the Lagrangian), all quantities below are independent of the choice of $(x,x')$, and we take $(x,x') = (e_1, e_2)$. Let $a_1, \dots, a_{2n} \in \F_2^n$ denote the rows of $A$, define the \emph{quadrant} $Q = \{\, i : \langle a_i, x\rangle = \langle a_i, x'\rangle = 1 \,\}$ and $t = |Q|$, and define the subset moments
\[
m_j \;:=\; \E_{A \sim \mathcal{A}_n}\!\left[\binom{t}{j}\Big/\binom{2n}{j}\right]
\;=\; \Pr\bigl[\text{$j$ uniformly chosen distinct rows all lie in } Q\bigr].
\]
For the \emph{unconstrained} ensemble ($A$ uniform over $\F_2^{2n\times n}$, rows i.i.d.\ uniform) one has $m_j = (1/4)^j$ exactly. The question is how far the isotropic ensemble deviates.

\paragraph{Reduction to column pairs.} With $(x,x')=(e_1,e_2)$, row $i$ lies in $Q$ iff the $i$-th bits of the first two columns $c_1, c_2$ of $A$ are both $1$. Thus $m_j$ depends only on the joint distribution of $(c_1, c_2)$, which by Witt's extension theorem (transitivity of $\mathrm{Sp}(2n,\F_2)$ on ordered pairs of distinct non-zero vectors spanning a totally isotropic $2$-space) is uniform over the set of ordered \emph{isotropic pairs}
\[
\mathcal{P} = \{(c_1,c_2) : c_1, c_2 \in \F_2^{2n}\setminus\{0\},\; c_1 \neq c_2,\; \Omega(c_1,c_2)=0\},
\qquad |\mathcal{P}| = (2^{2n}-1)(2^{2n-1}-2).
\]
By linearity of expectation over $j$-subsets $S \subseteq [2n]$,
\[
m_j = \E_{S}\,\Pr_{(c_1,c_2)}\bigl[\,c_1|_S = c_2|_S = \mathbf{1}\,\bigr],
\]
so it suffices to count, for each orbit of $j$-subsets $S$ under the symplectic coordinate symmetries, the number $q_S$ of isotropic pairs that are all-ones on $S$.

\paragraph{The counting lemmas.} Write $u := 2^{2n-2}$, let $v_0$ be the indicator vector of $S$, $W_S = \{w : w|_S = 0\}$, and $V_S = v_0 + W_S$ (so $|V_S| = 2^{2n-|S|}$ and $0 \notin V_S$). For fixed $c_1 \in V_S$, the condition $\Omega(c_1, c_2) = 0$ on $c_2 = v_0 + w_2 \in V_S$ is the affine condition $\Omega(c_1, w_2) = \Omega(c_1, v_0)$ on $w_2 \in W_S$. Two short computations in characteristic 2 show: (i) if $c_1 \in W_S^{\perp_\Omega}$ the condition is automatically satisfied by all $w_2 \in W_S$, and (ii) $w_2 = w_1$ (i.e.\ $c_2 = c_1$) always satisfies it. Since $W_S^{\perp_\Omega} = \{c : \operatorname{supp}(c) \subseteq \sigma(S)\}$, where $\sigma$ is the fixed-point-free involution pairing coordinate $i$ with $i \pm n$, membership of $c_1 \in V_S \cap W_S^{\perp_\Omega}$ requires $S \subseteq \operatorname{supp}(c_1) \subseteq \sigma(S)$.

\begin{lemma}[Symplectic-pair $2$-subsets]
\label{lem:sym2}
If $S = \{i, i+n\}$, then $q_{\mathrm{sym2}} = u(u-1)/2$.
\end{lemma}

\begin{lemma}[Generic $2$-subsets]
\label{lem:gen2}
If $S = \{i,j\}$ with $j \neq \sigma(i)$, then $q_{\mathrm{gen2}} = u(u-2)/2$.
\end{lemma}

\begin{lemma}[$3$-subsets]
\label{lem:three}
For every $3$-subset $S$ (either orbit), $q_{3} = u(u-4)/8$. In particular $q_3 = 0$ at $n=2$ ($u=4$), explaining $m_3 = 0$ there.
\end{lemma}

\begin{proof}[Proof of Lemmas~\ref{lem:sym2}--\ref{lem:three}]
In each case $q_S = \sum_{c_1 \in V_S}(|V_S \cap c_1^{\perp_\Omega}| - 1)$, the $-1$ removing $c_2 = c_1$ (always isotropic with itself), and $|V_S \cap c_1^{\perp_\Omega}|$ equals $|V_S|$ if $c_1 \in W_S^{\perp_\Omega}$ and $|V_S|/2$ otherwise, by the affine-condition classification above.

\emph{Lemma~\ref{lem:sym2}.} Here $|S|=2$, $|V_S| = u$, and $\sigma(S) = S$, so $c_1 \in V_S \cap W_S^{\perp_\Omega}$ forces $S \subseteq \operatorname{supp}(c_1) \subseteq S$, i.e.\ $c_1 = v_0 = e_i + e_{i+n}$: exactly one such $c_1$, contributing $u-1$; the remaining $u-1$ choices contribute $u/2-1$ each. Total $(u-1) + (u-1)(u/2-1) = u(u-1)/2$.

\emph{Lemma~\ref{lem:gen2}.} Now $S \cap \sigma(S) = \emptyset$ (a $2$-set is $\sigma$-invariant only if it is a symplectic pair), so $S \subseteq \operatorname{supp}(c_1) \subseteq \sigma(S)$ is impossible: no $c_1 \in V_S$ lies in $W_S^{\perp_\Omega}$, and $q = u(u/2-1) = u(u-2)/2$.

\emph{Lemma~\ref{lem:three}.} A $3$-set cannot be $\sigma$-invariant ($\sigma$ is a fixed-point-free involution, so invariant sets have even size), so again $V_S \cap W_S^{\perp_\Omega} = \emptyset$ for \emph{both} orbits; with $|V_S| = u/2$, $q = (u/2)(u/4-1) = u(u-4)/8$.
\end{proof}

\begin{theorem}[Closed forms for $m_2, m_3$]
\label{thm:mj-closed}
With $u = 2^{2n-2}$ and $P = (2^{2n}-1)(2^{2n-1}-2)$,
\begin{align*}
m_2 \;&=\; \frac{n\,q_{\mathrm{sym2}} + \bigl(\binom{2n}{2}-n\bigr)q_{\mathrm{gen2}}}{\binom{2n}{2}\,P}
\;=\; \frac{(2n-1)u^2-(4n-3)u}{4(2n-1)(4u^2-5u+1)},\\
m_3 \;&=\; \frac{q_3}{P} \;=\; \frac{u(u-4)}{16(4u^2-5u+1)}.
\end{align*}
Both 3-subset orbits contribute equally, so $m_3$ is independent of $n$ given $u$. Asymptotically
\[
\frac{1}{16} - m_2 = \frac{3}{64u} + O(u^{-2}), \qquad
\frac{1}{64} - m_3 = \frac{11}{256u} + O(u^{-2}),
\]
i.e.\ both moments converge to their unconstrained values $(1/4)^j$ at rate $\Theta(4^{-n})$.
\end{theorem}

\begin{proof}
Average $q_S$ over the orbit decomposition of the $2$- and $3$-subsets (Lemmas~\ref{lem:sym2}--\ref{lem:three}): $n$ symplectic pairs and $\binom{2n}{2}-n$ generic pairs for $j=2$; $n(2n-2)$ pair-containing and $8\binom{n}{3}$ generic triples for $j=3$; divide by $P$. The rational simplifications are routine and were machine-verified against exact enumeration for $n=2,\dots,6$ and a blind prediction-then-enumeration test at $n=7$ ($P=134{,}176{,}770$).
\end{proof}

The same orbit decomposition closes \emph{all} subset moments.

\begin{theorem}[All subset moments]
\label{thm:mj-general}
With $u=2^{2n-2}$, $D_j=2^{2n-j}$, and $P=(2^{2n}-1)(2^{2n-1}-2)$, the $j$-th subset moment is, for every $1\le j\le 2n$,
\[
m_j = \frac{\binom{2n}{j}\bigl(\tfrac12 D_j^2 - D_j\bigr) + \mathbf 1_{[j\ \mathrm{even}]}\binom{n}{j/2}\,\tfrac12 D_j}{\binom{2n}{j}\,P}.
\]
This reduces to \Cref{thm:mj-closed} at $j=2,3$, vanishes for $j\ge 2n-1$ (then $D_j\le 2$), and for each fixed $j$ satisfies $m_j = (1/4)^j + \Theta(4^{-n})$ --- approaching from above at $j=1$, where $m_1=u/(4u-1)$, and from below for every $j\ge 2$.
\end{theorem}

\begin{proof}
$\binom{t}{j}/\binom{2n}{j}$ is the probability that a uniform $j$-subset $S$ of coordinates lies in the common support of $(c_1,c_2)$, so $m_j=\binom{2n}{j}^{-1}P^{-1}\sum_{|S|=j}N(S)$ with $N(S)$ the number of ordered isotropic pairs in $V_S=\{v:v|_S=\mathbf 1\}$. Writing $c_i=v_0+w_i$ with $w_i\in W_S=\{w:w|_S=0\}$, we have $\langle c_1,c_2\rangle=L(w_1)+L(w_2)+\langle w_1,w_2\rangle$ with $L(w)=\langle v_0,w\rangle$. A character-sum evaluation, splitting $W_S$ into its radical (of dimension $b$, the number of symplectic pairs meeting $S$ in exactly one coordinate) and a non-degenerate complement, gives $N(S)=\tfrac12 D_j^2-\tfrac12 D_j$ when $S$ is a union of full symplectic pairs ($b=0$; there are $\binom{n}{j/2}$ such $S$ for even $j$ and none for odd $j$) and $N(S)=\tfrac12 D_j^2-D_j$ otherwise. Summing over the two orbit types yields the formula for every $1\le j\le 2n$; it was machine-verified against direct enumeration for all $j$ at $n=2,3,4$ and a blind test at $n=7$. The later joint generating function \Cref{thm:joint-gf} rederives the all-$j$ formula as the specialization $x_{11}\mapsto 1+x$.
\end{proof}

\paragraph{First moment.} For completeness, the first moment is elementary: the row marginal of $\mathcal{A}_n$ is the mixture $\mu_{\mathrm{row}} = \tfrac{1}{2^n+1}\,\delta_0 + \tfrac{2^n}{2^n+1}\,\mathrm{Unif}(\F_2^n \setminus \{0\})$. Indeed, row $i$ of $A$ is zero iff the Lagrangian $L = \operatorname{colspan}(A)$ lies in the coordinate hyperplane $\{x : x_i = 0\} = e_{\sigma(i)}^{\perp_\Omega}$, which by $L = L^{\perp_\Omega}$ happens iff $e_{\sigma(i)} \in L$ --- an event of probability exactly $1/(2^n+1)$ by $\mathrm{Sp}$-transitivity on non-zero vectors; uniformity over non-zero rows follows from $\mathrm{GL}(n,2)$-invariance. This gives $m_1 = 2^{2n-2}/(2^{2n}-1) = \tfrac14 + \Theta(4^{-n})$.

\begin{corollary}[Fixed-size bundles are LPN-like]
\label{cor:bundle}
Let $\sigma^2 = (1-2p)^2/(p(1-p))$ and consider the $k$-row bundle moment
$V_k = \sum_{j=0}^{k}\binom{k}{j}\sigma^{2j} m_j$, which controls the pairwise SQ correlation of $k$-row bundle examples drawn from the isotropic ensemble. For every fixed $k$, using the closed forms of \Cref{thm:mj-general},
$V_k = (1+\sigma^2/4)^k\bigl(1 + O_k(4^{-n})\bigr)$:
the isotropic ensemble is indistinguishable from the unconstrained LPN ensemble by any constant-size statistical test, up to an exponentially small advantage.
\end{corollary}

The closed forms also control the \emph{maximal} block size $k = 2n$, where the bundle consists of all rows of $A$.

\begin{proposition}[Maximal-block variance multiplier]
\label{prop:vmax}
Let $\sigma^2 = (1-2p)^2/(p(1-p))$, $X = 4^n$, and $P = (X-1)(X-4)/2$. Then
\[
V_{2n} \;=\; 1 + \frac{\tfrac{X^2}{2}\bigl[(1+\tfrac{\sigma^2}{4})^{2n}-1\bigr] \;-\; X\bigl[(1+\tfrac{\sigma^2}{2})^{2n}-1\bigr] \;+\; \tfrac{X}{2}\bigl[(1+\tfrac{\sigma^4}{4})^{n}-1\bigr]}{P}.
\]
At $p = 1/4$ (so $\sigma^2 = 4/3$) this simplifies to
\[
V_{2n} = \frac{X^4 - 2X\cdot 25^n + X\cdot 13^n - 4X\cdot 9^n + 4\cdot 9^n}{9^n\,(X-1)(X-4)},
\qquad
\frac{V_{2n}-V_{2n}^{\mathrm{iid}}}{V_{2n}^{\mathrm{iid}}} = -2\Bigl(\frac{25}{64}\Bigr)^{n} + O(4^{-n}),
\]
where $V_{2n}^{\mathrm{iid}} = (1+\sigma^2/4)^{2n} = (16/9)^n$ is the unconstrained value: the deviation is \emph{negative} and exponentially small even at the maximal block.
\end{proposition}

\begin{proof}
At $k=2n$ the bundle sum is exact rather than an average: $V_{2n} = \sum_{S \subseteq [2n]} \sigma^{2|S|}\Pr[S \subseteq Q] = \sum_{j}\binom{2n}{j}\sigma^{2j} m_j$. Substituting \Cref{thm:mj-general} and writing $D_j = X/2^j$ splits the sum into three binomial series --- $\sum_j \binom{2n}{j}(\sigma^2/4)^j$ from the $D_j^2$ part, $\sum_j \binom{2n}{j}(\sigma^2/2)^j$ from the $D_j$ part, and $\sum_i \binom{n}{i}(\sigma^4/4)^i$ from the even-$j$ orbit part --- which evaluate to the three bracketed terms. The expansion at $p=1/4$ follows by dividing by $9^n(X-1)(X-4)$ and expanding in $X^{-1}$; the closed form was additionally checked by exact-fraction summation for $n \le 14$ and four noise rates, and against direct enumeration at $n = 2,3,4$.
\end{proof}

Two remarks. First, \Cref{prop:vmax} controls one functional (the variance multiplier) at growing block size, while the individual moments near $j = 2n$ deviate by $\Theta(1)$ relative factors ($m_j = 0$ for $j \ge 2n-1$, whereas $(1/4)^j > 0$) --- their weight in $V_{2n}$ is simply exponentially small. The full \emph{law} of the quadrant count is, however, exactly recoverable (\Cref{prop:tdist} below). Second, the suppression is consistent in sign at every scale: exact enumeration shows the bundle correlation is \emph{sub-multiplicative} in $k$ --- the isotropic coupling slightly \emph{reduces} correlation relative to independent rows --- and \Cref{prop:vmax} now gives this suppression a closed form at the maximal block, $-2(25/64)^n(1+o(1))$ relative at $p=1/4$: the conditioning helps rather than harms the lower-bound side.

\begin{proposition}[Exact law of the quadrant count]
\label{prop:tdist}
For every $n \ge 1$ and $0 \le \ell \le 2n$,
\[
\Pr[t = \ell] \;=\; \sum_{j \ge \ell} (-1)^{j-\ell} \binom{j}{\ell} B_j,
\qquad B_j = \binom{2n}{j} m_j,
\]
with $m_j$ from \Cref{thm:mj-general} --- an exact rational expression for every $n$ (e.g.\ $(11/45,\,4/9,\,14/45,\,0,\,0)$ at $n=2$). Consequently \emph{every} statistic of the pair-level quadrant count --- all moments, all cumulants, $V_k$ for every $k$ --- is exactly computable. Moreover the total-variation distance to the unconstrained law $\mathrm{Bin}(2n,1/4)$ decays at the exact rate $2^{-(n+1)}$:
\[
2^n \cdot \mathrm{TV}\bigl(\Pr[t = \cdot],\, \mathrm{Bin}(2n,\tfrac14)\bigr) \;=\; \tfrac12 + O(2^{-n}).
\]
\end{proposition}

\begin{proof}
$B_j = \E[\binom{t}{j}]$, so the inversion formula is the standard binomial (inclusion--exclusion) transform; it is exact because \Cref{thm:mj-general} is. For the rate: splitting $\delta_j := B_j - \binom{2n}{j}4^{-j}$ along the three terms of \Cref{thm:mj-general} and inverting each piece gives, with $X = 4^n$ and $C = 1/((X-1)(X-4))$,
\[
\Delta_\ell = (X^2C - 1)\,\beta_\ell - 2XC\,q_\ell + (-1)^\ell XC\,r_\ell - \mathbf 1_{[\ell=0]}\tfrac{4}{X-4},
\]
where $\beta_\ell, q_\ell$ are the $\mathrm{Bin}(2n,\tfrac14), \mathrm{Bin}(2n,\tfrac12)$ laws and $r_\ell = [z^\ell]\,\bigl(\tfrac{5+2z+z^2}{4}\bigr)^n$. All $r_\ell$ are positive with $\sum_\ell r_\ell = 2^n$, so the alternating piece contributes exactly $XC \cdot 2^n = 2^{-n}(1+O(4^{-n}))$ to $\sum_\ell|\Delta_\ell|$, while the remaining pieces have total mass $\le |X^2C-1| + 2XC + \tfrac{4}{X-4} = O(4^{-n})$. Hence $\sum_\ell |\Delta_\ell| = 2^{-n} + O(4^{-n})$ and the claim follows. The laws and the decomposition were checked against direct enumeration at $n = 2, 3, 4$ and exactly for $n \le 10$ ($2^n\,\mathrm{TV} = 0.491, \dots, 0.498$): they sum to one, are non-negative, vanish for $\ell \ge 2n-1$, and re-contract to every $B_j$.
\end{proof}

The joint four-category composition closes as well, completing the pairwise level.

\begin{theorem}[Joint composition generating function]
\label{thm:joint-gf}
Let $t_{ab}$ ($a,b \in \{0,1\}$) count the coordinates of the ordered isotropic pair $(c_1, c_2)$ with $((c_1)_i, (c_2)_i) = (a,b)$. Then
\[
\E\Bigl[\prod_{ab} x_{ab}^{\,t_{ab}}\Bigr]
= \frac{1}{P}\Bigl[\tfrac12\bigl(T^{2n} + S^{\,n}\bigr) - A^{2n} - B^{2n} - C^{2n} + 2x_{00}^{2n}\Bigr],
\]
where $T = x_{11}+x_{10}+x_{01}+x_{00}$, $S = T^2 - 4(x_{10}x_{01} + x_{10}x_{11} + x_{01}x_{11})$, $A = x_{00}+x_{01}$, $B = x_{00}+x_{10}$, $C = x_{00}+x_{11}$, and $P = (2^{2n}-1)(2^{2n-1}-2)$. Every pairwise statistic of the isotropic pair is therefore exactly computable; \Cref{thm:mj-general} and \Cref{prop:tdist} are the specializations $x_{11} \mapsto 1+x$ (resp.\ $x$) with the other variables set to $1$.
\end{theorem}

\begin{proof}
Write $\mathbf 1_{\Omega(c_1,c_2)=0} = \tfrac12\sum_{\lambda \in \F_2}(-1)^{\lambda\,\Omega(c_1,c_2)}$ and sum over all $(c_1, c_2) \in V \times V$. The trivial character factorizes per coordinate and gives $T^{2n}$; the non-trivial character factorizes per symplectic coordinate pair, where the $4\times4$ sign matrix $(-1)^{a d + b c}$ contracts the per-pair sum to exactly $S$, giving $S^n$. Subtracting the excluded loci $c_1 = 0$, $c_2 = 0$, $c_1 = c_2$ (which contribute $A^{2n}, B^{2n}, C^{2n}$, all isotropic) and adding back their common intersection $\{(0,0)\}$ twice (inclusion--exclusion: the three pairwise intersections coincide) yields the formula. Verified coefficient-by-coefficient against direct enumeration at $n = 2, 3, 4$.
\end{proof}

\begin{corollary}[Disagreement count]
\label{cor:disagree}
$\Pr[t_{10} + t_{01} = k] = \binom{2n}{k}/(2^{2n}-1)$ for $1 \le k \le 2n$ (and $0$ at $k=0$): equivalently, $c_1 + c_2$ is exactly uniform over the non-zero vectors. (Directly: for each $v \neq 0$ the count $\#\{c_1 : \Omega(c_1, v) = 0,\ c_1 \notin \{0, v\}\} = 2^{2n-1}-2$ is $v$-independent.)
\end{corollary}

The first object at the multi-secret level closes by the same method. For an ordered triple $(c_1,c_2,c_3)$ that is pairwise isotropic and linearly independent (the law of three columns of $A$, after the same transitivity argument), classify coordinates by $\tau \in \F_2^3$ and let $t_\tau$ be the category counts.

\begin{theorem}[Triple composition generating function]
\label{thm:triple-gf}
For a subspace $L \subseteq \F_2^3$ put $T_L = \sum_{\tau \in L} x_\tau$,
\[
S_{\lambda, L} = \sum_{u, v \in L} (-1)^{\sum_{i<j} \lambda_{ij}(u_i v_j + u_j v_i)}\, x_u x_v \quad (\lambda \in \F_2^3),
\qquad
G_L = \tfrac18\Bigl(T_L^{2n} + \sum_{\lambda \neq 0} S_{\lambda, L}^{\,n}\Bigr).
\]
Then
\[
\E\Bigl[\prod_\tau x_\tau^{\,t_\tau}\Bigr]
= \frac{1}{P_3}\Bigl(G_{\F_2^3} - \sum_{H \,\mathrm{hyperplane}} G_H + 2\sum_{\ell\,\mathrm{line}} G_\ell - 8\, G_{\{0\}}\Bigr),
\qquad
P_3 = (2^{2n}-1)(2^{2n-1}-2)(2^{2n-2}-4).
\]
\end{theorem}

\begin{proof}
The three pairwise-isotropy constraints are imposed by the character average $\tfrac18 \sum_{\lambda}(-1)^{\sum \lambda_{ij}\Omega(c_i,c_j)}$; the trivial character factorizes per coordinate ($T_L^{2n}$) and each non-trivial one per symplectic coordinate pair ($S_{\lambda,L}^n$). A triple fails linear independence exactly when its coordinate categories lie in a proper subspace $L \subset \F_2^3$ (each linear relation $\alpha \cdot (c_1,c_2,c_3) = 0$ confines the categories to $\alpha^{\perp}$), so M\"obius inversion over the subspace lattice of $\F_2^3$ --- with $\mu = (-1)^c 2^{\binom{c}{2}}$ at corank $c$, i.e.\ $(+1, -1, +2, -8)$ --- extracts the independent triples. Verified coefficient-by-coefficient against direct enumeration at $n = 3$ ($22{,}680$ triples) and $n = 4$ ($1{,}927{,}800$); every pair-marginal reproduces \Cref{thm:joint-gf}. (The ensemble is empty at $n = 2$: $P_3 = 0$, as no isotropic $3$-space exists in $\F_2^4$.)
\end{proof}

Both \Cref{thm:joint-gf} and \Cref{thm:triple-gf} are the $k=2,3$ instances of a single construction valid for every $k$.

\begin{theorem}[$k$-tuple composition generating function]
\label{thm:kfold-gf}
The number of ordered, pairwise-isotropic, linearly independent $k$-tuples in $\F_2^{2n}$ is
\[
P_k(n) = \prod_{i=0}^{k-1}\bigl(2^{2n-i} - 2^i\bigr),
\]
which vanishes for $k > n$ (no isotropic $(n{+}1)$-tuple exists). For $1 \le k \le n$ the joint $2^k$-category composition generating function of such a $k$-tuple is
\[
\E\Bigl[\prod_{\tau \in \F_2^k} x_\tau^{\,t_\tau}\Bigr]
= \frac{1}{P_k(n)} \sum_{L \subseteq \F_2^k} \mu_{k-\dim L}\, G_L,
\qquad
\mu_c = (-1)^c 2^{\binom{c}{2}},
\]
where $G_L = 2^{-\binom{k}{2}}\bigl(T_L^{2n} + \sum_{\lambda \neq 0} S_{\lambda,L}^{\,n}\bigr)$ with $T_L, S_{\lambda,L}$ as in \Cref{thm:triple-gf} (now $\lambda \in \F_2^{\binom{k}{2}}$ indexing the $\binom{k}{2}$ pair constraints).
\end{theorem}

\begin{proof}
The count is the product of the choices for each $c_i$: avoiding $\operatorname{span}(c_1,\dots,c_{i-1})$ ($2^i$ vectors) inside its symplectic perp ($2^{2n-i}$ vectors, since the previous $c_j$ are isotropic and independent), giving $2^{2n-i} - 2^i$. The generating function is the same character average over the $\binom{k}{2}$ pairwise constraints (trivial character $\to T_L^{2n}$, each non-trivial $\to S_{\lambda,L}^n$) followed by M\"obius inversion over the subspace lattice of $\F_2^k$, exactly as at $k=3$. Verified: $P_k$ against enumeration for $k \le 3$; the construction reproduces \Cref{thm:joint-gf,thm:triple-gf} at $k=2,3$; and at $k=4$ ($n=4$, $P_4 = 46{,}267{,}200$) all pair- and triple-marginals reproduce the proven $k=2,3$ laws (a consistency test requiring no full enumeration).
\end{proof}

A clean specialization: the all-ones $k$-fold quadrant count $t_{1^k}$ deviates from its unconstrained law $\mathrm{Bin}(2n, 2^{-k})$ by exact total variation $707/5760$ at $(n,k)=(2,2)$, $1096511/27525120$ at $(3,3)$, etc.\ --- exponentially small, the same structural-closeness picture as the pairwise case.

What this does not close: the multi-secret level at $k = \Theta(n)$, and any SQ statement built on these joint laws. Reproducible enumeration code and exact-fraction data for all claims in this section are included with the code repository.

\section{Quantum Analysis}
\label{sec:quantum}

\subsection{Grover Search}
\label{subsec:grover}

The secret space has size $|\Lagr(2n,\F_2)| = 2^{n^2/2 + O(n)}$. A Grover search over this space requires $2^{n^2/4 + O(n)}$ iterations, each costing at least $\Omega(mn)$ to evaluate an LSN sample set. The total cost $2^{n^2/4 + O(n)}$ dominates all classical bounds for $n \geq 4$, so Grover is not a practical threat.

\subsection{Quantum BKW}
\label{subsec:qbkw}

A quantum version of BKW can speed up the inner Walsh--Hadamard transform and the search for low-weight parity checks, but no exponential quantum improvement over classical BKW is known. A naive chart-linearized BKW that treats the Lagrangian as a $\Theta(n^2)$-parameter secret suggests cost $2^{\Theta(n^2 / \log n)}$ \cite{BKW03}; this is attack evidence, not a proved lower bound against BKW-style adaptations, and should not be read as a formal security-parameter guarantee (the parameter table is SQ-model evidence only).

\subsection{Quantum SQ}
\label{subsec:quantum-sq}

Because LSN samples are classical bits, any quantum algorithm must either access them through classical queries or perform a quantum walk over the sample space. A Grover search over the $2^{n^2/2}$-sized secret space still requires $2^{n^2/4}$ iterations (\Cref{subsec:grover}), far above the security parameter. Thus the exponential SQ lower bound provides evidence (not proof) that no efficient quantum algorithm solves LSN.

\begin{conjecture}[Quantum LSN]
\label{conj:quantum}
In the classical-sample membership-oracle model, any quantum algorithm solving LSN$_{n,m,p}$ with constant $p$ requires $\Omega(2^n)$ samples/queries. (Coherent-oracle-access variants are a separate model, not claimed here.)
\end{conjecture}

\subsection{Dihedral Hidden Subgroup}
\label{subsec:dhsp}

The Lagrangian Grassmannian is not a group; it lacks a transitive abelian structure. Consequently, no reduction to the dihedral hidden-subgroup problem or its generalizations is known.

\subsection{What Is and Is Not Covered}
\label{subsec:quantum-covered}

We summarize the quantum attack landscape.

\paragraph{Considered quantum approaches.} The following quantum approaches have been analyzed and give no known efficient attack on LSN at constant noise:
\begin{itemize}
\item \textbf{Hidden Subgroup Problem.} The Lagrangian Grassmannian lacks a group structure with coset decomposition; no reduction to DHSP or abelian HSP is known.
\item \textbf{Grover search.} Complexity $2^{n^2/4 + O(n)}$ (over the $2^{n^2/2}$-sized secret space) exceeds all classical bounds.
\item \textbf{Quantum BKW.} Known BKW-style attacks, even with the standard quantum speed-ups in their subroutines, give no cost below the SQ-guided parameter scale; this is attack evidence, not a quantum lower bound.
\end{itemize}

\paragraph{Not covered.} We do \emph{not} claim a quantum lower bound for LSN. Quantum hardness is treated as a \textbf{conjecture} (\Cref{conj:quantum}), analogous to the status of LWE and LPN. A rigorous quantum query lower bound would require new techniques beyond the SQ model and is left as future work.

\section{Reduction Barriers and the Standard-Model Gap}
\label{sec:barriers}

\Cref{tab:barriers} summarizes the status of natural reductions.

\begin{table}[ht]
\centering
\caption{Reduction landscape for LSN.}
\label{tab:barriers}
\begin{tabular}{@{}>{\raggedright\arraybackslash}p{4.8cm}l>{\raggedright\arraybackslash}p{6.5cm}@{}}
\hline
Class (model) & Status & Mechanism \\
\hline
Real-linear queries (SQ) & BLOCKED & $\E[q]$ $L$-independent; F$_2$-linear $\leq (1-2p)2^{-n}$ (\Cref{thm:linear-sq}) \\
Poly degree $< n$ (SQ) & BLOCKED & RM distance (correlation $\leq 2^{-n}$) \\
Single-sample adaptive deg-$D$ (SQ) & BLOCKED & selection $\neq$ composition \\
Fixed $B$, any rank, $m \ge cn$ & ruled out & D.1 + Fannes: SD $\geq d-1/(mn)$ \\
Public $B$, $\rho \ge (3/2+\epsilon)n$ & ruled out & transport/Gram (constructive) \\
Adaptive $B$, conditional-uniform & ruled out & reachability (\Cref{thm:any-b-uniformity}) \\
Adaptive $B$, marginal-uniform, deterministic & ruled out & support counting: $|\operatorname{supp}(BA)|\le|\mathcal{A}_n|$ (\Cref{thm:deterministic-marginal-adaptive}) \\
Adaptive $B$, marginal-uniform, randomized & \textbf{open} & \Cref{lem:m1} proven half; \Cref{lem:m2} conditional (\Cref{open:marginal-adaptive}) \\
Non-linear / multi-sample & OPEN & no proof or construction known \\
\hline
\end{tabular}
\end{table}

Lu \etal~\cite{LPQR26} prove that linear reductions cannot reduce sympLPN to LPN (their \S2.4 and Appendix~D): for any \emph{fixed} $B \in \F_2^{m\times 2n}$ the matrix $BA$ has entropy at most $(1-d)mn$ for a constant $d$, so randomizing the public matrix forces a high-entropy $B$, which drives the output error weight past the Shannon-converse threshold. Per their Appendix~D, the proof targets the primary cryptographic regime $m=\Theta(n)$ (output rows $m=(1+\epsilon)n$); for $m=\omega(n)$ it shows error weight larger than $1/2-\delta$ for any constant $\delta$, which they note is not sufficient to fully rule out decodability, though they state they expect the barrier to extend to all $m=\operatorname{poly}(n)$. Our \Cref{thm:linear-sq} gives a complementary single-query barrier of $(1-2p)2^{-n}$ at any noise rate.

\subsection{Full-rank linear reductions are impossible}
\label{subsec:full-rank-linear}

LPQR26 leave open the case $m=\omega(n)$ (their bound shows error weight $>1/2-\delta$, which they note is insufficient to fully rule out decodability). We close the full-rank sub-case unconditionally, via a transport lemma that LPQR26 do not use.

Consider a linear reduction that maps a sympLPN instance $(A, y)$ to an LPN instance $(U, z)$ by outputting $U = BA$ for a public matrix $B \in \F_2^{m \times 2n}$. The input columns satisfy $S_A = A^{\top}\Omega A = 0$.

\begin{theorem}[Transport lemma for full-rank linear reductions]
\label{thm:transport-fullrank}
Let $B \in \F_2^{m \times 2n}$ with $\operatorname{rank}(B) = 2n$ and $m \ge 2n$. Let $B^{+} \in \F_2^{2n \times m}$ be any left inverse $(B^{+}B = I_{2n})$ and define $M := (B^{+})^{\top}\Omega B^{+} \in \F_2^{m \times m}$. Then for every isotropic-column matrix~$A$,
\[
(BA)^{\top} M (BA) \;=\; 0.
\]
A uniformly random matrix $C \in \F_2^{m \times n}$ satisfies $C^{\top} M C = 0$ with probability at most $2^{-n^2/2 + O(n)}$. Consequently, the output distribution of any fixed full-rank linear reduction is supported on a proper algebraic subset of $\F_2^{m \times n}$ and has statistical distance $1 - 2^{-\Theta(n^2)}$ from the uniform matrix distribution used in LPN.
\end{theorem}

\begin{proof}
The identity follows from
\begin{align*}
(BA)^{\top} M (BA)
&= A^{\top} B^{\top} (B^{+})^{\top} \Omega B^{+} B A \\
&= A^{\top} (B^{+}B)^{\top} \Omega (B^{+}B) A
= A^{\top}\Omega A = 0.
\end{align*}

For the probability bound, $M$ has rank $2n$ (because $B^{+}$ has rank $2n$ and $\Omega$ is non-degenerate). Choose an invertible $P \in \F_2^{m \times m}$ such that $P^{\top} M P = \operatorname{diag}(J,\dots,J,0,\dots,0)$ with $n$ copies of $J = \begin{psmallmatrix}0&1\\1&0\end{psmallmatrix}$. Writing $P^{-1}C = [C_1; C_2]$ with $C_1 \in \F_2^{2n \times n}$, the condition $C^{\top} M C = 0$ is equivalent to $C_1^{\top} \operatorname{diag}(J,\dots,J) C_1 = 0$, i.e. the column space of $C_1$ is totally isotropic in the $2n$-dimensional symplectic space $(\F_2^{2n}, \operatorname{diag}(J,\dots,J))$. The number of $n$-dimensional totally isotropic subspaces of $\F_2^{2n}$ is $|\Lagr(2n,\F_2)| = \prod_{i=1}^n (2^i+1) = 2^{n(n+1)/2+O(1)}$, while the total number of $n$-dimensional subspaces is the Gaussian binomial ${2n \brack n}_2 = 2^{n^2+O(n)}$. Hence a full-rank $C_1$ yields a totally isotropic column space with probability $2^{-n^2/2+O(n)}$; stratifying over $r=\operatorname{rank}(C_1)\le n$, the rank-$r$ contribution is at most $\Pr[\operatorname{rank}(C_1)=r]\cdot\#\{\text{totally isotropic }r\text{-subspaces}\}/{2n\brack r}_2$, whose maximum over $0\le r\le n$ is $2^{-n^2/2+O(n)}$ (attained up to $2^{O(n)}$ factors near $r=n$); summing over $r\le n$ preserves the bound, so a uniformly random $C_1$ (equivalently, a uniformly random $C$) yields a totally isotropic column space with probability $2^{-n^2/2+O(n)}$.
\end{proof}

\Cref{thm:transport-fullrank} blocks full-rank linear reductions at \emph{every} $m \ge 2n$ and \emph{every} noise rate. The near-full-rank stratum needs a distributional input: a fixed low-corank bilinear form, pulled back through a random $C$, almost never produces a low-rank Gram.

\begin{lemma}[Random-Gram corank tail]
\label{lem:random-gram-corank}
Let $M \in \F_2^{m \times m}$ be symmetric (resp.\ alternating) of rank $R$, and let $C \in \F_2^{m \times n}$ be uniform. If $R \ge n + k$ for some integer $1 \le k \le n$, then
\[
\Pr\bigl[\operatorname{rank}(C^{\top} M C) \le n-k\bigr]
\;=\; \Pr\bigl[\operatorname{corank}(C^{\top} M C) \ge k\bigr]
\;\le\; 2^{-\Omega(k^2)},
\]
where in the symmetric case $1 \le k \le n$, while in the alternating case $C^{\top}MC$ has even rank, forcing $\operatorname{corank} \ge n \bmod 2$, so the bound is asserted for $n \bmod 2 < k \le n$ (the corank values not fixed by parity).
\end{lemma}

\begin{proof}
The form $(u,v)\mapsto u^{\top}Mv$ factors through the quotient $\pi:\F_2^m \twoheadrightarrow \overline V := \F_2^m/\ker M$ (with $\dim \overline V = R$), on which $M$ induces a \emph{nondegenerate} symmetric (resp.\ alternating) form $\overline M$ with $u^{\top}Mv = \pi(u)^{\top}\overline M\,\pi(v)$. Thus $C^{\top}MC = \overline C^{\top}\overline M\,\overline C$ where $\overline C := \pi(C) \in \F_2^{R \times n}$ is uniform (a surjective linear image of a uniform matrix), and
\[
\operatorname{corank}(C^{\top}MC) \;=\; \bigl(n - \operatorname{rank}\overline C\bigr) \;+\; \dim\operatorname{rad}\!\bigl(\overline M|_{\operatorname{Col}\overline C}\bigr).
\]
For the first term, $\Pr[\operatorname{rank}\overline C \le n-j] \le {n \brack j}_2\, 2^{-jR} < 4\cdot 2^{j(n-j)}2^{-jR} = 4\cdot 2^{-j(R-n+j)} \le 4\cdot 2^{-jk}$ (a $\mathrm{GL}$ corank bound, using ${n \brack j}_2 < 4\cdot 2^{j(n-j)}$ --- the absolute constant being $\prod_{i\ge 1}(1-2^{-i})^{-1} < 4$ --- and $R-n \ge k$). For the second term we prove a uniform radical tail by a first moment. Condition on $\operatorname{rank}\overline C = n-j$; then $W := \operatorname{Col}\overline C$ is a uniform $(n-j)$-subspace of $(\overline V, \overline M)$. If $\dim(W\cap W^{\perp_{\overline M}})\ge i$ then $W\cap W^{\perp}$ contains an $i$-dimensional \emph{isotropic} subspace $U$ (necessarily with $U\subseteq W\subseteq U^{\perp}$), so with $Z_i$ the number of such $U$,
\[
\Pr[\dim(W\cap W^{\perp})\ge i] \;\le\; \E[Z_i] \;=\; I(i)\cdot\frac{{R-2i \brack (n-j)-i}_2}{{R \brack n-j}_2},
\]
where $I(i)$ is the number of $i$-dimensional isotropic subspaces of $(\overline V, \overline M)$. In the alternating case $\overline V$ is symplectic of rank $R=2r$ and $I(i)={r\brack i}_2\prod_{\ell=r-i+1}^{r}(2^{\ell}+1)$, while the Gaussian-binomial ratio is $2^{-i(R-i)+O(i)}$; multiplying gives $\E[Z_i]=2^{-i(i-1)/2+O(i)}=2^{-\Theta(i^2)}$, uniformly in $j$; the parity $\dim\operatorname{rad}\equiv n\pmod 2$ holds because $\overline M|_W$ is then alternating of even rank. The symmetric case is analogous: the number of totally isotropic $i$-subspaces of a nondegenerate symmetric form over $\F_2$ is at most $2^{\,i(R-i)-\Theta(i^2)}$ --- the orthogonal (quadratic) polar-space order (standard finite classical-polar-space enumeration, e.g.\ \cite{BCN89}; the type $O^{\pm}$ shifts only lower-order terms), which a non-alternating symmetric form only undercuts --- so the first moment is again $\E[Z_i]=2^{-\Theta(i^2)}$ (an upper bound, with no parity restriction). This first moment is exact at the maximal radical dimension: it equals $3/7$, $3/31$, $45/3937$ for $(n,R,i)=(2,4,2),(3,6,3),(4,8,4)$, matching direct enumeration. Splitting $\{\operatorname{corank}\ge k\}$ by $j = n-\operatorname{rank}\overline C$ and applying this with $i=k-j$ gives $\sum_{j=0}^{k} 4\cdot 2^{-jk}\,2^{-\Theta((k-j)^2)} \le 4(k+1)\,2^{-\Omega(k^2)} = 2^{-\Omega(k^2)}$, since $jk + \Theta((k-j)^2) \ge \Omega(k^2)$ uniformly on $0 \le j \le k$.
\end{proof}

The next theorem closes the near-full-rank stratum.

\begin{theorem}[Transport lemma for near-full-rank linear reductions]
\label{thm:transport-nearfull}
Let $B \in \F_2^{m \times 2n}$ with $\operatorname{rank}(B) = \rho = 2n-c$ where $1 \le c \le n-1$. Let $K = \ker B$ ($\dim K = c$). Then there exists a matrix $M \in \F_2^{m \times m}$ such that for every isotropic-column matrix $A$,
\[
\operatorname{rank}\bigl((BA)^{\top} M (BA)\bigr) \;\le\; 2c - \operatorname{rank}(\Omega|_K) \;\le\; 2c.
\]
In particular, an adversarial reduction can choose $K$ isotropic (possible for all $c \le n$), in which case $M$ is alternating and the Gram rank of the output is at most $2c$ deterministically. Moreover, in the regime $n - 2c \ge 2$ (equivalently $\rho \ge \tfrac{3}{2}n + 1$) a uniformly random matrix $C \in \F_2^{m \times n}$ satisfies $\operatorname{rank}(C^{\top} M C) \le 2c$ only with probability $2^{-\Omega((n-2c)^2)}$, so the rank-$(\le\!2c)$ test distinguishes the reduction's output from uniform. When $n - 2c \le 1$ the test is vacuous --- for $2c \ge n$ because every $n \times n$ Gram has rank $\le n \le 2c$, and for $2c = n-1$ with $n$ odd because the \emph{alternating} output Gram has even rank $\le n-1 = 2c$ --- and that regime is covered instead by the entropy-deficiency branch of the rank stratification below.
\end{theorem}

\begin{proof}
Rank-factorize $B = \tilde{B} P$ with $\tilde{B} \in \F_2^{m \times \rho}$ full column rank and $P \in \F_2^{\rho \times 2n}$ full row rank. Then $(BA)^{\top} M (BA) = (PA)^{\top} (\tilde{B}^{\top} M \tilde{B}) (PA)$, and $Q := \tilde{B}^{\top} M \tilde{B}$ ranges over all of $\F_2^{\rho \times \rho}$ (lift: $M := (\tilde{B}^{+})^{\top} Q \tilde{B}^{+}$ for a left inverse $\tilde{B}^{+}$). The set of matrices $N$ of the form $P^{\top}QP$ is therefore exactly the set of bilinear forms vanishing on $K$ from both sides. Fix a basis adapted to $K \oplus V$ with $\dim V = 2n-c$. In this basis the standard symplectic form $\Omega$ becomes a block matrix with blocks $\Omega_{KK}, \Omega_{KV}, \Omega_{VK}, \Omega_{VV}$. Adding a form $N$ that kills $K$ on both sides changes only the $VV$ block; the blocks $\Omega_{KK}, \Omega_{KV}, \Omega_{VK}$ remain fixed. We therefore obtain a \emph{minimal-rank completion} problem for the matrix
\[
E = \begin{pmatrix} \Omega_{KK} & \Omega_{KV} \\ \Omega_{VK} & E_{VV} \end{pmatrix},
\]
where $E_{VV}$ is free. The classical formula for the minimal rank of such a completion (classical minimal-rank-completion folklore) gives
\[
\min_{E_{VV}} \operatorname{rank}(E)
= \operatorname{rank}\begin{pmatrix} \Omega_{KK} & \Omega_{KV} \end{pmatrix}
+ \operatorname{rank}\begin{pmatrix} \Omega_{KK} \\ \Omega_{VK} \end{pmatrix}
- \operatorname{rank}(\Omega_{KK}).
\]
Since $\Omega$ is invertible, any $c$ rows (resp.\ columns) of $\Omega$ are linearly independent, so the two row- and column-ranks are exactly $c$. Hence
\[
\min_{E_{VV}} \operatorname{rank}(E) = 2c - \operatorname{rank}(\Omega|_K).
\]
The minimum is achieved constructively by $E_{VV} = \Omega_{VK} \cdot G \cdot \Omega_{KV}$ with $G$ any generalized inverse of $\Omega_{KK}$ (over $\F_2$, obtained from a rank factorization). A symmetric g-inverse always exists: $G^* := G^{\top} \Omega_{KK} G$ satisfies $\Omega_{KK} G^* \Omega_{KK} = \Omega_{KK}$ and $(G^*)^{\top} = G^*$, yielding a symmetric $E_{VV}$ and hence a symmetric matrix $M$ representing the optimal form. In the adversarially-relevant case of \emph{isotropic} $K$ ($\Omega|_K = 0$, so $\Omega_{KK}=0$ and the optimal completion takes $E_{VV}=0$), the form $N = B^{\top}MB$ representing $M$ is $\Omega$ with its $K$-indexed rows and columns zeroed (the zero-padded restriction $\Omega|_V$). Writing $E := \Omega + N$ for the structured-output form --- of minimal rank $2c$, which is what bounds the reduction's own Gram $A^{\top}NA = A^{\top}EA$ (using $A^{\top}\Omega A = 0$) --- subadditivity on $\Omega = N + E$ gives $\operatorname{rank} N \ge \operatorname{rank}\Omega - \operatorname{rank} E = 2n-2c$ (and the lift preserves rank: $\operatorname{rank} M = \operatorname{rank} Q = \operatorname{rank}(P^{\top}QP) = \operatorname{rank} N$, since $\tilde B^{+}$ and $P$ have full row rank $\rho$). Thus the small-rank object is $E$, not $M$: $M$ itself is \emph{alternating} of rank $\ge 2n-2c \ge n+\Omega(n)$ (for $c \le (\tfrac12-\epsilon)n$), and the output Gram $C^{\top}MC$ is alternating too (zero diagonal over $\F_2$, rather than a general symmetric matrix).

For the probability bound, restrict to the regime $n - 2c \ge 2$ (the rank-$(\le\!2c)$ test is vacuous otherwise, as noted in the statement). The output Gram $C^{\top}MC$ is symmetric (general $K$) or alternating (isotropic $K$), and by the previous paragraph $M$ has rank $R \ge 2n-2c = n + (n-2c)$. Applying \Cref{lem:random-gram-corank} with $k = n-2c \ge 2$ (so the hypothesis $R \ge n+k$ holds, and since $n-2c \equiv n \!\!\pmod 2$ we have $k > n \bmod 2$, clearing the alternating parity floor) yields
\[
\Pr\bigl[\operatorname{rank}(C^{\top}MC) \le 2c\bigr] \;=\; \Pr\bigl[\operatorname{corank}(C^{\top}MC) \ge n-2c\bigr] \;\le\; 2^{-\Omega((n-2c)^2)}.
\]
The complementary distinguishing advantage is $1 - 2^{-\Omega(n^2)}$ precisely when $n - 2c = \Omega(n)$, i.e. $\rho \ge (3/2+\epsilon)n$ for a constant $\epsilon>0$ --- the near-full stratum recorded below.
\end{proof}

\paragraph{Rank stratification.} \Cref{thm:transport-fullrank} and \Cref{thm:transport-nearfull}, together with the LPQR26 entropy-deficiency argument for the low-rank branch (in the stated linear-sample regimes), stratify the fixed/public-$B$ linear-reduction landscape by rank:
\begin{itemize}
\item $\rho = 2n$ (full rank): \textbf{BLOCKED} --- output Gram is identically zero (\Cref{thm:transport-fullrank}).
\item $(3/2+\epsilon)n \le \rho < 2n$ (near-full): \textbf{BLOCKED} --- output Gram rank $\le 2c \le (1-\epsilon)n$ (\Cref{thm:transport-nearfull}), while the uniform Gram has rank $\approx n$ w.h.p.
\item $\rho \le (3/2-\epsilon)n$ (mid/low rank): \textbf{ruled out} for fixed $B$ at $m \ge cn$ --- the rank test dies, but the output matrix $BA$ is statistically far from uniform. By LPQR26's Theorem~D.1, $H(BA) \le (1-d)mn$ for a constant $d > 0$; the Fannes--Csisz\'{a}r continuity inequality gives $\operatorname{SD}(BA, \operatorname{Uniform}) \ge d - 1/(mn)$. At $n=41$, $m=82$ ($m=2n$), the frame entropy $H(A) \le (3/2)n^2 + n/2$ gives $d \ge 1/4 - 1/(4n) \approx 0.244$, so this distance is at least $0.24$. Since any LPN solver expects a uniform matrix, the reduction fails regardless of rank. The rank-stratification theorems above retain value as \emph{constructive} distinguishers (explicitly exhibiting the defect), whereas the Fannes argument is non-constructive but covers all ranks uniformly.
\end{itemize}
This stratification applies to \emph{fixed/public} $B$ (per-realization). For $m < cn$ the Fannes bound is not triggered; this corner case is cryptographically irrelevant because any reduction targeting a \emph{solvable} instance must supply $m \gtrsim n/(1-H_2(p'))$ samples (below this syndrome-decoding information threshold the instance cannot be solved at all). If the reduction draws $B$ from a high-entropy distribution conditioned on $BA$ being uniform, the matrix is randomized---but then the label noise $Be$ is killed by piling-up at $p=1/4$ (row weight $w$ gives bias $2^{-w}$). The remaining regime is the secret-$B$ model of \Cref{lem:affine-coset-bias}, where the quantitative trade-off is governed by the Gram-rank transition.

\paragraph{B-visibility and the secret-$B$ regime.} The quantifier order of LPQR26's Theorem~D.2 is subtle: $B$ may be chosen as a function of $A$ provided only that the \emph{induced distribution} of $BA$ is uniform. This creates a \textbf{secret-$B$} regime in which the reduction matrix is not public. Our transport theorems (which require the realized $B$ to be available to the distinguisher) do not apply here.  The residual open question was whether label-signal bounds together with conditional $BA$-uniformity suffice to close the linear-reduction landscape in the conditional model at constant noise.  \Cref{thm:any-b-uniformity} resolves the \textbf{conditional-uniformity} case: when $BA$ is uniform conditioned on each fixed $A$, the reachability bound forces high-weight rows regardless of $B$'s visibility.  The \textbf{marginal-uniformity} case, where $B=g(A)$ yields uniform $BA$ marginally over $A$ but not conditionally, is not resolved: low-weight rows that vary with $A$ can evade the per-instance bound because $\bigcup_A R_w(A)$ may cover $\F_2^n$.  This marginal-adaptive corner remains open (Open Problem~\ref{open:marginal-adaptive}). In the \textbf{public-$B$} regime, by contrast, randomizing $B$ does not escape our rank-stratification barrier because the detector operates on the realized matrix. The secret-$B$ regime is loosely analogous to the gap between our membership-LSN and the stabilizer-decoding LSN of \cite{KLP+25,LPQR26}: in both settings the adversary does not fully observe the structure that generates the samples.

The label-signal side of the secret-$B$ regime is controlled by the following affine-coset bias bound.

\begin{lemma}[Affine-coset bias, expectation form]
\label{lem:affine-coset-bias}
Let $A \in \F_2^{D \times n}$ have rank $n$ and isotropic columns, let $c \in \F_2^n$, and let $b$ be uniform over the affine subspace $\{b \in \F_2^D : b^{\!\top} A = c^{\!\top}\}$. For $e \sim \operatorname{Bernoulli}(p)^{\otimes D}$,
\[
  \bigl|\E_{b,e}[(-1)^{b^{\!\top} e}]\bigr| \;\le\; 2^{-n} \cdot W_N(1-2p),
\]
where $W_N$ is the weight enumerator of $N = \operatorname{nullspace}(A^{\!\top})$ and $D = 2n$. In expectation over a uniformly random isotropic $A$,
\[
  \E_A\bigl|\E_{b,e}[(-1)^{b^{\!\top} e}]\bigr|
  \;\le\; 2^{-n} + (1-p)^{2n}
  \;=\; 2^{-n} + \Bigl(\frac{9}{16}\Bigr)^n \quad\text{at }p=1/4.
\]
When the coset offset $b_0=0$ every term in the MacWilliams sum is positive, so the bound is an equality.
\end{lemma}

\begin{proof}
Write $b = b_0 + Nz$ with $z \sim \operatorname{Uniform}(\F_2^n)$ and $N$ a basis of $\operatorname{nullspace}(A^{\!\top})$. Then $\E_b[(-1)^{b^{\!\top} e}] = (-1)^{b_0^{\!\top} e} \cdot \mathbf{1}_{e \in \operatorname{colspace}(A)}$, so $|\E_{b,e}[(-1)^{b^{\!\top} e}]| = |\sum_{x \in \operatorname{colspace}(A)} (-1)^{b_0^{\!\top} x}\, p^{|x|}(1-p)^{D-|x|}| \le \sum_{x \in \operatorname{colspace}(A)} p^{|x|}(1-p)^{D-|x|} = (1-p)^D \cdot W_{\operatorname{colspace}(A)}(p/(1-p))$, with equality when $b_0 = 0$ (all terms positive). The MacWilliams identity with $t = p/(1-p)$ gives $(1-p)^D W_{\operatorname{colspace}(A)}(t) = 2^{-n} W_N(1-2p)$. For the expectation, $N = \Omega \cdot L$ is the image of a uniform Lagrangian $L$, so $\Pr[x \in N] = 1/(2^n+1)$ for every non-zero $x$ by $\mathrm{Sp}(2n)$-transitivity. Hence $\E_A[W_N(1/2)] = 1 + ((3/2)^D-1)/(2^n+1) \le 1+(9/8)^n$, yielding the claimed bound.
\end{proof}

The expectation bound can be upgraded to a high-probability statement.  Write $W := W_N(1/2)-1 = \sum_{v\neq 0} \mathbf{1}_{v\in N}\,2^{-|v|}$.  The exact first two moments are
\[
  \E[W] = \frac{(3/2)^{2n}-1}{2^n+1}, \qquad
  \operatorname{Var}[W] = p(1{-}p)\,D + q\,S_0 - p^2 T,
\]
with $p=\tfrac{1}{2^n+1}$, $q=\tfrac{1}{(2^{n-1}+1)(2^n+1)}$, $D=(5/4)^{2n}-1$, $T=((3/2)^{2n}-1)^2-D$, and $S_0=\tfrac12(T+C)$.  The character sum
\[
  C = \!\!\sum_{\substack{v\neq v'\\ v,v'\neq 0}}\!\! (-1)^{\Omega(v,v')}2^{-|v|-|v'|}
    = (7/4)^{2n} - 2(3/2)^{2n} - (5/4)^{2n} + 2
\]
factorises over the $n$ symplectic coordinate blocks, each block contributing
\[
\sum_{a,b,x,y\in\{0,1\}}(-1)^{ay+bx}\,2^{-(a+b+x+y)} = (7/4)^2
\]
(see \Cref{app:krawtchouk} for the full derivation).  Consequently $\operatorname{Var}[W]/\E[W]^2 = O\!\bigl((50/81)^n\bigr)$, and Chebyshev's inequality gives the following.

\begin{lemma}[Affine-coset bias, w.h.p.]
\label{lem:affine-coset-bias-whp}
For a uniformly random full-rank isotropic $A \in \F_2^{2n \times n}$ and uniformly random noise $e$,
\[
  \bigl|\E_{e}\bigl[(-1)^{b^{\!\top} e}\bigr]\bigr|
  \;\le\; \bigl(2^{-n} + (9/16)^n\bigr)\bigl(1+o(1)\bigr)
  \qquad\text{with probability } 1 - 2^{-\Omega(n)}.
\]
\end{lemma}

Thus, except with exponentially small probability, the per-row label bias of the secret-$B$ regime is bounded by $(9/16)^n(1+o(1))$; the secret-$B$ trade-off is governed by the Gram-rank transition, not by bias amplification.

\paragraph{Low-weight reachability barrier (any-$B$).}
The preceding affine-coset bound controls the per-row label signal when the reduction is forced into a high-rank affine subspace.  A complementary barrier applies to \emph{any} $B$ model (adaptive or fixed, public or secret) under the \textbf{conditional-uniformity} hypothesis ($\operatorname{SD}(BA\mid A)\le\delta$ for each fixed $A$).  It closes the linear-reduction landscape in this conditional model at constant noise $p=1/4$ for all $m=\operatorname{poly}(n)$.  The counting argument is independent of $A$, but the hypothesis is per-instance: the marginal-adaptive corner, where $B=g(A)$ yields marginally-uniform $BA$ via low-weight rows that vary with $A$, remains open (see \Cref{open:marginal-adaptive}).

Let $A\in\F_2^{2n\times n}$ be any isotropic basis and let $B\in\F_2^{m\times 2n}$ be any matrix (possibly depending on $A$).  For $w\ge 0$ define the \emph{reachability set}
\[
R_w(A) \;:=\; \{\,b^{\!\top} A \;:\; b\in\F_2^{2n},\; |b|\le w \,\} \subseteq \F_2^{n}.
\]
Every row of $BA$ that is produced with a row of $B$ of weight at most $w$ must lie in $R_w(A)$.  The cardinality of $R_w(A)$ is bounded independently of $A$:

\begin{lemma}[Reachability bound]
\label{lem:reachability}
For every $A\in\F_2^{2n\times n}$ and every $w\ge 0$,
\[
|R_w(A)| \;\le\; \sum_{j=0}^{w} \binom{2n}{j}.
\]
\end{lemma}

\begin{proof}
Each $b\in\F_2^{2n}$ with $|b|\le w$ contributes at most one element $b^{\!\top}A$ to $R_w(A)$.  The number of such $b$ is exactly $\sum_{j=0}^{w}\binom{2n}{j}$.
\end{proof}

For $w=\lfloor 0.19n\rfloor$ and large $n$, the binomial sum is at most $2^{2n\cdot H(0.095)}\le 2^{0.906n}$ (since $2H(0.095)=0.9059\ldots$), where $H(\cdot)$ is the binary entropy function.  Hence
\begin{equation}
\label{eq:reach-mass}
\frac{|R_w(A)|}{2^{n}} \;\le\; 2^{-0.094n}.
\end{equation}
This ratio is exponentially small, regardless of $A$.

\begin{theorem}[Any-$B$ uniformity forcing]
\label{thm:any-b-uniformity}
Let $A\in\F_2^{2n\times n}$ be any isotropic basis and let $B\in\F_2^{m\times 2n}$ be any matrix (possibly depending on~$A$).  Let $w=\lfloor 0.19n\rfloor$ and suppose that the conditional distribution satisfies $\operatorname{SD}(BA\mid A,\,\operatorname{Uniform}_{\F_2^{m\times n}})\le\delta$ for every $A$ (so that, by data processing, each row marginal $b_i^{\top}A$ is $\delta$-close to $\operatorname{Uniform}_{\F_2^{n}}$).  Then for every row $b_i$ of~$B$,
\[
\Pr\bigl[\,|b_i|\le w \mid A\,\bigr] \;\le\; \frac{|R_w(A)|}{2^{n}} + \delta \;\le\; 2^{-0.094n}+\delta.
\]
In particular, if $m\le\operatorname{poly}(n)$ and $\delta=o(1/m)$, then with probability $1-o(1)$ over the randomness of the reduction \emph{every} row of $B$ has weight $>w$.

The bound is \emph{monotone in the uniformity gap}: if $\operatorname{SD}(BA\mid A)\le\delta$ then the same conclusion holds with $\delta$ replaced by the actual gap.  The theorem does \emph{not} apply to the marginal model where only $\operatorname{SD}(BA,\text{Uniform})\le\delta$ is required (averaged over $A$): low-weight rows that vary with $A$ can yield marginally-uniform output while evading the per-instance reachability bound.
\end{theorem}

\begin{proof}
If $|b_i|\le w$ then $c_i=b_i^{\!\top}A\in R_w(A)$.  Under a uniform distribution on $\F_2^{n}$ the probability of $R_w(A)$ is $|R_w(A)|/2^{n}\le 2^{-0.094n}$.  Since $\operatorname{SD}(BA\mid A{=}a,\,\operatorname{Uniform}_{\F_2^{m\times n}})\le\delta$ for the fixed $A=a$ (so the row marginal $b_i^{\top}A$ is $\delta$-close to $\operatorname{Uniform}_{\F_2^{n}}$), the actual probability is at most $2^{-0.094n}+\delta$.  A union bound over $m$ rows gives
\[
\Pr[\exists i:\,|b_i|\le w]\;\le\; m\bigl(2^{-0.094n}+\delta\bigr)\;=\;o(1)
\]
when $m=\operatorname{poly}(n)$ and $\delta\le 1/(4mn)$.
\end{proof}

\Cref{thm:any-b-uniformity} applies to adaptive $B$, secret $B$, and any distribution over $B$, \emph{provided the conditional-uniformity hypothesis holds}: $BA$ must be close to uniform conditioned on each fixed $A$. In that model, once $BA$ is forced to be uniform (otherwise the output is distinguishable from standard LPN), every row of $B$ must have weight $>0.19n$. The theorem does \emph{not} cover the marginal-adaptive corner (\Cref{open:marginal-adaptive}).

The row-weight lower bound immediately amplifies the noise.

\begin{corollary}[Noise amplification]
\label{cor:noise-amp}
Under the hypotheses of \Cref{thm:any-b-uniformity}, every output label $a_i=\langle b_i,e\rangle+\langle c_i,x\rangle$ has effective noise parameter
\[
p' \;\ge\; \frac{1-(1-2p)^{w+1}}{2} \;=\; \frac{1-2^{-(w+1)}}{2}
\]
at $p=1/4$.  For $w=\lfloor 0.19n\rfloor$ this gives $p'\ge \tfrac12-2^{-0.19n-1}$.
\end{corollary}

\begin{proof}
The label noise is the XOR of $|b_i|$ independent $\operatorname{Bernoulli}(p)$ variables.  The bias of XOR of $w'$ i.i.d. $\operatorname{Bernoulli}(p)$ bits is $(1-2p)^{w'}$.  Since $|b_i|>w$, the bias is at most $(1-2p)^{w+1}=2^{-(w+1)}$, which translates to the stated $p'$.
\end{proof}

Finally, the amplified noise forces an information-theoretic sample lower bound.

\begin{corollary}[Recovery barrier]
\label{cor:recovery-barrier}
No algorithm can recover $x\in\F_2^{n}$ from $m$ samples carrying \emph{independent} per-sample noise of rate $p'\ge \tfrac12-2^{-0.19n-1}$ with probability $>1/3$ unless
\[
m \;=\; \Omega\bigl(n\,\cdot\,2^{0.38n}\bigr).
\]
In particular, any linear reduction from symplectic LPN to standard LPN that produces uniform queries requires super-polynomially many output samples.
\end{corollary}

\begin{proof}
Each sample carries at most $1-h_2(p')$ bits of information about $x$, where $h_2$ is the binary entropy function; since the per-sample noises are independent, these informations add, giving $I(x;\text{transcript})\le m\,(1-h_2(p'))$.  For $p'=\tfrac12-\epsilon$ with $\epsilon=2^{-0.19n-1}$, we have $1-h_2(p')=\Theta(\epsilon^2)=\Theta(2^{-0.38n})$.  To resolve $n$ bits of entropy requires $m\ge n/\Theta(2^{-0.38n})=\Omega(n\cdot 2^{0.38n})$.
\end{proof}

Together, \Cref{lem:reachability}, \Cref{thm:any-b-uniformity}, \Cref{cor:noise-amp}, and \Cref{cor:recovery-barrier} close the linear-reduction landscape \emph{in the conditional-uniformity model}: any linear reduction that makes $BA$ uniform \emph{conditioned on each $A$} is forced into high-weight rows, which amplifies the noise to $p'\ge \tfrac12-2^{-\Theta(n)}$, and recovery then requires $m=2^{\Theta(n)}$.  At $m=\operatorname{poly}(n)$ the reduction fails in this model, for every $B$ model, every isotropic $A$, and every noise rate $p=\Theta(1)$.  (\Cref{cor:recovery-barrier} treats the output noise as independent; the genuinely correlated noise $Be$, confined to a space of dimension $\le 2n$, is confronted directly in the marginal-adaptive analysis below --- \Cref{lem:confined-g} and \Cref{prop:diffuse-solvable} --- rather than through this information bound.)  The marginal-adaptive corner, where $B$ may depend on $A$ in a way that yields uniform $BA$ marginally but not conditionally, is not covered by this argument.

A sharper, unconditional barrier comes from an information-theoretic floor:

\begin{proposition}[Entropy floor]
\label{prop:entropy-floor}
Any reduction from $\mathrm{LSN}(n)$ to $\mathrm{LPN}(k)$ that preserves the secret with recovery probability $\alpha(n)$ must satisfy
\[
k \;\geq\; \alpha(n)\,H(L) - O(1) \;=\; \alpha(n)\,\Theta(n^2).
\]
In particular, constant success probability requires $k=\Omega(n^2)$, and success probability $1-o(1)$ requires $k \geq H(L)-o(n^2)$.
\end{proposition}

\begin{proof}
The LSN secret space is $\Lagr(2n,\F_2)$ with $|\Lagr| = \prod_{i=1}^n (2^i+1)$. Hence $H(L) = \log_2|\Lagr| = n(n+1)/2 + \Theta(1)$, where the additive constant is $\log_2\prod_{i=1}^{\infty}(1+2^{-i}) \approx 1.25$. If a reduction outputs an $\mathrm{LPN}(k)$ instance from which $L$ is recovered with probability $\alpha$, Fano's inequality gives $k \geq \alpha H(L) - H_2(1-\alpha) = \alpha H(L) - O(1)$.
\end{proof}

\paragraph{Attack-cost consequence of the entropy floor.} For near-perfect recovery of $L$ (i.e., $\alpha=1-o(1)$ in Proposition~\ref{prop:entropy-floor}), the entropy floor forces $k \geq H(L)-o(n^2)$. For the parameter sets in \Cref{tab:parameters} this gives $k \gtrsim 862$ ($\lambda{=}80$), $2146$ ($\lambda{=}128$), $4754$ ($\lambda{=}192$), or $8386$ ($\lambda{=}256$), matching the $|\Lagr|$ column. The best known attack on $\mathrm{LPN}(k)$ (BKW \cite{BKW03}) runs in time $2^{\Theta(k/\log k)}$. We caution that this is an \emph{asymptotic} rate: the exponent $k/\log_2 k$ hides the constant in $\Theta(\cdot)$, polynomial factors, and the noise-rate dependence, so no concrete bit-security figure can be read off by substituting these dimensions --- it indicates only that the $k=\Theta(n^2)$ dimensions forced by the entropy floor place BKW asymptotically well beyond the corresponding $80$/$128$/$192$/$256$-bit budgets. For context, BKW on $\mathrm{LPN}(\Theta(n^2))$ runs in $2^{\Theta(n^2/\log n)}$, asymptotically smaller than brute-force on LSN ($2^{\Theta(n^2)}$) but still far above the parameter ranges of interest. No concrete-security claim is made (cf.\ the claims posture).

\paragraph{Single-sample adaptivity does not increase degree.} A natural question is whether \emph{adaptive} degree-$D$ feature-map reductions circumvent the static degree barrier. For the restricted model in which the reduction is \emph{single-sample}---it may choose each $\phi_j$ adaptively, but applies it only to the raw samples $(a,b)$, with the choice of $\phi_j$ made as a function of previous answers---the answer is negative: in each round the feature map $\phi_j(a,b)$ has degree at most $D$, and the effective degree of the entire transcript is bounded by $D$. When $D < n$, the Reed--Muller obstruction of \Cref{sec:decoders} still applies: correlation with $\mathbf{1}_L$ is at most $2^{-n}$, requiring exponentially many samples. Thus single-sample adaptive degree-$D$ feature-map reductions are blocked for every $D < n$, independent of the number of rounds.

The remaining open question is whether a \emph{multi-sample} or \emph{genuinely non-linear} adaptive reduction exists. Such a reduction could apply non-linear functions to \emph{fresh} samples generated from previous outputs, or could use adaptively chosen non-linear queries in a way that raises the algebraic degree of the transcript. We do not claim impossibility for this richer model. We note a complementary point: a reduction in this richer model would likely improve random self-reductions for LPN itself, and so would be of independent interest even if it placed LSN within an existing assumption.

\paragraph{The symplectic structure is reduction-inert.}
A sharper observation comes from the structure of the LSN samples themselves. The deterministic quadratic constraint $S_A = 0$ (the $\Omega$-Gram of the public matrix columns) is \emph{public} and \emph{$x$-free}: it depends only on the public matrix $A$, not on the secret $x$ or the noise. Consequently, $S_A=0$ adds no $x$-information to an attacker and provides no usable handle for a reductionist trying to embed LPN into LSN. This mirrors the accepted status of Ring-LWE, which is regarded as a lattice-family assumption by \emph{source} (ring structure native to the problem) rather than by reduction to plain LWE. We formalize this as:

\begin{conjecture}[Source novelty]
\label{conj:source}
The following three invariants of LSN have no analogue in standard LPN, and each obstructs every natural reduction class that has been proposed --- by being destroyed under the reduction (items 1 and 3) or, for the $x$-free constraint (item 2), by carrying no reduction-usable information:
\begin{enumerate}
\item \textbf{Symplectic-Fourier self-duality.} The indicator $\mathbf{1}_L$ satisfies $\mathcal{F}_\Omega[\mathbf{1}_L] = 2^n \mathbf{1}_L$, where $\mathcal{F}_\Omega$ is the symplectic Fourier transform. A standard linear function $f(x)=\langle s,x\rangle$ has no such fixed-point property; its Fourier transform is a Dirac mass, not a function.
\item \textbf{$x$-free quadratic constraint (sympLPN).} In the sympLPN variant (\Cref{def:symplpn}) the public matrix $A$ satisfies $S_A=0$ (the $\Omega$-Gram of columns is identically zero). This is a deterministic, public, quadratic constraint on the \emph{public parameters} alone. The batch-LSN distribution with pseudorandom columns (\Cref{def:lsn-batch}) draws pseudorandom rather than isotropic columns, so this invariant does not apply there; it is a structural distinction of the sympLPN variant studied by LPQR26. LPN has no analogous public structure linking its samples. Being $x$-free, this constraint is moreover \emph{reduction-inert}: public and independent of the secret, it offers a reductionist no handle (cf.\ the discussion preceding this conjecture), so it resists reductions by inertness rather than by being destroyed.
\item \textbf{Stabilizer degeneracy.} Every LSN secret $L$ is a stabilizer state, and a \emph{uniformly random} secret has stabilizer (minimum Pauli) weight $\Omega(n)$ with high probability (a vanishing fraction, such as the standard Lagrangian $\F_2^n\times 0$ with its weight-$1$ $X$-generators, are low-weight exceptions). The secret space is the Lagrangian Grassmannian, with entropy $H(L)=\Theta(n^2)$, not an $O(n)$-bit vector space; LPN secrets are bare linear functions.
\end{enumerate}
Each invariant is preserved by the symplectic group $\mathrm{Sp}(2n)$ but is absent from any standard LPN sample, and each proposed reduction class is barred by a documented barrier: linear reductions by the entropy/error-weight argument of Lu \etal\ \cite{LPQR26}, degree-$D<n$ polynomial reductions by the Reed--Muller distance of $\mathbf{1}_L$ (which has degree exactly $n$, so no lower-degree polynomial can represent it), and entropy-compressing reductions by the $\Omega(n^2)$ secret-entropy floor (Proposition~\ref{prop:entropy-floor}). We do not claim a one-to-one correspondence between invariants and barriers; rather, no proposed reduction class survives all three.
\end{conjecture}

\Cref{conj:source} is falsifiable: a single explicit reduction (membership-LSN, batch-LSN, or sympLPN) $\to$ LPN that preserves any of the three invariants would refute the conjecture in its current form. It is not decidable by reduction analysis alone; it rests on the absence of such a reduction and the structural barriers documented above.

\paragraph{The scrambling barrier at constant noise.} The LPQR26 linear barrier (Theorem D.2) applies at any noise rate $p=\omega(1/n)$, and at our constant noise $p=1/4$ its error-amplification step is in fact \emph{stronger}: a randomizing row of weight $w$ leaves per-bit bias $(1-2p)^w = 2^{-w}$ (piling-up), versus $\approx 1-2wp$ in the low-noise regime. The residual gap is therefore not about noise: linear reductions consuming $m=\omega(n)$ samples are ruled out in the conditional-uniformity model for all $m=\operatorname{poly}(n)$ by the reachability barrier (\Cref{thm:any-b-uniformity}).  The marginal-adaptive corner, in which $B=g(A)$ yields marginally-uniform $BA$ via low-weight rows that vary with $A$, is not covered by this argument and remains open (Open Problem~\ref{open:marginal-adaptive}).  What remains outside both regimes is only non-linear scramblers. For (ii) an algebraic obstruction remains: $S_A=0$ is a deterministic quadratic relation among the public vectors (our observation, distinct from the LPQR entropy argument). Direct inspection of the $n=3$ case shows that simple scramblers fall into one of three classes: a linear scrambler \emph{transports} the quadratic relation ($S'_{A'}=0$ with respect to the conjugated form $M^{-\top}\Omega M^{-1}$, so it remains detectable), a non-linear scrambler destroys the linear label structure required by LPN, and a random scrambler destroys the secret entirely. None of these simple classes simultaneously destroys $S_A=0$ and creates a fresh LPN secret-bit structure. The residual question --- whether a sufficiently complex non-linear scrambler exists --- is the randomized marginal-adaptive question governed by \Cref{lem:m2}.

A Fourier reformulation sharpens this residual. For any boolean scrambler $z:\F_2^{2n}\to\F_2$, the correlation of $(-1)^{z(y)}$ with a candidate LPN bit $\langle c,x\rangle$ on $y=Ax+e$ is exactly
\[
\mathrm{bias}(c)=\E_{x,e}\bigl[(-1)^{z(Ax+e)}(-1)^{\langle c,x\rangle}\bigr]=\sum_{S:\,A^{\top}S=c}\hat z(S)\,(1-2p)^{|S|}
\]
(expand $(-1)^z$ in characters; $\E_x$ annihilates the terms with $A^{\top}S\neq c$ and $\E_e$ contributes $(1-2p)^{|S|}$), so a usable LPN bit is Fourier mass of the scrambler on \emph{low-weight} $S$ inside a coset of the self-dual code $\ker(A^{\top})=\Omega\,\mathrm{Col}(A)$. This recasts the search over non-linear maps as a coset-Fourier-concentration problem and isolates where non-linearity can help. The optimal single-bit scrambler is the maximum-likelihood bit-decoder $z^\ast(y)=\operatorname{sign}\sum_{A^{\top}S=c}(1-2p)^{|S|}(-1)^{\langle S,y\rangle}$, so linearity is \emph{not} per-bit optimal. At constant noise, however, a Cauchy--Schwarz/Parseval bound caps the reach of non-linearity: $\max_z|\mathrm{bias}(c)|\le\bigl(\sum_{S:\,A^{\top}S=c}(1-2p)^{2|S|}\bigr)^{1/2}$, and the average over $c$ of the quantity under the root is $\bigl(\tfrac12(1+(1-2p)^2)^2\bigr)^{n}$ ($\approx 0.78^{n}$ at $p=\tfrac14$) $\to 0$, so for all but an exponentially small fraction of targets the best achievable bias is exponentially small, and only $c$ with an anomalously low-weight coset leader are reachable. (This is specific to constant noise: at $p<0.11$ the rate-$\tfrac12$ code lies below capacity, an ML decoder recovers $x$ outright, and \emph{any} $c$ becomes reachable --- a regime outside ours.) The achievable low-noise targets are thus essentially $\{A^{\top}S:|S|\text{ small}\}$, a property of $A$: indeed the achievable coset-leader-weight \emph{profile} of an isotropic $A$ is itself a detectable, $n$-robust distributional signature against a generic full-rank matrix (total variation $\approx 0.2$, rising to $\approx 0.23$ by $n=5$, well above the finite-sample baseline; \emph{evidence}, exact small $n$). But the reduction's output does not betray it: the output matrix $C=BA$ is decoder-independent (an ML decoder for a target $c$ leaves $c$, hence $C$, unchanged), and the natural quadratic-form invariant $\dim\{M:\,C^{\top}MC=0\}=\binom{2n}{2}-\binom{n}{2}$ coincides for isotropic and generic $A$ --- it depends only on $\operatorname{rank}A=n$ (the constraint is $M|_{\mathrm{Col}A}=0$), while the \emph{transported} form needs the hidden $A$. The detectable obstruction is therefore the correlated noise $Be$, not the matrix $C$: the residual cell collapses onto \Cref{lem:m2}, which non-linearity --- confined at constant noise to low-weight cosets by the bias bound above --- does not relax.

\paragraph{The marginal-adaptive corner: entropy-support bound.}
The gap between conditional and marginal uniformity is bridged by an entropy-support argument that does not require per-instance uniformity.  The idea is simple: the isotropic basis $A$ carries entropy $H(A)=\Theta(n^{2})$, while each low-weight row of $B$ can only ``spend'' $O(n)$ bits of entropy on its output query.  If the total output entropy is $nm$ (uniformity), the budget forces most rows to have high weight.

\begin{lemma}[Entropy-support bound for low-weight rows]
\label{lem:m1}
Let $A\in\F_2^{2n\times n}$ be a random isotropic basis and let $B=g(A,R)\in\F_2^{m\times 2n}$ be any randomized function of $A$.  Let $C=BA$ and suppose $\operatorname{SD}(C,\,\operatorname{Uniform}_{\F_2^{m\times n}})\le\delta$.  Let $w=\lfloor 0.19n\rfloor$.  Then the expected number of rows of $B$ with weight $\le w$ satisfies
\[
\E\bigl[\,\bigl|\{\,i:|b_i|\le w\,\}\bigr|\,\bigr]
\;\le\;
16n+11\delta m+\frac{11m}{n}+O(1),
\]
\end{lemma}

\begin{proof}
By the Fannes--Csisz\'{a}r continuity bound (Cover--Thomas \S17.3.3), $\operatorname{SD}(C,\text{Uniform})\le\delta$ implies $H(C)\ge nm-\delta nm-1$.  Conditioning on $A$ and using the chain rule,
\[
H(C)\;\le\;H(A)+H(C\mid A)\;=\;H(A)+\sum_{i=1}^{m}H(c_i\mid A,c_{<i})
\;\le\;H(A)+\sum_{i=1}^{m}H(c_i\mid A).
\]
Let $T_i$ be the indicator of the event $|b_i|\le w$.  Then
\begin{align*}
H(c_i\mid A)
&\;\le\;H(T_i,c_i\mid A)
\;=\;H(T_i\mid A)+H(c_i\mid T_i,A) \\
&\;\le\;1+\Pr[T_i=1]\cdot0.906n+\Pr[T_i=0]\cdot n,
\end{align*}
where the first term is at most $1$ because $T_i$ is binary, the second uses \Cref{lem:reachability} ($|R_w(A)|\le2^{0.906n}$), and the third uses the trivial bound $H(c_i\mid A)\le n$.  Simplifying,
\[
H(c_i\mid A)\;\le\;1+n-\Pr[T_i=1]\cdot0.094n.
\]
Summing over $i$ and letting $k=\sum_i T_i$,
\[
\sum_{i=1}^{m}H(c_i\mid A)
\;\le\;m+nm-\E[k]\cdot0.094n.
\]
The entropy of $A$ is at most the log-support of the isotropic-frame distribution, $H(A)\le (3/2)n^{2}+n/2+O(1)$.  Combining with the Fannes--Csisz\'{a}r bound $H(C)\ge nm-\delta nm-1$ and the chain rule,
\[
nm-\delta nm-1
\;\le\;(3/2)n^{2}+n/2+m+nm-\E[k]\cdot0.094n+O(1),
\]
which rearranges to $\E[k]\cdot0.094n\le(3/2)n^{2}+n/2+m+\delta nm+O(1)$.  Since $(3/2)n^{2}+n/2$ divided by $0.094n$ is at most $16n+O(1)$, and $\E[k]$ is already the (global) expected number of low-weight rows, this is the claimed bound.
\end{proof}

\Cref{lem:m1} bounds the \emph{expected} number of low-weight rows in the marginal-adaptive model using only marginal uniformity and the entropy of $A$.  When $m\ge Cn$ with $C>16$ and $\delta=o(1)$, the lemma gives at least $m-16n-o(m)$ rows of weight $>0.19n$ \emph{in expectation}, regardless of how $B$ depends on $A$.  The rowwise noise-amplification of \Cref{cor:noise-amp} applies to those high-weight rows; the recovery bound of \Cref{cor:recovery-barrier}, however, assumes \emph{independent} per-sample noise, so the genuinely correlated $Be$ is deferred to \Cref{lem:m2}.

\begin{conjecture}[Distribution-mismatch death]
\label{lem:m2}
\textbf{Status: heuristic, not yet established.}  The obstacle is that the output noise components $e'_i=\langle b_i,e\rangle$ are $m$ linear images of the fixed $2n$-bit symplectic-LPN noise vector $e$; they therefore live in a subspace of dimension at most $2n$ and are heavily correlated when $m=\omega(n)$.  Both standard repair strategies hit walls: (i)~a recovery-reduction argument needs \Cref{cor:recovery-barrier} to be robust to correlated noise, but the $\le 2n$ leak-rows could over-determine $x\in\F_2^{n}$ (and the solver does not know $B$); (ii)~a noise-entropy argument $H(Be)\le H(e)=O(n)$ holds for fixed $B$ but fails for random high-entropy $B$ (since $Be$ can be near-uniform).  The $\le 2n$-dimensional noise structure is both why the corner is plausibly closed and why no clean argument closes it yet.

\emph{If} the following statement could be proved, it would close the marginal-adaptive corner: Let $p'\in(0,1/2)$ with $1-2p'\ge 1/\operatorname{poly}(n)$.  Consider any reduction outputting $m$ samples $(c_i,a_i)$ with $a_i=\langle c_i,x\rangle+e'_i$ where $e'_i$ has bias $\le 2^{-0.19n}$ on all but $\ell$ rows with $\mathbb E[\ell]\le 16n+o(m)$ (the expectation \Cref{lem:m1} delivers).  Then
\[
\operatorname{SD}\bigl(\,(C,y),\;\mathrm{LPN}_{p'}\,\bigr)
\;\ge\;
1-o(1)
\]
provided $m\ge \frac{4n}{(1-2p')^{2}}\cdot\operatorname{polylog}(n)$.
\end{conjecture}

\begin{theorem}[Deterministic marginal-adaptive barrier]
\label{thm:deterministic-marginal-adaptive}
Let $A\sim\operatorname{Uniform}(\mathcal{A}_n)$ be a uniform ordered isotropic frame --- the public object the reduction actually receives (\Cref{def:symplpn}) --- let $B=g(A)\in\F_2^{m\times 2n}$ be any \emph{deterministic} function of $A$, and set $C=BA$. Then for every noise rate $p'$,
\[
\operatorname{SD}\bigl((C,y),\;\mathrm{LPN}_{p'}\bigr)
\;\ge\;
\operatorname{SD}\bigl(C,\;\operatorname{Uniform}_{\F_2^{m\times n}}\bigr)
\;\ge\;
1-\frac{|\mathcal{A}_n|}{2^{mn}},
\]
where $\log_2|\mathcal{A}_n|=\log_2\!\bigl(|\Lagr(2n,\F_2)|\cdot|\mathrm{GL}(n,\F_2)|\bigr)=\tfrac32 n^2+\tfrac n2+O(1)$ is the frame entropy. This is $1-o(1)$ as soon as $mn$ exceeds that entropy (e.g.\ for any $m\ge 2n$). The deterministic marginal-adaptive cell is therefore closed unconditionally.
\end{theorem}

\begin{proof}
$C=g(A)A$ is a deterministic function of the ordered frame $A$, so its support has size at most $|\mathcal{A}_n|=|\Lagr(2n,\F_2)|\cdot|\mathrm{GL}(n,\F_2)|$ (the reduction may depend on the frame, not merely the subspace it spans), whereas the LPN public matrix is uniform on $2^{mn}$ values; the statistical distance between a distribution supported on $K$ points and the uniform distribution on $N\ge K$ points is at least $1-K/N$. Since $(C,y)$ determines $C$, the data-processing inequality carries the bound to $\operatorname{SD}((C,y),\mathrm{LPN}_{p'})$. For $n=2$, $|\mathcal{A}_2|=15\cdot 6=90$, so the bound is $1-90/4^m$: it exceeds $\tfrac12$ at $m=4$ ($1-90/256=0.648$) and $\to 1$ as $m$ grows, by exact enumeration.
\end{proof}

The deterministic barrier leaves the \emph{randomized} marginal-adaptive cell, whose hard core is the correlated, low-dimensional noise $Be$ (\Cref{lem:m2}). A broad slice of that cell --- every reduction whose conditional noise occupies sub-exponentially many low-dimensional cages --- closes unconditionally, with the asymptotic-in-$n$ rate settled, by an elementary support-confinement argument.

\begin{lemma}[Confined-noise closure]
\label{lem:confined-g}
Consider a randomized marginal-adaptive reduction outputting $C=BA$ and $y=Cx+Be$ with $\operatorname{SD}(C,\operatorname{Uniform})\le\delta$. For each output $C$, the conditional law of the noise $Be$ is supported on the union of the \emph{cages} $\operatorname{Col}(B)$ --- the column space of each matrix $B$ the reduction can produce alongside $C$, necessarily of dimension $\le 2n$ (as $B\in\F_2^{m\times 2n}$) and containing $\operatorname{Col}(C)=\operatorname{Col}(BA)$, and containing every $Be$ since the noise is $B$ applied to $e$ --- let $N_C$ be their number and $\bar N=\mathbb{E}_C[N_C]$. Then for every $p'\in(0,\tfrac12)$,
\[
\operatorname{SD}\bigl((C,y),\,\mathrm{LPN}_{p'}\bigr)\;\ge\;1-\bar N\cdot 2^{2n}(1-p')^{m}-\delta .
\]
In particular, once $m\ge cn$ with $c>2/\log_2(1/(1-p'))$ and the reduction uses only $\bar N=2^{o(n)}$ cages, $\operatorname{SD}=1-o(1)$. The \emph{fixed transverse generator} case --- the action of $B$ off $L=\operatorname{Col}(A)$ a constant $G$, so $N_C\equiv 1$ --- gives the clean bound $1-2^{2n}(1-p')^m-\delta$ ($c>4.82$ at $p'=\tfrac14$, $c\to2$ as $p'\to\tfrac12$); the asymptotic-in-$n$ rate is settled for every bounded-cage sub-class.
\end{lemma}

\begin{proof}
As $C$ is within $\delta$ of uniform in both distributions, $\operatorname{SD}((C,y),\mathrm{LPN}_{p'})\ge \mathbb{E}_{C}\bigl[\operatorname{SD}(y\mid C,\,y'\mid C)\bigr]-\delta$ with $C$ uniform. Fix a full-rank $C$ (the rank-deficient $C$ are a $2^{-\Theta(m)}$ fraction and only help). The secret $x$ is uniform and independent of the noise, so $P(y\mid C)=2^{-n}\nu(y+\operatorname{Col}(C))$, where $\nu$ is the law of $Be$, and $P(y'\mid C)=2^{-n}\mu(y+\operatorname{Col}(C))$ with $\mu=\operatorname{Bernoulli}(p')^{m}$; both are constant on cosets of $\operatorname{Col}(C)$, whence $\operatorname{SD}(y\mid C,y'\mid C)=\operatorname{SD}\bigl(\nu\bmod\operatorname{Col} C,\ \mu\bmod\operatorname{Col} C\bigr)$. Every realization has $Be\in\operatorname{Col}(B)$ for the matrix $B$ it used (since $Be=Cw+B|_{L'}v$ with $w,v\in\F_2^n$ the $L$- and $L'$-coordinates of $e$, and $\operatorname{Col}(C)\subseteq\operatorname{Col}(B)$), so $\nu$ is supported on the union $\mathcal{V}=\bigcup_{i=1}^{N_C}V_i$ of the $N_C$ cages $V_i=\operatorname{Col}(B_i)$, each of dimension $\le 2n$. Hence, by a union bound,
\[
\sum_S\min(\nu_S,\mu_S)\ \le\ \mu(\mathcal{V})\ \le\ \sum_{i=1}^{N_C}\Pr_{e''\sim\operatorname{Bernoulli}(p')^m}\!\bigl[e''\in V_i\bigr]\ \le\ N_C\cdot 2^{2n}(1-p')^{m},
\]
the last step because $\dim V_i\le 2n$ and every Bernoulli mass $(p')^{\mathrm{wt}}(1-p')^{m-\mathrm{wt}}\le(1-p')^{m}$ for $p'<\tfrac12$. Therefore $\operatorname{SD}(y\mid C)\ge 1-N_C\cdot 2^{2n}(1-p')^{m}$ for every full-rank $C$, and averaging over $C$ gives the claim.
\end{proof}

The lemma reduces the open core to a single combinatorial quantity: the number $\bar N$ of low-dimensional noise-cages a reduction can hang on one output $C$. It closes every reduction with sub-exponentially many cages ($\bar N=2^{o(n)}$) --- the fixed-generator case ($N_C\equiv 1$), and any bounded-diversity transverse randomization. The genuine residual --- the open core of \Cref{lem:m2} --- is therefore exactly the reductions that \emph{diffuse} the noise across $2^{\Theta(n)}$ cages, by letting the transverse generator $G=g(A,R)$ depend on $A$ or on private randomness (uniform-$B$ is the extreme: $\operatorname{Col}(B)$ then ranges over essentially all $2n$-dimensional spaces above $\operatorname{Col}(C)$). For those the support bound is vacuous. The residual obstruction for the diffuse half is then \Cref{lem:m2} itself --- the output must stay far from a \emph{usable} $\mathrm{LPN}_{p'}$ --- and \emph{not} an information bound on $x$: a mechanism-level computation at $n=2,3,4$ indicates that the secret in fact leaks $\Theta(n)$ bits ($I(x;y\mid C)=\Theta(n)$), since the syndrome rank is capped at $n$ (the identity $\operatorname{rank}(HB)=\operatorname{rank}(B)-\operatorname{rank}(BA)\le n$) and the complementary $n$ message coordinates stay biased even given the syndrome. That leakage is symplectic- rather than LPN-structured, so it cannot drive an LPN solver (\Cref{open:marginal-adaptive}); the operative barrier remains statistical distance. The cell thus splits three ways, by the \emph{type} of the output law: the \emph{confined} part (few cages; closed by \Cref{lem:confined-g}), the \emph{near-uniform} diffuse part (closed by \Cref{prop:diffuse-solvable} below --- in particular the extreme uniform-$B$ strategy), and a \emph{concentrated} diffuse residual (outputs neither cage-confined nor near-uniform), which remains governed by \Cref{lem:m2}. We \emph{conjecture} this residual cell is also empty --- no linear reduction maps symplectic LPN to a usable LPN instance --- but cannot yet prove it; the obstruction is genuinely \Cref{lem:m2}. The supporting intuition is a win--win: a noise spread enough to evade \Cref{lem:confined-g} pushes the output toward uniform, hence (by \Cref{prop:diffuse-solvable}) away from any \emph{solvable} usable $\mathrm{LPN}_{p'}$; while a noise concentrated enough to mimic a solvable target must absorb the heavy rows forced by \Cref{lem:m1} into a codeword, re-masking the secret by an $\Omega(n)$-entropy shift (the symplectic leakage above). In fact, on the bounded-rate part of the residual band --- $c=m/n$ with $m\ll 2^n$ and $\delta_{\mathrm{GV}}(1/c)<\tfrac{7}{16}$ (equivalently $1-1/c<H_2(\tfrac7{16})$, i.e.\ $c<88.5$) --- the concentrated horn can be pushed to the level of the noise \emph{rate}. The exact per-coordinate output rate is $q=\frac1m\sum_i\frac{1-2^{-\operatorname{wt}(r_i)}}{2}$ in the row weights of $B$ (since $(Be)_i=\langle r_i,e\rangle$, $e\sim\operatorname{Bernoulli}(\tfrac14)$), so a \emph{usable} (solvable) target rate needs $q\le\delta_{\mathrm{GV}}(1/c)$, where $H_2(\delta_{\mathrm{GV}})=1-1/c$. A deterministic row-dependency $\sum_{i\in S}r_i=0$ forces $\sum_{i\in S}C_i=\bigl(\sum_{i\in S}r_i\bigr)A=0$ for \emph{every} $A$ (the rows of $C=BA$ are $C_i=r_iA$), so $C$ carries a weight-$|S|$ dependency; but a uniform $m\times n$ matrix (constant $c$, $m\ll 2^n$) has, except with probability $\sum_{w\le K\log n}\binom{m}{w}2^{-n}=o(1)$, no row-dependency of weight $\le K\log n$. Now $q\le\delta_{\mathrm{GV}}(1/c)<\tfrac7{16}$ makes weight-$\ge3$ rows (cost $\ge\tfrac7{16}$ per coordinate) unaffordable in bulk, forcing at least $m\bigl(1-\tfrac{16}{7}\delta_{\mathrm{GV}}(1/c)\bigr)=\alpha(c)\,n$ rows to weight $\le2$, with $\alpha(c)=c\bigl(1-\tfrac{16}{7}\delta_{\mathrm{GV}}(1/c)\bigr)$; on the constant-rate subrange where $\alpha(c)>2$ --- numerically $3.8\lesssim c\lesssim 56.9$ --- the count exceeds $(1+\varepsilon)(2n{+}1)$. Representing weight-$1$ rows as edges to an auxiliary root and weight-$2$ rows as ordinary edges on $2n$ vertices, the augmented graph then has average degree $>2+\varepsilon$ unless an immediate repeated-edge or zero-row dependency already occurs (a repeated unit or weight-$2$ row is itself a weight-$2$ dependency; a zero row gives a zero $C$-row, of which a uniform $C$ has only $m2^{-n}$), so the irregular Moore bound~\cite{AHL02} yields an $O_\varepsilon(\log n)$ cycle --- in every case an $O(\log n)$ dependency in $C$ that the uniform matrix lacks. This rules out light matched-rate outputs only in that $\alpha(c)>2$ subrange ($q>\delta_{\mathrm{GV}}(1/c)$ is forced there); outside it the rate argument does \emph{not} force $q>\delta_{\mathrm{GV}}(1/c)$, so even the light-rate exclusion remains part of the open \Cref{lem:m2}. Within the subrange the residual is whether the resulting \emph{heavy} ($q>\delta_{\mathrm{GV}}$), $\le 2n$-confined noise is distinguishable from a usable $\mathrm{LPN}_{p'}$ --- the over-dispersion question, on which the min-syndrome-weight scalar statistic saturates near the covering radius at the marginally-solvable $p'\approx\delta_{\mathrm{GV}}$; moreover an exact computation finds the statistic's over-dispersion---the ratio of its variance under the confined $\le 2n$-dimensional noise to that under the matched-rate target---\emph{decreasing} toward $1$ as $n$ grows for consistent-rank weight-$2$ generators (from $\approx 3.9$ at $n=2$ to $\approx 1.6$ at $n=5$, exact through $n=4$), so its distinguishing advantage vanishes and this scalar statistic cannot by itself close \Cref{lem:m2}---any closure must come from the full statistical distance. Finite-$n$ evidence from that full distance points the same way: for a fixed heavy weight-$2$ generator (the first $m$ weight-$2$ rows of $\F_2^{2n}$ in integer order) at $m=4n$, the per-output $\operatorname{SD}\bigl(\nu\bmod\operatorname{Col}C,\ \mu_{\delta_{\mathrm{GV}}}\bmod\operatorname{Col}C\bigr)$ to the \emph{marginally-solvable} $p'=\delta_{\mathrm{GV}}$ LPN target, averaged over Lagrangians, is $0.850,\,0.939,\,0.973,\,0.992$ for $n=2,3,4,5$ ($n\le 3$ exact; $n=4,5$ Monte-Carlo over Lagrangians), increasing over the tested range and consistent with the rigorous confined-cell closure of \Cref{lem:confined-g} and with \Cref{lem:m2}. Within each cage, moreover, the Bernoulli pushforward converges to the cage-uniform law in Lagrangian-averaged $L_2$ at the explicit rate $\mathbb{E}_L\!\sum_{v\in\Omega L\setminus 0}(1-2p)^{2\operatorname{wt}(v)}=(25/32)^n-\tfrac{(25/32)^n+1}{2^n+1}\to 0$ (isotropic MacWilliams; exact at $n\le4$), so the residual obstruction is the cage-\emph{level} distribution rather than any within-cage structure. After \Cref{lem:confined-g} and \Cref{prop:diffuse-solvable}, only the high-diffusion concentrated regime ($\bar N=2^{\Omega(n)}$ cages, non-uniform), where that support-confinement turns vacuous, is left. Neither horn is captured by the support-counting or spectral arguments available here --- which is precisely why \Cref{lem:m2} stays open.

The upper half of this $\Theta(n)$ leakage is in fact unconditional, and pins the secret-recovery (search) barrier exactly:

\begin{lemma}[Secret-entropy floor: search non-recovery]
\label{lem:secret-floor}
For every $B\in\F_2^{m\times 2n}$ with $\operatorname{rank}(BA)=n$, secret $x\sim\operatorname{Uniform}(\F_2^n)$, and noise $e\sim\operatorname{Bernoulli}(p)^{2n}$ independent of $x$, the output $(C,y)=(BA,\,Cx+Be)$ satisfies
\[
H(x\mid C,y)\ \ge\ \bigl(2H_2(p)-1\bigr)n,\qquad\text{equivalently}\qquad I(x;y\mid C)\ \le\ 2\bigl(1-H_2(p)\bigr)n .
\]
At $p=\tfrac14$: $H(x\mid C,y)\ge 0.623\,n$ and $I(x;y\mid C)\le 0.377\,n$. The bound holds for \emph{every} $B$ --- the $\operatorname{rank}(B)$ dependence cancels --- so no marginal-adaptive linear reduction \emph{recovers} $x$.
\end{lemma}

\begin{proof}
Fix $(A,B)$, so $C=BA$ is determined. Write $\sigma=Be\bmod\operatorname{Col}(C)\in\operatorname{Col}(B)/\operatorname{Col}(C)$ for the syndrome and $v\in\F_2^n$ for the in-code coordinate, so $Be=Cv+r_\sigma$ with $r_\sigma$ a fixed coset representative; the pair $(\sigma,v)$ determines $Be$. Since $y=Cx+Be=C(x+v)+r_\sigma$, both $\sigma=y\bmod\operatorname{Col}(C)$ and the sum $x+v$ (from $C(x+v)=y+r_\sigma$, $C$ injective) are determined by $(C,y)$; as $x$ is uniform and independent of $e$, $x+v$ is uniform and carries no information about $v$ given $\sigma$. Hence
\[
H(x\mid C,y)=H(v\mid \sigma)=H(Be)-H(\sigma).
\]
Now $H(Be)=H(e)-H(e\mid Be)\ge 2nH_2(p)-(2n-\operatorname{rank}B)$, since the fibre of $e\mapsto Be$ is a coset of $\ker B$ of dimension $2n-\operatorname{rank}(B)$; and $H(\sigma)\le\dim\bigl(\operatorname{Col}(B)/\operatorname{Col}(C)\bigr)=\operatorname{rank}(B)-n$ by the rank identity $\operatorname{rank}(HB)=\operatorname{rank}(B)-\operatorname{rank}(BA)$ used above. Subtracting, the $\operatorname{rank}(B)$ terms cancel:
\[
H(x\mid C,y)\ \ge\ \bigl[2nH_2(p)-(2n-\operatorname{rank}B)\bigr]-\bigl(\operatorname{rank}(B)-n\bigr)\ =\ \bigl(2H_2(p)-1\bigr)n.
\]
Averaging over the frame ensemble preserves the bound, and $I(x;y\mid C)=H(x\mid C)-H(x\mid C,y)=n-H(x\mid C,y)$ gives the mutual-information form. (At $n=2$ the floor $0.623n$ is matched in form by the leakiest --- full-rank --- $B$, for which $H(x\mid C,y)=nH_2(p)=1.6226$; rank-deficient $B$ only raise it. Exact check, $n=2$.)
\end{proof}

Thus the leakage obeys $\Omega(n)\le I(x;y\mid C)\le 2(1-H_2(p))n$: it is genuinely $\Theta(n)$, yet a constant fraction $(2H_2(p)-1)n$ of the secret's entropy \emph{never} leaves the transcript. This makes rigorous the ``re-masking by an $\Omega(n)$-entropy shift'' horn above, and confirms that the leakage --- symplectic- not LPN-structured --- cannot drive secret recovery; any residual reduction must therefore work at the level of statistical distance (\Cref{lem:m2}), not secret recovery.

\begin{proposition}[Diffuse closure at solvable sample size]
\label{prop:diffuse-solvable}
Suppose a marginal-adaptive reduction's output is $o(1)$-close to uniform, $\operatorname{SD}\bigl((C,y),\operatorname{Uniform}\bigr)=o(1)$. Then for every usable $p'$ ($1-2p'\ge1/\operatorname{poly}(n)$) and every sample size $m$ at which $\mathrm{LPN}_{p'}$ is solvable with a diverging \emph{smooth} capacity margin --- $m(1-H_2(p'))-n-\omega\bigl(\sqrt{m\,V(p')}+\log m\bigr)\to\infty$, where $V(p')=p'(1-p')\bigl(\log_2\tfrac{1-p'}{p'}\bigr)^2$ is the noise varentropy (for $p'$ bounded away from $\tfrac12$ this is the first-order $\bigl(m-n/(1-H_2(p'))\bigr)(1-H_2(p'))\to\infty$) ---
\[
\operatorname{SD}\bigl((C,y),\mathrm{LPN}_{p'}\bigr)\ \ge\ 1-2^{\,n-m(1-H_2(p'))+\omega(\sqrt{m\,V(p')})}-o(1)\ =\ 1-o(1).
\]
Uniform-$B$ satisfies the hypothesis: $\nu_{Be\mid C}\bmod\operatorname{Col}(C)=q_{\mathrm{graph}}(n)\,\delta_0+(1-q_{\mathrm{graph}}(n))\,\operatorname{Uniform}(\F_2^{m-n})$ with $q_{\mathrm{graph}}(n)\to0$, since for $e\notin L$ the rows of $[\,C\mid Be\,]=B[\,A\mid e\,]$ are i.i.d.\ uniform on $\F_2^{n+1}$ (as $[A\mid e]$ has rank $n+1$). The extreme diffuse strategy is therefore closed.
\end{proposition}

\begin{proof}
A useful reduction must output a \emph{solvable} $\mathrm{LPN}_{p'}$ instance, so $m\ge n/(1-H_2(p'))$: below this the $2^n$ noise balls of radius ${\sim}p'm$ have total volume $2^{n}\cdot 2^{mH_2(p')}>2^m$, so $x$ retains $\Omega(n)$ bits of conditional entropy --- not even \emph{partial} recovery sufficient to break the source is possible --- and no reduction results. At such $m$, fix any slowly diverging $t_m\to\infty$ with $t_m\sqrt{m\,V(p')}=o\bigl(m(1-H_2(p'))-n\bigr)$ (possible by the smooth-margin hypothesis). The law of $\mathrm{LPN}_{p'}$ given $C$ is then supported, up to $o(1)$ mass, on its \emph{smooth} typical set of size $\le 2^{\,n+mH_2(p')+t_m\sqrt{m\,V(p')}}$ --- the window $t_m\sqrt{m\,V(p')}$ diverges relative to the standard deviation $\sqrt{m\,V(p')}$ of the information density $\log_2(1/\text{mass})$, so by Chebyshev it captures all but $o(1)$ of the mass (this second-order term is negligible for constant $p'$ but must be carried when $p'\to\tfrac12$). A distribution $o(1)$-close to uniform places mass at most $2^{\,n-m(1-H_2(p'))+t_m\sqrt{m\,V(p')}}+o(1)=o(1)$ on any such set, so the overlap is $o(1)$ and $\operatorname{SD}=1-o(1)$. For usable $p'$ ($1-2p'\ge1/\operatorname{poly}(n)$, e.g.\ $1-2p'=1/n$) the smooth margin is met at $m=\operatorname{poly}(n)$; only $p'$ super-polynomially close to $\tfrac12$ ($1-2p'<1/\operatorname{poly}(n)$) is excluded as not usable.
\end{proof}

The vanishing of all distinguishers at the small ratios $m/n=O(1)$ examined elsewhere is consistent with \Cref{prop:diffuse-solvable}: those sample sizes lie \emph{below} the solvability threshold $n/(1-H_2(p'))$ (for $p'=\tfrac14$, $m/n<5.3$), where the LPN target is itself near-uniform, so the two near-uniform laws agree --- but such instances are unsolvable and yield no reduction. The closure lives at the larger sample size any useful reduction must reach.

\begin{theorem}[Randomized marginal-adaptive linear-reduction barrier --- \textbf{conditional}]
\label{thm:marginal-adaptive}
\textbf{Depends on \Cref{lem:m2}.}  Let $A\sim\operatorname{Uniform}(\mathcal{A}_n)$ be a uniform ordered isotropic frame (a uniform ordered basis of a uniform Lagrangian) and let $B=g(A,R)\in\F_2^{m\times 2n}$ be any randomized adaptive matrix (using fresh private randomness $R$).  \Cref{lem:m1} shows that if $\operatorname{SD}(BA,\text{Uniform})\le\delta$ with $\delta=o(1)$ and $m\ge Cn$ ($C>16$), then \emph{in expectation} at least $m-16n-o(m)$ rows have weight $>0.19n$ (all but $16n+o(m)$ rows; this is $m-o(m)$ when $m=\omega(n)$), hence effective noise bias $\le 2^{-0.19n}$.  If \Cref{lem:m2} holds, then for every $p'$ with $1-2p'\ge 1/\operatorname{poly}(n)$ and in the sample regime $m\ge 4n/(1-2p')^2\cdot\operatorname{polylog}(n)$, $\operatorname{SD}((C,y),\mathrm{LPN}_{p'})=1-o(1)$, so no marginal-adaptive linear reduction maps symplectic LPN to a usable standard LPN instance in that regime.
\end{theorem}

\Cref{thm:marginal-adaptive} is a \textbf{conditional} barrier: \Cref{lem:m1} is proven, but \Cref{lem:m2} remains open because of the correlated-noise obstacle described above.  The deterministic marginal-adaptive case is already closed unconditionally by \Cref{thm:deterministic-marginal-adaptive}. The status map is therefore: fixed-$B$ \textbf{ruled out}, public-$B$ \textbf{ruled out}, conditionally-uniform adaptive \textbf{ruled out}, deterministic marginal-adaptive \textbf{ruled out}, randomized marginal-adaptive \textbf{open} --- with its bounded-cage slice (conditional noise on $2^{o(n)}$ cages, including every fixed-transverse-generator reduction) now \textbf{ruled out} by \Cref{lem:confined-g}, leaving only the diffuse slice ($2^{\Theta(n)}$ cages), which remains governed by the still-open statistical-distance bound \Cref{lem:m2} (\emph{not} by an information bound on $x$, which is $\Theta(n)$ rather than $o(n)$; see \Cref{open:marginal-adaptive}).  The remaining regime $m=\Theta(n)$ with $m<Cn$ is covered by LPQR26 Appendix~D --- their \emph{fixed-$B$} entropy-deficiency bound \textbf{Theorem~D.1} (their Theorem~D.2 instead treats randomized $B$ with $BA$ induced-uniform) --- for fixed $B$ of any rank, including low-rank.

\section{Open Problems}
\label{sec:open}

We state precisely the problems whose resolution would close the standard-model gap. We first record, as a numbered proposition, one sub-question our techniques settle completely in the negative --- whether public rerandomization can manufacture \emph{fresh} samples --- since the corresponding open problem (item~6 below) is only its residual.

\begin{proposition}[Sample-freshness identity for label-preserving splits]
\label{prop:freshness}
Fix a noise rate $p \in (0,\tfrac12)$ (the constant-noise regime, e.g.\ the LSN rate $p=\tfrac14$; at $p=\tfrac12$ the labels are uninformative and the equality condition below degenerates), and split a single LSN sample $(x,b)$ (secret Lagrangian $L$, label $b = \mathbf 1_L(x)\oplus e$ with $e \sim \mathrm{Bernoulli}(p)$) into two via public bijections $f_1, f_2$ of $\F_2^{2n}$ as $(x,b)\mapsto((f_1 x, b),(f_2 x, b))$. Then the statistical distance between this pair and two independent fresh same-secret samples is exactly
\[
\mathrm{SD} \;=\; 1 - \frac{2p(1-p) + (1-2p)^2\,A(f_1,f_2)}{4^{n}},
\qquad A(f_1,f_2) = \Pr_{L,x}\bigl[\mathbf 1_L(f_1 x) = \mathbf 1_L(f_2 x)\bigr] \le 1,
\]
hence $\mathrm{SD} \ge 1 - (p^2+(1-p)^2)/4^n$ with equality iff $f_1 = f_2$. The same lower bound holds verbatim for label-flipping splits $(x,b)\mapsto((f_1 x,\,b\oplus h_1(x)),(f_2 x,\,b\oplus h_2(x)))$ with arbitrary public $h_i:\F_2^{2n}\to\F_2$, with equality only at the literal duplicate $f_1=f_2,\ h_1=h_2$. For the symplectic-orbit family $f_i = S_i \in \mathrm{Sp}(2n,\F_2)$ (label-preserving, rerandomizing the secret to $S_i\cdot L$), comparing instead against fresh samples \emph{for the rerandomized secrets} $S_i\cdot L$ --- the natural target of rerandomization --- the membership values $\mathbf 1_{S_iL}(S_ix)=\mathbf 1_L(x)$ coincide, so the distance attains the duplicate value
\[
\mathrm{SD} \;=\; 1 - \frac{p^2+(1-p)^2}{4^{n}} \;=\; 1 - \frac{5}{8\cdot 4^{n}} \quad (p = \tfrac14),
\]
i.e.\ $123/128$ at $n=2$, $507/512$ at $n=3$, tending to $1$ (numerically the same as the same-secret lower bound).
\end{proposition}

\begin{proof}
For the same-secret comparison, the transformed pair carries one shared label $c=\mathbf 1_L(x)\oplus e$ across both points $f_1x,f_2x$, so it overlaps the product fresh law only where the two fresh labels agree and are consistent with a single noise draw; evaluating this overlap yields the displayed identity, with $A$ the probability that the membership values $\mathbf 1_L(f_1x),\mathbf 1_L(f_2x)$ coincide. Since $\mathbf 1_L(f_1 x)\oplus\mathbf 1_L(f_2 x)$ is non-constant over $L$ whenever $f_1 x \neq f_2 x$, $A = 1$ forces $f_1 = f_2$; the label-flipping extension replaces $A$ by $A' = \Pr_{L,x}[\mathbf 1_L(f_1 x)\oplus\mathbf 1_L(f_2 x) = h_1(x)\oplus h_2(x)]$, and $A'=1$ forces $f_1=f_2$ and $h_1=h_2$. For the symplectic-orbit family against fresh \emph{rerandomized-secret} samples, $\mathbf 1_{S_iL}(S_ix)=\mathbf 1_L(x)$ makes the two membership values coincide identically, collapsing the overlap to the duplicate value $(p^2+(1-p)^2)/4^n$. All identities are confirmed by exact enumeration over symplectic, affine, and random bijections at $n=2$ (all $720$ symplectic maps, and exactly across the whole range $p\in(0,\tfrac12)$) and $n=3$; the identity has no small-$p$ degeneracy.
\end{proof}

\begin{proposition}[Label-preserving valid-output public maps are secret rerandomizations]
\label{prop:valid-output}
Let $n\ge 2$ and let $g$ be a \emph{label-preserving} public bijection --- $g(x,b)=(\varphi(x),b)$ for a bijection $\varphi$ of $\F_2^{2n}$ --- that sends LSN samples to LSN samples ($g_*(D_L)=D_{L'}$ for every Lagrangian $L$ and some $L'=L'(L)$). Then $g(x,b)=(Sx,b)$ for a fixed $S\in\mathrm{Sp}(2n,\F_2)$ (so $L'=S\cdot L$); no such map manufactures freshness beyond secret rerandomization. (A general valid-output bijection reduces to this label-preserving case --- its label marginal must equal $\mathrm{Bernoulli}(p)$ for every secret simultaneously, which forbids mixing the label into the point coordinate --- a marginal reduction we take as a structural sublemma rather than prove in full.)
\end{proposition}

\begin{proof}
By hypothesis $g(x,b)=(\varphi(x),b)$ is label-preserving, and the valid-output condition forces $\varphi$ to carry every Lagrangian to a Lagrangian. (The $n=2$ enumeration below is over the $16!$ point-space permutations $\varphi$, not the full $32!$ bijections of $\F_2^{2n}\times\F_2$; the reduction of a general bijection to this label-preserving form is the marginal sublemma noted in the statement.) Such a $\varphi$ is a collineation of the symplectic polar space $W(2n{-}1,2)$ (rank $n\ge 2$); by the fundamental theorem of polar geometry its automorphism group over $\F_2$ --- where no field automorphism intervenes --- is exactly $\mathrm{Sp}(2n,\F_2)$ \cite{Tits74,BCN89}, so $\varphi = S \in \mathrm{Sp}(2n,\F_2)$. At $n=2$ this needs no linearity assumption: among all $16!$ permutations of $\F_2^4$, exactly the $720=|\mathrm{Sp}(4,\F_2)|$ symplectic maps send every Lagrangian to a Lagrangian (exhaustive check).
\end{proof}

\begin{enumerate}
\item \textbf{Statistical dimension under $S_A=0$ for sympLPN \cite{LPQR26}.}\label{open:statdim} The full SQ proof is blocked by the deterministic quadratic symplectic relation $S_A=0$ (the $\Omega$-Gram of columns) that characterises the sympLPN variant (\Cref{def:symplpn}). Prove that conditioning on $S_A=0$ does not reduce the statistical dimension below $2^{\Omega(n)}$. \Cref{sec:moments} resolves the one-pair quadrant-count moments $m_j$ exactly for every $1 \le j \le 2n$, establishing that no statistical test of constant bundle size can exploit the conditioning; what remains open is whether these exact moment computations extend to the full growing-dimensional SQ/statistical-dimension regime --- many secrets and general tests rather than this one-pair statistic.

\item \textbf{LPN(low-noise) $\to$ LSN.} In the cryptographically relevant regime $p=O(1/\sqrt{n})$, the reduction of \cite{KLP+25} is vacuous. Does a non-vacuous reduction exist, or can its non-existence be proved?

\item \textbf{LWE $\to$ LSN.} A reduction from LWE (or worst-case lattice problems) to LSN would provide the gold-standard foundation currently missing.

\item \textbf{Quantum lower bound.} Prove \Cref{conj:quantum}: any quantum algorithm solving LSN requires $\Omega(2^n)$ queries.

\item \textbf{Pencil-extremality conjecture (sharp constant).} \label{open:pencil-sharp} \Cref{thm:main-sq-cond} is now \emph{unconditional}, via the big-cell theorem together with the augmentation/averaging transfer (absolute constant $24$). What remains open is the \emph{sharp} \Cref{conj:pencil}: every subset of size at least $|\Lagr(2n,\F_2)|/2^{2n-c}$ has average correlation $\leq 5\,C_n$. The $k\le 2$ pencils are the scale-relevant witnesses (ratio tending to $4$), so $5$ is a slack universal target rather than a sharp value attained by the pencils; establishing it would sharpen the VSTAT parameter of \Cref{thm:main-sq-cond} from $1/(72\rho_{\mathrm{avg}}^{\mathrm{SQ}})$ to $1/(15\,\rho_{\mathrm{avg}}^{\mathrm{SQ}})$.

\item \textbf{Sample freshness for rerandomized LSN.} Exact secret rerandomization is possible (apply a public $S \in \mathrm{Sp}(2n)$ to each sample's public part, mapping $L \mapsto S\cdot L$ while preserving labels), but it cannot manufacture \emph{fresh} samples: \Cref{prop:freshness} settles the symplectic-orbit and all label-preserving (and label-flipping) split families negatively at every $n$, with statistical distance $1 - (p^2+(1-p)^2)/4^n \to 1$ and equality only at the literal duplicate. Two refinements narrow the family. First, by \Cref{prop:valid-output}, every \emph{label-preserving} valid-output public bijection is a pure secret rerandomization $(x,b)\mapsto(Sx,b)$, $S\in\mathrm{Sp}(2n,\F_2)$, and so cannot manufacture freshness (a general valid-output bijection reduces to the label-preserving case by the marginal sublemma of \Cref{prop:valid-output}). Second, a direct computation shows that even \emph{$b$-dependent} label-preserving bijections $g_i(x,b)=(\varphi_{i,b}(x),b)$ obey $\mathrm{SD} = 1 - (p^2+(1-p)^2)/4^n + \tfrac{(1-2p)^2}{2\cdot 4^n}(2 - A_0 - A_1) \ge 1 - (p^2+(1-p)^2)/4^n$, with equality iff $g_0=g_1$. What remains \emph{open} is whether any transformation with non-product output structure (necessarily not valid-output) yields fresh samples; a positive answer would give tight multi-user security bounds for LSN-based constructions, while a fully general negative answer would explain why hybrid arguments lose a factor of the number of users.

\item \textbf{Membership-LSN $\leftrightarrow$ stabilizer-decoding LSN bridge.} Relate our membership formulation (\Cref{def:lsn}; secret Lagrangian, no public matrix) to the LSN of \cite{KLP+25,LPQR26} (public stabilizer matrix $[A \mid B]$ with a Lagrangian $A$-block, noisy-codeword samples, secret logical string). For the \emph{sympLPN} formulation this cumulative goal is already met: \Cref{prop:bernoulli-inherits} attaches the external constant-rate-LPN hardness to the very Bernoulli sympLPN object that \Cref{cor:symplpn-sq} and \Cref{thm:symplpn-depol} bound, so there the external hardness and our SQ lower bounds describe one and the same object. The open bridge concerns specifically the \emph{membership} formulation: a reduction in either direction would let the external hardness results (constant-rate LPN $\le$ their LSN via \cite[Thm.~6.4 and Lemma~6.5]{KLP+25}; their search/decision and single-sample equivalences --- \cite[Thm.~1.9 and Cor.~5.6]{KLP+25} --- carry this hardness to the \emph{single-sample} Search variant) and our SQ lower bounds support a \emph{single} object there as well; absent such a bridge, the membership and stabilizer-decoding hardness cases remain parallel rather than cumulative. A reduction THEIR $\le$ OUR must embed the $k$-bit logical secret into our $\Theta(n^2)$-bit Lagrangian secret without the junk register that \cite{KLP+25} (Lemma~6.5) exploits to bypass the analogous dimension mismatch: from the single-sample hardness carrier it meets an information-budget obstruction shared with \Cref{open:marginal-adaptive} (one noisy codeword carries $O(n)$ secret-correlated bits, while membership recovery needs $\Omega(n^2)$ labels); the multi-sample variant instead meets a frame-alignment obstruction (the secret arrives in a fresh public frame per sample, whereas ours is one fixed Lagrangian). The reverse direction is blocked by code visibility (a public $[A \mid B]$ would encode our secret $L$). None of these is an impossibility result; each blocks only the natural maps.

\item \textbf{Absence of a symplectic worst-to-average self-reduction.} Like LPN, and unlike LWE, LSN has no known worst-case-to-average-case reduction (\Cref{subsec:two-forms} treats LSN as a primitive assumption); we record a structural obstruction to the natural symplectic route. The group $\mathrm{Sp}(2n,\F_2)$ acts transitively on Lagrangians, so for $g\sim\mathrm{Unif}(\mathrm{Sp}(2n,\F_2))$ the map $L\mapsto gL$ \emph{perfectly} randomizes the secret subspace; the obstruction is purely metric. Applying $g$ carries a sample $y=Ax+e$ to $gy=(gA)x+ge$ (with $gA$ spanning the uniform Lagrangian $gL$), but $ge$ is no longer $\mathrm{Bernoulli}(p)^{2n}$. In the Walsh domain $\widehat{g_*B_p}(u)=(1-2p)^{\mathrm{wt}(g^{\top}u)}$, so correcting $g_*B_p$ back to $\mathrm{Bernoulli}(p')^{2n}$ by convolving with independent fresh noise $P'$ would require $\widehat{P'}(u)=(1-2p')^{\mathrm{wt}(u)}/(1-2p)^{\mathrm{wt}(g^{\top}u)}\le 1$ for every $u$, i.e.\ $1-2p'\le(1-2p)^{W(g)}$ with $W(g):=\max_{u\neq 0}\mathrm{wt}(g^{\top}u)/\mathrm{wt}(u)$. For a generic $g\in\mathrm{Sp}(2n,\F_2)$ one has $W(g)=\Theta(n)$ --- immediate, since $g^{\top}$ maps a weight-$1$ vector to a column of $g^{\top}$, a generic non-zero vector of weight $\Theta(n)$, while always $W(g)\le 2n$ --- driving the corrected rate $p'\to 1/2$ and destroying the secret. The weight-preserving subgroup ($W=1$, the symplectic monomial maps) is \emph{not} transitive on Lagrangians: exhaustive enumeration of all $720$ elements of $\mathrm{Sp}(4,\F_2)$ shows that single-application reachability with $W\le 2$ covers only $10$ of the $15$ Lagrangians, while the remaining $5$ require $W=3$ (corrected rate $0.4375$ at $p=1/4$). Leaving the noise uncorrected is circular: a public $g$ lets the distinguisher apply $g^{-1}$ and recover the worst-case instance verbatim. Whether an intermediate subgroup $H$ with $\mathrm{(monomial)}<H\le\mathrm{Sp}(2n,\F_2)$ that is transitive on Lagrangians with bounded $W$, or a non-symplectic randomization, escapes the barrier is open; we conjecture none does. The root cause is a Hamming--symplectic incompatibility: Hamming weight is not $\mathrm{Sp}$-invariant. Since $\mathrm{Sp}(2n,\F_2)$ is transitive on $\F_2^{2n}\setminus\{0\}$, \emph{every} $\mathrm{Sp}$-invariant law is constant on the non-zero vectors --- the weight-blind one-parameter family $q\,\delta_0+(1-q)\,\mathcal U(\F_2^{2n}\setminus\{0\})$ --- and the symplectic twirl of $\mathrm{Bernoulli}(p)^{2n}$ lands at $q=(3/4)^{2n}$, carrying no weight signal; no weight-graded $\mathrm{Sp}$-invariant noise exists, so there is no discrete analogue of the Gaussian self-duality underlying Regev's lattice reduction. This is consistent with the isotropic-constraint barrier that \cite{LPQR26} report for adapting LPN-style security proofs, and positions LSN with LPN among assumptions trusted by sustained cryptanalysis rather than by a self-reduction backbone. The obstruction is stated for the natural correction by an \emph{independent} fresh-noise convolution (and rejection sampling onto $\mathrm{Bernoulli}(p')$ has acceptance $2^{-\Theta(n)}$); whether a \emph{correlated} or non-linear correction, or an intermediate transitive subgroup $H$ with bounded $W$, escapes it is open (\emph{evidence}: the $n=2$ enumeration and the elementary $W(g)=\Theta(n)$; \emph{conjecture}: that no usable intermediate~$H$ or correction exists).

\item \textbf{Randomized marginal-adaptive linear reductions (sharpened).} \label{open:marginal-adaptive} In the isotropic-to-LPN reduction model, the distinguisher receives the public matrix $C=BA$ and the noisy output $y=Cx+e$. The deterministic case is closed by \Cref{thm:deterministic-marginal-adaptive}, so the remaining question concerns reductions that use fresh private randomness $R$. The operative obstruction is the statistical-distance bound \Cref{lem:m2}: the output $(C,y)$ must stay far from a \emph{usable} $\mathrm{LPN}_{p'}$. We flag a natural but \emph{misdirected} mutual-information route. Exact computation of $I(x;y\mid C)$ for the marginal-uniform output (closed form over the ordered-basis isotropic ensemble) is sublinear and \emph{saturating} in the number of samples while staying well below $H(x)$; the finite-ratio values $0.102,0.077,0.053$ at $m/n=2$ ($n=2,3,4$) reflect this per-sample saturation, \emph{not} an $o(n)$ decay in $n$. As $m$ grows the leakage does not vanish: a mechanism-level computation at $n=2,3,4$ shows the syndrome rank is capped at $n$ (the identity $\operatorname{rank}(HB)=\operatorname{rank}(B)-\operatorname{rank}(BA)\le n$), so the complementary $n$ message coordinates retain bias $\approx 1-2p$ even conditioned on the syndrome, giving $I(x;y\mid C)=\Theta(n)$ (\emph{evidence}; the exact constant is open). Hence the target ``$I(x;y\mid C)=o(n)$'' is almost certainly \emph{unattainable}---and it is also unnecessary: the $\Theta(n)$ leakage is symplectic-structured, not LPN-structured, so $H(x\mid y,C)\ge(2H_2(p)-1)n$ (\Cref{lem:secret-floor}) rules out high-probability \emph{full recovery} of $x$ from the transcript --- but this is only a Shannon-entropy floor, not a proof of statistical distance from LPN. The no-go therefore rests on \Cref{lem:m2} (statistical distance), not on an information bound; (for context, prior Fisher-information and total-variation approaches bounded $I(x;y)$, but since $x$ is independent of $C$ one has $I(x;y)\le I(x;y\mid C)$, so they did not even bound the conditional quantity).

The status map for the linear-reduction landscape is: fixed-$B$ \textbf{ruled out}, public-$B$ \textbf{ruled out}, conditionally-uniform adaptive \textbf{ruled out}, deterministic marginal-adaptive \textbf{ruled out}, randomized marginal-adaptive \textbf{open}.  \Cref{lem:m1} shows that any marginal-adaptive $B=g(A)$ with uniform $BA$ must have, \emph{in expectation}, $m-16n-o(m)$ rows of weight $>0.19n$ (hence effective noise bias $\le 2^{-0.19n}$), but \Cref{lem:m2} is not yet established because the output noise $e'=Be$ is $m$ linear images of a fixed $2n$-bit symplectic-LPN noise vector and therefore lives in a subspace of dimension at most $2n$.  The decisive question: \emph{Does the $\le 2n$-dimensional noise structure of $Be$ survive a high-entropy secret $B$, or can adaptive $B$ fake $m$ independent $\operatorname{Bernoulli}(p')$ from $2n$ noise bits undetectably in $C$?}  A proof either way would close the marginal-adaptive corner. Two concrete strategies are provably non-reducing. The \emph{uniform-$B$-per-$A$} instance (each $B$ drawn uniformly given $A$) has output statistical distance $\to 1$ from matched-rate LPN at fixed $n$: writing the output mixture $q\,P_{\mathrm{graph}} + (1-q)\,P_{\mathrm{full}}$, a Hellinger tensorization gives $\mathrm{SD} \ge 1 - C(n)\,\rho(n)^m$ with $\rho(n) = 1 - 2^{-n}(1 - \sqrt{(1-p_{\mathrm{eff}})/2} - \sqrt{p_{\mathrm{eff}}/2}) < 1$, and the optimal scalar distinguisher is the minimum syndrome weight $\min_w \mathrm{wt}(y + Cw)$ (which captures over $99\%$ of the exact distance numerically, while the membership test saturates at $q$). A row-\emph{correlated} marginal-uniform family (rows equal with probability $\lambda$) moves the distance the wrong way --- correlation collapses the support dimension of $Be$, making the output \emph{more} distinguishable, not less. Both observations point to the same residual: any reducing $B$ would have to keep $Be$'s joint law product-like while staying marginal-uniform, which neither strategy achieves. Moreover, the $W=0$ spike is a \emph{$B$-agnostic} provable distinguisher: for \emph{every} marginal-uniform $B$, whenever $e \in \operatorname{Col}(A)$ --- i.e.\ $e = Aw$, an event of probability $q_{\mathrm{graph}}(n)$ that depends only on the noise prior, not on $B$ or its private randomness --- one has $Be = BAw = Cw$, so $y = C(x+w)$ is a codeword and $\min_w \mathrm{wt}(y+Cw) = 0$; since real LPN has $\Pr[W=0] \le 2^n(1-p')^m$, the output leaks at rate $\ge q_{\mathrm{graph}}(n) - \mathrm{negl}$ regardless of $B$. This is a constant gap at each fixed $n$ but vanishes asymptotically ($q_{\mathrm{graph}}(n) \to 0$), so it rules out reductions at any fixed dimension without by itself giving an asymptotic bound. A sharper, \emph{non-vanishing} distinguisher exists in the \emph{shared-sample} model. If a reduction must answer $k\ge 2$ samples with a common public matrix $C$ (equivalently a fixed $B$), stack the outputs $Y = [\,y^{(1)} \mid \dots \mid y^{(k)}\,]$ and apply a parity-check $H$ of $C$ ($HC=0$): then $HY = H(CX + BE) = HBE$. Since $C=BA$ gives $\operatorname{Col}(C)=\operatorname{Col}(BA)\subseteq\operatorname{Col}(B)$ and $\ker H=\operatorname{Col}(C)$, the noise image obeys the sharp identity $\operatorname{rank}(HB)=\operatorname{rank}(B)-\operatorname{rank}(BA)$, so
\[
\operatorname{rank}(HY) \le \operatorname{rank}(HB) \le n \quad \text{deterministically, for every } B
\]
(the naive confinement bound $2n$ is loose by a factor of two, since $H$ already annihilates the $n$-dimensional $\operatorname{Col}(C)\subseteq\operatorname{Col}(B)$). For full-rank $B$ the identity gives $\operatorname{rank}(HB)=n$ exactly; the bound can collapse to $0$ only when $\operatorname{rank}(B)\le n$, which forces $\operatorname{Col}(B)=\operatorname{Col}(C)$ and hence $y\in\operatorname{Col}(C)$ for every output --- i.e.\ the already-provable $W=0$ spike --- so neither branch escapes detection. Real LPN, whose noise is full-entropy (and whose rate $\to 1/2$ under \Cref{lem:m1}), has $\operatorname{rank}(HY) \to \min(m-n, k)$; for $m \ge 2n$ and $k > n$ this exceeds $n$ with probability $1 - o(1)$, a detection advantage that does \emph{not} vanish in $n$. Thus the shared-sample (fixed-$B$) case is closed asymptotically by this \emph{multi-sample} ($k\ge 2$) statistic --- but for the single block of \Cref{lem:m2} ($k=1$, one noise vector $e$) the stack has a single column and $\operatorname{rank}(HY)\in\{0,1\}$ is trivial, leaving that object to its low-dimensional-support signal (statistical distance $\to 1$ at fixed $n$ by the Hellinger tensorization above, now corroborated by an \emph{exact} evaluation of the min-syndrome-weight scalar statistic at \emph{solvable} sample sizes $m\ge n/(1-H_2(p'))$, for $n=2$, via a coset-leader-weight transfer matrix (\emph{evidence}, exact at $n=2$): its advantage is bounded away from $0$ and rising in $m$ there, the regime the earlier $m\le 4$ enumerations sat below; asymptotic-in-$n$ rate open). The genuine residual is the \emph{fresh-$B$} model --- the \emph{multi-sample} regime, in which the reduction is handed independent sympLPN samples (the multi-sample variant $\mathrm{LSN}^m$ of \cite{KLP+25}, a shared secret across fresh public bases supplied by the oracle, \emph{not} a public rerandomization of a single instance, which the sample-freshness item above bars), so that each sample carries an independent basis $A^{(i)}$ and matrix $B^{(i)}$ and the noises $B^{(i)}e^{(i)}$ no longer share one $\le 2n$-dimensional space and the stacked-rank invariant does not apply; there, exact computation (uniform-$B$, row-correlated, column-pair, and $k$-sample rank families, $n=2$, $m \le 80$) shows every tested distinguisher leaks only at a rate $\to 0$ in $n$. Whether the fresh-$B$ marginal-adaptive reduction is asymptotically distinguishable remains the open core.
\end{enumerate}

\subsection{Limitations of the Present Work}
\label{subsec:limitations}

We state explicitly the gaps a reader should weigh against the results above.

\begin{enumerate}
\item \textbf{No worst-case foundation.} LSN hardness is a primitive assumption: there is no reduction from a worst-case problem in the style of Regev's reduction for LWE, and no non-vacuous reduction from low-noise LPN or from LWE is known (Open Problems~2--3).

\item \textbf{SQ bounds are evidence, not proof.} Exponential SQ lower bounds also hold for LPN and LWE; they bound a restricted (if natural and historically predictive) class of algorithms and do not distinguish hardness families. The unconditional bound of \Cref{thm:main-sq-uncond} moreover protects only a promise subfamily.

\item \textbf{One result is conditional.} The marginal-adaptive barrier (\Cref{thm:marginal-adaptive}) depends on \Cref{lem:m2}, which remains unproven; the dependency is stated where it occurs. (The $2^{2n-O(1)}$ SQ bound is now unconditional via the big-cell theorem and averaging transfer, with absolute constant $24$; only its sharp constant rests on the pencil-extremality conjecture \Cref{conj:pencil}.)

\item \textbf{Moment results close every fixed-order multi-secret composition; the $k=\Theta(n)$ level is open.} The closed forms of \Cref{sec:moments} are exact for every order $j$, \Cref{prop:vmax} controls the variance multiplier at the maximal block, \Cref{prop:tdist} recovers the full exact law of the quadrant count (with the proven TV rate $2^{-(n+1)}$), and \Cref{thm:kfold-gf} closes the joint composition of every fixed-size secret tuple ($k \le n$; \Cref{thm:joint-gf,thm:triple-gf} are $k=2,3$). What remains open is the $k=\Theta(n)$ regime and the SQ statements built on these joint laws.

\item \textbf{Quantum hardness is conjectural.} We analyze the natural quantum attack families (\Cref{sec:quantum}) and find no known efficient attack at constant noise, but prove no quantum query lower bound; \Cref{conj:quantum} has the same epistemic status as the corresponding conjectures for LPN and LWE.
\end{enumerate}

\appendix

\section{Full Proof of Pairwise Correlation}
\label{app:corr}

Here we give a self-contained derivation of \Cref{lem:exact-corr}. Let $D_0$ be the distribution on $\F_2^{2n} \times \F_2$ in which $a$ is uniform and $b \sim \operatorname{Bernoulli}(p)$ independently. For a fixed Lagrangian $L$, the LSN distribution $D_L$ has density
\[
D_L(a,b) = \frac{1}{2^{2n}} \Bigl( (1-p)\mathbf{1}_{b=\mathbf{1}_L(a)} + p\mathbf{1}_{b \neq \mathbf{1}_L(a)} \Bigr).
\]
The likelihood ratio is
\[
\frac{D_L(a,b)}{D_0(a,b)} = 1 + \mathbf{1}_L(a) \cdot \beta \cdot \Bigl(\mathbf{1}_{b=1} - \mathbf{1}_{b=0}\frac{p}{1-p}\Bigr),
\]
where $\beta = (1-2p)/p$. Therefore
\[
\ell_L(a,b) = \mathbf{1}_L(a) \cdot \beta \cdot \Bigl(\mathbf{1}_{b=1} - \mathbf{1}_{b=0}\frac{p}{1-p}\Bigr).
\]
For two secrets $L,L'$ with $j = \dim(L \cap L')$,
\begin{align*}
\langle D_L, D_{L'} \rangle
&= \E_{D_0}[\ell_L \ell_{L'}] \\
&= \E_a[\mathbf{1}_L(a)\mathbf{1}_{L'}(a)] \cdot
   \beta^2 \E_b\Bigl[\Bigl(\mathbf{1}_{b=1} - \mathbf{1}_{b=0}\frac{p}{1-p}\Bigr)^2\Bigr] \\
&= 2^{j-2n} \cdot \beta^2 \Bigl( p + (1-p)\frac{p^2}{(1-p)^2} \Bigr) \\
&= 2^{j-2n} \cdot \frac{(1-2p)^2}{p(1-p)}.
\end{align*}
Setting $p=1/4$ gives the coefficient $(1/2)^2 / (1/4 \cdot 3/4) = (1/4)/(3/16) = 4/3$ exactly.

\section{Q-Binomial Distance Distribution}
\label{app:qbin}

The number of $k$-dimensional subspaces of $\F_2^n$ is the Gaussian binomial coefficient
\[
{n \brack k}_2 = \prod_{i=0}^{k-1} \frac{2^{n-i}-1}{2^{k-i}-1}.
\]
To count Lagrangians $L'$ with $\dim(L\cap L')=k$ \emph{exactly}, fix the core $W=L\cap L'$, a $k$-dimensional (isotropic) subspace of $L$. Then $L'\supseteq W$ and, crucially, $L'/W$ must be \emph{transverse} to $L/W$ in the rank-$(n-k)$ symplectic quotient $W^\perp/W\cong\F_2^{2(n-k)}$ (otherwise the intersection would exceed $k$); the transverse Lagrangians of $L/W$ form its big cell, of size $2^{(n-k)(n-k+1)/2}$ --- not the larger $\prod_{i=1}^{n-k}(2^i+1)$ of \emph{all} Lagrangians containing $W$, which would count intersections of dimension $\ge k$. Since $L$ has ${n\brack k}_2$ subspaces $W$ of dimension $k$ and each $L'$ determines its core $W$ uniquely (no overcounting), the exact probability is
\[
\Pr[j=k] = \frac{{n \brack k}_2 \cdot 2^{(n-k)(n-k+1)/2}}{\prod_{i=1}^n (2^i+1)}.
\]
The moment generating function $C_n = \E[2^j]$ converges rapidly to $2$ as $n \to \infty$.

\section{Full Proof of the Affine-Coset Bias Bound}
\label{app:krawtchouk}

Let $N \subset \F_2^{2n}$ be a uniform random Lagrangian.  For $\lambda \in (0,1)$ define the weight enumerator $W_N(\lambda) := \sum_{v \in N} \lambda^{|v|}$.  In \Cref{lem:affine-coset-bias-whp} we need $\lambda=1/2$.

\paragraph{Exact moments.}  Write $p := \Pr[v \in N \setminus \{0\}] = 1/(2^n+1)$ and $q := \Pr[v,v' \in N \setminus \{0\} \mid v \neq v',\; \Omega(v,v')=0] = 1/((2^{n-1}+1)(2^n+1))$.  Let $X_v := \mathbf{1}_{v \in N}$ and $W := W_N(1/2)-1 = \sum_{v \neq 0} X_v\,2^{-|v|}$.  Then
\[
\E[W] = \frac{(3/2)^{2n}-1}{2^n+1}, \qquad
\Var[W] = p(1-p)\,D + q\,S_0 - p^2 T,
\]
where $D = (5/4)^{2n}-1$, $T = ((3/2)^{2n}-1)^2-D$, and $S_0 = \tfrac12(T+C)$ with
\[
C = \sum_{\substack{v\neq v'\\ v,v'\neq 0}} (-1)^{\Omega(v,v')}2^{-|v|-|v'|}
  = (7/4)^{2n} - 2(3/2)^{2n} - (5/4)^{2n} + 2.
\]

\paragraph{Block factorisation of $C$.}  Decompose $v=(a,b)$, $v'=(x,y)$ with $a,b,x,y \in \F_2^n$.  The symplectic form factorises coordinate-wise: $\Omega(v,v') = \sum_{i=1}^n (a_i y_i + b_i x_i)$.  Hence
\[
\sum_{v,v' \in \F_2^{2n}} (-1)^{\Omega(v,v')} 2^{-|v|-|v'|}
= \prod_{i=1}^n \underbrace{\sum_{a,b,x,y\in\{0,1\}} (-1)^{ay+bx} 2^{-(a+b+x+y)}}_{=\,49/16 = (7/4)^2}
= (7/4)^{2n}.
\]
Removing the terms with $v=0$ or $v'=0$ (each contributes $(3/2)^{2n}$) and adding back the double-counted $(0,0)$ term yields $C$ as above.

\paragraph{Asymptotic decay.}  Keeping only dominant exponentials in each symbol,
\begin{align*}
p(1-p)D &= (25/32)^n + O((25/64)^n), \\
qS_0 &= (81/64)^n + (49/64)^n + O((9/16)^n), \\
p^2T &= (81/64)^n + O((9/16)^n).
\end{align*}
The $(81/64)^n$ terms cancel exactly because $\tfrac{q}{2}-p^2 = -1/(2(2^{n-1}+1)(2^n+1)^2)$.  Therefore
\[
\Var[W] = (25/32)^n + (49/64)^n + O((9/16)^n) + O((25/64)^n).
\]
Since $\E[W] = (9/8)^n(1+o(1))$, we obtain $\Var[W]/\E[W]^2 = O((50/81)^n)$.

\paragraph{Chebyshev concentration.}  Choose $\varepsilon_n = (50/81)^{n/4}$.  Then
\[
\Pr\bigl[|W-\E[W]| \ge \varepsilon_n \E[W]\bigr]
\le \frac{\Var[W]}{\varepsilon_n^2 \E[W]^2}
= O((50/81)^{n/2}) = 2^{-\Omega(n)}.
\]
On this high-probability event,
\[
\frac{W_N(1/2)-1}{2^n} = \frac{(9/8)^n}{2^n}(1+o(1)) = (9/16)^n(1+o(1)),
\]
which yields the bound of \Cref{lem:affine-coset-bias-whp}.

\section{Reduction Barrier Summary}
\label{app:barrier}

This appendix summarizes the SQ/feature-map barriers; for the full linear-reduction landscape (transport, entropy-continuity, reachability, and the open marginal-adaptive cell) see \Cref{tab:barriers} and \Cref{sec:barriers}.

\paragraph{Linear.} For any linear classifier $h(a) = \langle w,a \rangle + t$, the advantage
\[
\Adv(h) = \bigl| \Pr_{D_L}[h(a)=1] - \Pr_{D_0}[h(a)=1] \bigr|
\]
is zero because $\Pr_{D_L}[h(a)=1] = \Pr_{D_0}[h(a)=1]$.

\paragraph{Polynomial degree $<n$.} The exact polynomial representation of $\mathbf{1}_L$ has degree exactly $n$ and $2^n$ monomials. After an invertible linear change of variables $\mathbf{1}_L$ is the monomial $z_1\cdots z_n$ of degree exactly $n$, whose distance from $\mathrm{RM}(n-1,2n)$ is $2^{-n}$; hence every polynomial of degree $d < n$ differs from $\mathbf{1}_L$ on at least a $2^{-n}$ fraction of inputs, so its correlation with $\mathbf{1}_L$ is at most $2^{-n}$.

\paragraph{Adaptive linear SQ.} For real-linear queries $q(a,b) = \langle w,a \rangle + cb$ we have $\E_{D_L}[q] = c(p + 2^{-n}(1-2p))$, independent of $L$. For F$_2$-linear queries the maximum advantage is $(1-2p)2^{-n}$ (\Cref{thm:linear-sq}). In either case the advantage is exponentially small in $n$, so adaptive linear SQ algorithms require exponentially many samples.

\section{Extension to $\F_q$}
\label{app:fq}

All results in the main text are stated over $\F_2$. Here we sketch the extension to arbitrary finite fields $\F_q$, which widens the parameter space and allows a trade-off between block length $n$ and field size $q$.

\subsection{Symplectic Spaces over $\F_q$}
\label{subsec:fq-symplectic}

Let $q$ be a prime power and $V = \F_q^{2n}$ equipped with a non-degenerate alternating bilinear form $\Omega(x,y) = x^T J y$ where $J = \begin{psmallmatrix} 0 & I_n \\ -I_n & 0 \end{psmallmatrix}$. A subspace $L \subset V$ is \emph{Lagrangian} if it is maximal isotropic; then $\dim L = n$ and $|\Lagr(2n,\F_q)| = \prod_{i=1}^n (q^i+1)$.

For fixed $L$ and $j = \dim(L \cap L')$, the number of Lagrangians $L'$ with intersection dimension $j$ is
\[
N_j(n,q) = {n \brack j}_q \cdot q^{(n-j)(n-j+1)/2},
\]
where ${n \brack j}_q$ is the Gaussian binomial coefficient. Summing over $j$ recovers $|\Lagr(2n,\F_q)|$.

\subsection{Exact Correlation over $\F_q$}
\label{subsec:fq-corr}

The label remains binary ($b \in \F_2$) even though the ambient space is $\F_q^{2n}$, so the SQ machinery of Feldman \etal\ applies without modification. The likelihood ratio calculation is identical with $2$ replaced by $q$:

\begin{lemma}[$\F_q$ pairwise correlation]
\label{lem:fq-corr}
For $L,L' \in \Lagr(2n,\F_q)$ with $j = \dim(L \cap L')$,
\[
\langle D_L, D_{L'} \rangle = \frac{(1-2p)^2}{p(1-p)} \cdot \frac{q^j}{q^{2n}}.
\]
\end{lemma}

\begin{lemma}[$\F_q$ average correlation]
\label{lem:fq-avg}
Let $C_{n,q} = \E[q^j]$ computed via the $q$-binomial distribution. Then
\[
\rho_{\mathrm{avg}} = \frac{(1-2p)^2}{p(1-p)} \cdot C_{n,q} \cdot q^{-2n}.
\]
Numerically, $C_{n,q} \to 2$ as $n \to \infty$ for every $q$ (the limit is independent of the field size).
\end{lemma}

\subsection{sympLPN over $\F_q$: the noise-universal $M_2$ law}
\label{subsec:fq-symplpn}

The general per-pair-noise SQ law of \Cref{thm:symplpn-depol} is field-agnostic. The $\F_q$ sympLPN problem has public isotropic $A \in \F_q^{2n\times n}$, secret $x\in\F_q^n$, and $\F_q$-valued sample $(A,\,Ax+e)$, with $e$ drawn per qudit i.i.d.\ from a law $\mu$ on $\F_q^2$ (a $q$-ary Pauli channel); the SQ framework of \cite{FGR+17} applies to the $\F_q$-valued output unchanged.

\begin{proposition}[$\F_q$ sympLPN: noise-universal correlations]
\label{prop:fq-symplpn-m2}
Let $\hat\mu(\sigma)=\sum_{e\in\F_q^2}\mu(e)\,\omega^{\langle\sigma,e\rangle}$ with $\omega=e^{2\pi i/q}$ (so $\hat\mu(0)=1$), and $M_2 = \sum_{\sigma\in\F_q^2}|\hat\mu(\sigma)|^2$. Then
\[
\langle D_x,D_x\rangle = M_2^{\,n}-1, \qquad \langle D_x,D_{x'}\rangle = -\frac{M_2^{\,n}-1}{q^{2n}-1}\quad(x\neq x'),
\]
and any SQ distinguisher of $\{D_x\}$ from the null needs $q^{\,c_p n-O(1)}$ queries at constant VSTAT strength with $c_p = 1-\log_q M_2$. The $q$-ary depolarizing channel ($1-p$ on the identity, $p/(q^2-1)$ on each of the $q^2-1$ nonzero Paulis) has $\hat\mu(\sigma\neq 0)=\tau_{d,q}:=1-\tfrac{pq^2}{q^2-1}$, hence $M_2 = 1+(q^2-1)\,\tau_{d,q}^2$, recovering $1+3\tau_d^2$ at $q=2$.
\end{proposition}

\begin{proof}
Identical to \Cref{thm:symplpn-depol} with $2$ replaced by $q$. Parseval on each $\F_q^2$ factor gives $\sum_e \mathrm{mass}(e)^2=\prod_i\bigl(\tfrac1{q^2}\sum_\sigma|\hat\mu(\sigma)|^2\bigr)=(M_2/q^2)^n$, so the diagonal is $q^{2n}(M_2/q^2)^n-1=M_2^n-1$, independent of $A$. Averaging the off-diagonal over $c=A(x{-}x')$, uniform on $\F_q^{2n}\setminus\{0\}$ by $\mathrm{Sp}(2n,\F_q)$-transitivity, factorizes per qudit into $\sum_\sigma|\hat\mu(\sigma)|^2\,\omega^{\langle\sigma,c_{\mathrm{qudit}}\rangle}$, whose uniform-$c$ average is $|\hat\mu(0)|^2=1$, collapsing the nonzero-$c$ average to $-(M_2^n-1)/(q^{2n}-1)$. Verified exactly at $q=2,3,5$ ($n=2$).
\end{proof}

\paragraph{Structural inertness over $\F_q$ (evidence).} The self-dual code view of \Cref{subsec:symplpn-sq} carries over: $\F_q$ sympLPN is syndrome decoding of a random symplectically self-dual $\F_q$-code. \emph{Evidence} indicates this code is no weaker than a generic random rate-$\tfrac12$ $\F_q$-code: in the tested small cases its minimum distance was not below the generic one --- exact enumeration at $q=3$ ($n\le 3$) and $q=5$ ($n=2$), and sampling at $q=3,\,n=4$ and $q=5,\,n=3$ (the Lagrangian mean exceeds the generic by a factor $\approx 1.04$--$1.07$, never smaller). This suggests, but does not prove, that information-set decoding gains no structural advantage over $\F_q$, as over $\F_2$.
\label{subsec:fq-params}

The average-case SQ lower bound \emph{would} become $q_{\mathrm{queries}} \geq q^{2n-O(1)}$, a leading SQ exponent of $2n \log_2(q)$; but over $\F_q$ this is \emph{conjectural} (\textbf{CONJECTURE}), since it requires the $\F_q$ analogue of the big-cell Delsarte certificate and the averaging transfer behind \Cref{thm:main-sq-cond}, which we do not establish here --- only the exact $\F_q$ correlations (\Cref{lem:fq-corr,lem:fq-avg}) and the noise-universal sympLPN bound (\Cref{prop:fq-symplpn-m2}) carry over unconditionally. As a \emph{heuristic asymptotic guide} (not a concrete-security estimate), \Cref{tab:fq-params} lists, on that conjectural basis, the block lengths $n$ at which the leading exponent reaches each target level.

\begin{table}[ht]
\centering
\caption{$\F_q$ block lengths $n$ at which the leading SQ exponent $2n\log_2 q$ reaches each level, for $p=1/4$ (heuristic asymptotic guide, not concrete bit-security).}
\label{tab:fq-params}
\begin{tabular}{@{}ccccc@{}}
\hline
Target exponent $\lambda$ & $\F_2$: $n$ & $\F_3$: $n$ & $\F_5$: $n$ & $\F_7$: $n$ \\
\hline
$80$  & 41  & 26 & 18 & 15 \\
$128$ & 65  & 41 & 28 & 24 \\
$192$ & 97  & 62 & 42 & 35 \\
$256$ & 129 & 82 & 56 & 46 \\
\hline
\end{tabular}
\end{table}

Larger $q$ reduces the required $n$ but increases the cost of each sample (field arithmetic in $\F_q$ versus $\F_2$). The optimal choice depends on the target platform.

\subsection{Open Questions}
\label{subsec:fq-open}

\begin{enumerate}
\item Extend to $q$-ary labels $y \in \F_q$ with noise over $\F_q$; this requires $q$-ary SQ lower bounds.
\item The proof holds for $q = p^k$ with $k > 1$ because symplectic theory is field-agnostic.
\item In characteristic $2$, alternating is equivalent to symmetric; in odd characteristic they differ. The intersection distribution formula remains valid in all characteristics.
\end{enumerate}

% Bibliography
\bibliographystyle{alpha}
\section*{Acknowledgements}
This work was conceived and directed by the author, who set the research agenda, posed and pursued the problems, made the final mathematical and cryptanalytic judgments, and accepted each result only after independent re-derivation and computational cross-checking. In support of this author-led investigation, large-language-model assistants were used for literature search, exploratory calculation, code prototyping, and drafting; the author reviewed and corrected their output and takes full responsibility for the accuracy, integrity, and final content of the work.

\paragraph{\textbf{Data and code availability.}} All supporting material --- Python enumeration and verification scripts, Rust implementations (cryptanalysis screens, decoder validation, and a reference scaffold), SageMath symbolic checks, exact-fraction data and test vectors, and the \LaTeX{} source --- is available at \url{https://github.com/Gattochoo/LSN-PQ} and archived at \href{https://doi.org/10.5281/zenodo.20805781}{\texttt{doi:10.5281/zenodo.20805781}}. In particular, the closed forms of \Cref{sec:moments} are reproducible from the enumeration and symbolic-verification scripts. This version of the paper corresponds to the tagged release \texttt{v3.3.1} (each tagged release carries its own Zenodo version DOI), pinning the exact source and scripts against later repository drift.

\paragraph{\textbf{Conflict of interest.}} The author declares no competing interests.

\section{The big-cell Delsarte certificate}\label{app:bigcell-cert}
This appendix records the explicit constant dual certificate behind \Cref{thm:bigcell}, so that the big-cell bound can be re-derived mechanically; the exact-rational verification scripts (\texttt{verify\_bigcell\_certificate.py} and \texttt{endtoend\_gate.py}, under \texttt{experiments/certificate/}) are part of the code archive of the data-and-code availability statement.

On the symmetric bilinear forms association scheme~\cite{Schmidt2018} (eigenspaces labelled by $(\mathrm{rank}\ r,\mathrm{type}\ \tau)$, with $P$-matrix $P[i,k]$), the size-constrained Delsarte program at density $\alpha\ge 2^{-2n+2}$ admits a \emph{constant} dual $F=\lambda+\mu\,\delta_0+\sum_j\eta_j E_j$ with $\lambda=8$ and
\[
\eta_i\equiv -5\cdot 2^{-2n}\ \text{on}\ S=\{(4,-)\}\cup\{(2a,\pm):a\ge3\}\cup\{(2b{+}1,0):b\ge2\},\qquad \eta_j=0\ \text{else},
\]
$\mu$ being fixed by the size constraint. Writing $g(k)=\sum_{i\in S}P[i,k]$ (which has the closed form $g(k)=-1-\sum_{\text{low-five}}P[i,k]$), dual feasibility becomes
\[
F(k)=8-\frac{5\,g(k)}{4^{n}}\ \ge\ 2^{\,n-\mathrm{rank}\,k}\quad\text{for every relation }k,\qquad F=2^{n}\ \text{at}\ k=(0,1).
\]
Substituting the closed form and setting $y=4^{\lfloor n/2\rfloor-r}$ ($r=\mathrm{rank}\,k$), the slack reduces to the positivity of five scalar quadratics,
\[
5y^2{-}87y{+}546,\quad 5y^2{-}117y{+}636,\quad 5y^2{-}234y{+}2544,\quad 10y^2{-}87y{+}273,\quad 10y^2{-}117y{+}318,
\]
on the powers of four $y\in\{1,4,16,\dots\}$. Two have negative discriminant (positive everywhere); the other three have real roots confined to $(8.6,14.8)$, $(17.2,29.6)$, $(4.3,7.4)$, none of which contains a power of four, so on the admissible $y$ they are bounded below by $44,80,10$ respectively. The objective tends to $\lambda+\mu/s\to 37/4$, i.e.\ $\rho_{BC}(S)/\rho_{\mathrm{avg}}\to 37/8=4.625<5$. Dual feasibility was checked exactly (rational arithmetic, all relations including the membership row) for $n=4,\dots,32$, and slack positivity through the closed form for $n=4,\dots,60$.

\begin{thebibliography}{99}

\bibitem[AAC+22]{AAC+22}
C.~Aguilar~Melchor, N.~Aragon, S.~Bettaieb, L.~Bidoux, O.~Blazy, J.-C.~Deneuville, P.~Gaborit, E.~Persichetti, G.~Z\'emor, \etal\
HQC.
NIST PQC Round 4 Submission, 2022.

\bibitem[ABD+21]{ABD+21}
M.~Albrecht, D.~Bernstein, T.~Chou, C.~Cid, J.~Gilcher, T.~Lange, V.~Maram, I.~von Maurich, R.~Misoczki, R.~Niederhagen, K.~Paterson, E.~Persichetti, C.~Peters, P.~Schwabe, N.~Sendrier, J.~Szefer, C.~Tjhai, M.~Tomlinson, and W.~Wang.
Classic McEliece.
NIST PQC Round 4 Submission, 2022.

\bibitem[ABW10]{ABW10}
B.~Applebaum, B.~Barak, and A.~Wigderson.
Public-key cryptography from different assumptions.
In \emph{STOC}, pages 171--180. ACM, 2010.

\bibitem[Ale11]{Ale11}
M.~Alekhnovich.
More on average case vs approximation complexity.
\emph{Computational Complexity}, 20(4):755--786, 2011.

\bibitem[BFJ+94]{BFJ+94}
A.~Blum, M.~Furst, J.~Jackson, M.~Kearns, Y.~Mansour, and S.~Rudich.
Weakly learning DNF and characterizing statistical query learning using Fourier analysis.
In \emph{STOC}, pages 253--262. ACM, 1994.

\bibitem[BFKL94]{BFKL94}
A.~Blum, M.~Furst, M.~Kearns, and R.~Lipton.
Cryptographic primitives based on hard learning problems.
In \emph{CRYPTO}, pages 278--291. Springer, 1994.

\bibitem[BKW03]{BKW03}
A.~Blum, A.~Kalai, and H.~Wasserman.
Noise-tolerant learning, the parity problem, and the statistical query model.
\emph{J.\ ACM}, 50(4):506--519, 2003.

\bibitem[BPR12]{BPR12}
A.~Banerjee, C.~Peikert, and A.~Rosen.
Pseudorandom functions and lattices.
In \emph{EUROCRYPT}, pages 719--737. Springer, 2012.

\bibitem[BCGI18]{BCGI18}
E.~Boyle, G.~Couteau, N.~Gilboa, and Y.~Ishai.
Compressing vector OLE.
In \emph{ACM CCS}, pages 896--912, 2018.

\bibitem[CRSS98]{CRSS98}
A.~R.~Calderbank, E.~M.~Rains, P.~W.~Shor, and N.~J.~A.~Sloane.
Quantum error correction via codes over GF(4).
\emph{IEEE Trans.\ Inform.\ Theory}, 44(4):1369--1387, 1998.

\bibitem[DKS17]{DKS17}
I.~Diakonikolas, D.~M.~Kane, and A.~Stewart.
Statistical query lower bounds for robust estimation of high-dimensional Gaussians and Gaussian mixtures.
In \emph{FOCS}, pages 73--84. IEEE, 2017.

\bibitem[DGM12]{DGM12}
N.~D\"ottling, J.~M\"uller-Quade, and A.~Nascimento.
IND-CCA secure cryptography based on a variant of the LPN problem.
In \emph{ASIACRYPT}, pages 485--503. Springer, 2012.

\bibitem[EKL22]{EKL2022}
D.~Ellis, G.~Kindler, and N.~Lifshitz.
An analogue of Bonami's lemma for functions on spaces of linear maps, and 2-2 games.
\emph{arXiv:2209.04243}, 2022.

\bibitem[EKM17]{EKM17}
A.~Esser, R.~K\"ubler, and A.~May.
LPN decoded.
In \emph{CRYPTO}, pages 486--514. Springer, 2017.

\bibitem[FGR+17]{FGR+17}
V.~Feldman, E.~Grigorescu, L.~Reyzin, S.~S.~Vempala, and Y.~Xiao.
Statistical algorithms and a lower bound for detecting planted cliques.
\emph{Journal of the ACM}, 64(2):8:1--8:37, 2017.

\bibitem[GJL14]{GJL14}
Q.~Guo, T.~Johansson, and C.~L\"ondahl.
Solving LPN using covering codes.
In \emph{ASIACRYPT}, pages 1--20. Springer, 2014.

\bibitem[GL22]{GL22}
A.~Gollakota and D.~Liang.
On the hardness of PAC-learning stabilizer states with noise.
\emph{Quantum}, 6:640, 2022. arXiv:2102.05174.

\bibitem[Got97]{Got97}
D.~Gottesman.
Stabilizer codes and quantum error correction.
PhD thesis, Caltech, 1997.

\bibitem[Gro96]{Gro96}
L.~K.~Grover.
A fast quantum mechanical algorithm for database search.
In \emph{STOC}, pages 212--219. ACM, 1996.

\bibitem[HKL+12]{HKL+12}
S.~Heyse, E.~Kiltz, V.~Lyubashevsky, C.~Paar, and K.~Pietrzak.
Lapin: An efficient authentication protocol based on ring-LPN.
In \emph{FSE}, pages 346--365. Springer, 2012.

\bibitem[Kea98]{Kea98}
M.~Kearns.
Efficient noise-tolerant learning from statistical queries.
\emph{J.\ ACM}, 45(6):983--1006, 1998.

\bibitem[Kir11]{Kir11}
P.~Kirchner.
Improved generalized birthday attack.
\emph{Cryptology ePrint Archive}, Report 2011/377, 2011.

\bibitem[KLLM24]{KLLM2024}
P.~Keevash, N.~Lifshitz, E.~Long, and D.~Minzer.
Hypercontractivity for global functions and sharp thresholds.
\emph{Journal of the American Mathematical Society}, 37:245--279, 2024. arXiv:1906.05568.

\bibitem[KLP+25]{KLP+25}
A.~Khesin, J.~Lu, A.~Poremba, A.~Ramkumar, and V.~Vaikuntanathan.
Average-Case Complexity of Quantum Stabilizer Decoding.
\emph{arXiv:2509.20697}, 2025.

\bibitem[LF06]{LF06}
E.~Levieil and P.-A.~Fouque.
An improved LPN algorithm.
In \emph{SCN}, pages 348--359. Springer, 2006.

\bibitem[LPQR26]{LPQR26}
J.~Lu, A.~Poremba, Y.~Quek, and A.~Ramkumar.
Post-quantum cryptography from quantum stabilizer decoding.
\emph{arXiv:2603.19110}, 2026.

\bibitem[Lyu05]{Lyu05}
V.~Lyubashevsky.
The parity problem in the presence of noise, decoding random linear codes, and the subset sum problem.
In \emph{APPROX-RANDOM}, pages 378--389. Springer, 2005.

\bibitem[AG04]{AG04}
S.~Aaronson and D.~Gottesman.
Improved simulation of stabilizer circuits.
\emph{Phys.\ Rev.\ A}, 70:052328, 2004. arXiv:quant-ph/0406196.

\bibitem[AHL02]{AHL02}
N.~Alon, S.~Hoory, and N.~Linial.
The Moore bound for irregular graphs.
\emph{Graphs and Combinatorics}, 18(1):53--57, 2002.

\bibitem[BCN89]{BCN89}
A.~E.~Brouwer, A.~M.~Cohen, and A.~Neumaier.
\emph{Distance-Regular Graphs}.
Ergebnisse der Mathematik (3) vol.~18. Springer, 1989.

\bibitem[McE78]{McE78}
R.~J.~McEliece.
A public-key cryptosystem based on algebraic coding theory.
\emph{DSN Progress Report 42--44}, pages 114--116, 1978.

\bibitem[NC00]{NC00}
M.~A.~Nielsen and I.~L.~Chuang.
Quantum Computation and Quantum Information.
Cambridge University Press, 2000.

\bibitem[PQS26]{PQS26}
A.~Poremba, Y.~Quek, and P.~Shor.
The Learning Stabilizers with Noise problem.
In \emph{ITCS}, LIPIcs vol.~362, pages 108:1--108:19, 2026. arXiv:2410.18953.

\bibitem[Pra62]{Pra62}
E.~Prange.
The use of information sets in decoding cyclic codes.
\emph{IRE Trans.\ Inform.\ Theory}, 8(5):5--9, 1962.

\bibitem[Reg05]{Reg05}
O.~Regev.
On lattices, learning with errors, random linear codes, and cryptography.
\emph{J.\ ACM}, 56(6):34, 2009.

\bibitem[Sch18]{Schmidt2018}
K.-U.~Schmidt.
Quadratic and symmetric bilinear forms over finite fields and their association schemes.
\emph{arXiv:1803.04274}, 2018.

\bibitem[Sho94]{Sho94}
P.~W.~Shor.
Algorithms for quantum computation: Discrete logarithms and factoring.
In \emph{FOCS}, pages 124--134. IEEE, 1994.

\bibitem[Tit74]{Tits74}
J.~Tits.
\emph{Buildings of Spherical Type and Finite BN-Pairs}.
Lecture Notes in Mathematics, vol.~386. Springer, 1974.

\end{thebibliography}

\end{document}
